% Options for packages loaded elsewhere
\PassOptionsToPackage{unicode}{hyperref}
\PassOptionsToPackage{hyphens}{url}
\documentclass[
  spanish,
]{article}
\usepackage{xcolor}
\usepackage[margin=2.5cm]{geometry}
\usepackage{amsmath,amssymb}
\setcounter{secnumdepth}{-\maxdimen} % remove section numbering
\usepackage{iftex}
\ifPDFTeX
  \usepackage[T1]{fontenc}
  \usepackage[utf8]{inputenc}
  \usepackage{textcomp} % provide euro and other symbols
\else % if luatex or xetex
  \usepackage{unicode-math} % this also loads fontspec
  \defaultfontfeatures{Scale=MatchLowercase}
  \defaultfontfeatures[\rmfamily]{Ligatures=TeX,Scale=1}
\fi
\usepackage{lmodern}
\ifPDFTeX\else
  % xetex/luatex font selection
\fi
% Use upquote if available, for straight quotes in verbatim environments
\IfFileExists{upquote.sty}{\usepackage{upquote}}{}
\IfFileExists{microtype.sty}{% use microtype if available
  \usepackage[]{microtype}
  \UseMicrotypeSet[protrusion]{basicmath} % disable protrusion for tt fonts
}{}
\makeatletter
\@ifundefined{KOMAClassName}{% if non-KOMA class
  \IfFileExists{parskip.sty}{%
    \usepackage{parskip}
  }{% else
    \setlength{\parindent}{0pt}
    \setlength{\parskip}{6pt plus 2pt minus 1pt}}
}{% if KOMA class
  \KOMAoptions{parskip=half}}
\makeatother
\usepackage{longtable,booktabs,array}
\usepackage{caption}
\captionsetup[table]{skip=6pt}
\newcounter{none} % for unnumbered tables
\usepackage{calc} % for calculating minipage widths
% Correct order of tables after \paragraph or \subparagraph
\usepackage{etoolbox}
\makeatletter
\patchcmd\longtable{\par}{\if@noskipsec\mbox{}\fi\par}{}{}
\makeatother
% Allow footnotes in longtable head/foot
\IfFileExists{footnotehyper.sty}{\usepackage{footnotehyper}}{\usepackage{footnote}}
\makesavenoteenv{longtable}
\ifLuaTeX
\usepackage[bidi=basic,shorthands=off]{babel}
\else
\usepackage[bidi=default,shorthands=off]{babel}
\fi
\ifLuaTeX
  \usepackage{selnolig} % disable illegal ligatures
\fi
\setlength{\emergencystretch}{3em} % prevent overfull lines
\providecommand{\tightlist}{%
  \setlength{\itemsep}{0pt}\setlength{\parskip}{0pt}}
\usepackage{bookmark}
\IfFileExists{xurl.sty}{\usepackage{xurl}}{} % add URL line breaks if available
\urlstyle{same}
% fallback for those not using the hyperref driver hyperxmp:
\makeatletter
\@ifundefined{xmpquote}{\newcommand{\xmpquote}[1]{#1}}{}
\makeatother
\hypersetup{
  pdflang={es},
  hidelinks,
  pdfcreator={LaTeX via pandoc}}

\author{}

% --- Corte de linea en identificadores y rutas largas --------------------
% Los nombres de pruebas, clases y rutas del repositorio se componen en
% \texttt{} y con frecuencia son mas anchos que la linea; sin esto pdfTeX
% no puede partirlos y el texto se sale del margen.
\usepackage[htt]{hyphenat}

% --- Simbolos unicode que pdfTeX clasico no resuelve por defecto ---------
\usepackage{newunicodechar}
\newunicodechar{≥}{\ensuremath{\geq}}
\newunicodechar{≤}{\ensuremath{\leq}}
\newunicodechar{→}{\ensuremath{\rightarrow}}
\newunicodechar{↔}{\ensuremath{\leftrightarrow}}
\newunicodechar{−}{\ensuremath{-}}
\newunicodechar{×}{\ensuremath{\times}}

\date{}

% --- Metadatos institucionales (usados en la portada) ---------------------
\newcommand{\universidad}{Universidad T\'ecnica Estatal de Quevedo}
\newcommand{\facultad}{Facultad de Ciencias de la Computaci\'on}
\newcommand{\carrera}{Carrera de Software}
\newcommand{\asignatura}{Aplicaciones Web}
\newcommand{\docente}{Guerrero Ulloa Gleiston Ciceron}
\newcommand{\periodoacademico}{2025 -- 2026}
\newcommand{\cursoparalelo}{A}
\newcommand{\equipo}{BMT}
\newcommand{\tituloproyecto}{BIOPET --- Sistema web de gesti\'on veterinaria}
\newcommand{\ciudadfecha}{Quevedo, 11 de septiembre de 2026}

\begin{document}

% === PORTADA INSTITUCIONAL ===============================================
\begin{titlepage}
\centering
\vspace*{0.6cm}
{\LARGE\bfseries \universidad \\[0.5cm]}
{\Large \facultad \\[0.2cm]}
{\large \carrera \\[0.2cm]}
{\large Proyecto de Fin de Carrera --- \asignatura \\[0.3cm]}
{\large Periodo acad\'emico: \periodoacademico \quad --- \quad Curso/Paralelo:
\cursoparalelo \\[1.2cm]}
\rule{\linewidth}{0.8pt}\\[0.4cm]
{\Huge\bfseries BIOPET \\[0.35cm]}
{\LARGE Sistema web de gesti\'on veterinaria \\[0.4cm]}
{\Large Especificaci\'on de Requisitos de Software (SRS) --- v1.0.0 \\[0.3cm]}
\rule{\linewidth}{0.8pt}
\vspace{1cm}
{\large \textbf{Docente-director:} \docente \\[0.5cm]}
{\large \textbf{Equipo:} \equipo \\[0.8cm]}
{\large\bfseries Integrantes \\[0.35cm]}
{\small
\renewcommand{\arraystretch}{1.5}
\begin{tabular}{@{}>{\centering\arraybackslash}p{5.2cm}
                >{\centering\arraybackslash}p{4.2cm}
                >{\centering\arraybackslash}p{5.4cm}@{}}
  Fred Adri\'an Beltr\'an Montiel & \texttt{fbeltranm@uteq.edu.ec} &
  ORCID: 0009-0007-8303-2137 \\
  Jaime Josu\'e Mariscal Cabrera & \texttt{jmariscalc@uteq.edu.ec} &
  ORCID: 0009-0007-2206-3941 \\
  Zaida Melissa Taipe Mora & \texttt{ztaipem@uteq.edu.ec} &
  ORCID: 0009-0000-1227-3258 \\
\end{tabular}
}
\vspace{1.2cm}
{\large Repositorio: \url{https://github.com/Grinjoww/ENTREGA-FINAL-BIOPET}}
\vfill
{\large \ciudadfecha}
\end{titlepage}

{
\setcounter{tocdepth}{3}
\tableofcontents
}
\section{Especificación de Requisitos de Software (SRS) ---
BIOPET}\label{especificaciuxf3n-de-requisitos-de-software-srs-biopet}

{\def\LTcaptype{none} % do not increment counter
\begin{longtable}[]{@{}
  >{\raggedright\arraybackslash}p{(\linewidth - 2\tabcolsep) * \real{0.5000}}
  >{\raggedright\arraybackslash}p{(\linewidth - 2\tabcolsep) * \real{0.5000}}@{}}
\toprule\noalign{}
\begin{minipage}[b]{\linewidth}\raggedright
Control documental
\end{minipage} & \begin{minipage}[b]{\linewidth}\raggedright
\end{minipage} \\
\midrule\noalign{}
\endhead
\bottomrule\noalign{}
\endlastfoot
\textbf{Versión documental} & v1.0.0 \\
\textbf{Revisión} & 2026-09-11 \\
\textbf{Fecha de emisión de esta revisión} & 11 de septiembre de 2026 \\
\textbf{Tag histórico de referencia} & \texttt{v1.\allowbreak{}0.\allowbreak{}0}
(\texttt{0d5cd525ce648cca\allowbreak{}7219da204e16fa62\allowbreak{}2e671a87}, 2026-08-18) \\
\textbf{Commit de revisión/corrección actual} &
\texttt{8130ee00b0083808\allowbreak{}fc60567018589e38\allowbreak{}302b375d} \\
\textbf{Repositorio} &
\url{https://github.com/Grinjoww/ENTREGA-FINAL-BIOPET} \\
\end{longtable}
}

\begin{quote}
\textbf{Aviso de correspondencia con el tag (importante).} El tag
\texttt{v1.\allowbreak{}0.\allowbreak{}0} es inmutable y \textbf{no se mueve, no se borra y no se
recrea} en esta revisión. Esta copia del SRS \textbf{no} corresponde
byte a byte al contenido del tag \texttt{v1.\allowbreak{}0.\allowbreak{}0}: el documento fue
corregido después de ese tag (verificable con
\texttt{git\ diff\ v1.\allowbreak{}0.\allowbreak{}0\ -\/-\ docs/\allowbreak{}requisitos/}). Por eso se declaran
por separado la \emph{versión documental} (v1.0.0, que no cambia porque
no se altera el alcance comprometido), el \emph{tag histórico de
referencia} y el \emph{commit de revisión} real sobre el que se generó
esta copia. Cualquier afirmación de evidencia reproducible de este
documento debe reproducirse sobre el commit indicado como ``commit de
revisión'', salvo cuando el propio requisito diga explícitamente que la
evidencia se generó sobre el commit del tag.
\end{quote}

\textbf{Conforme a:} ISO/IEC/IEEE 29148:2018 (estructura de SRS), INCOSE
Guide to Writing Requirements v4 (calidad de requisitos individuales y
de conjunto), criterios INVEST de Cohn (historias de usuario, aplicados
explícitamente en la sección 3.3), plantilla de Cockburn (casos de uso).

\textbf{Estado de aprobación:} documento preparado para revisión y
aprobación del docente-director. Esta copia
(\texttt{SRS-v1.0.0-para-firma.pdf}) se envía específicamente para su
firma electrónica; el estado de firma se actualizará en el documento
histórico (\texttt{SRS-v1.0.0.pdf}), con fecha real, únicamente cuando
exista una firma legítima recibida del docente --- nunca antes, y nunca
simulada o copiada de otro documento.

\textbf{Procedencia de este documento:} este archivo se reconstruye a
partir de dos fuentes: (1) el SRS de la Entrega 1A
(\texttt{PFC\_\allowbreak{}Entrega1\allowbreak{}A\_\allowbreak{}BMT.\allowbreak{}pdf}), que especificaba el backend sobre
ASP.NET Core 8/C\#, y (2) el código realmente implementado en las
Entregas 1B, Tercera Entrega y Unidad IV, sobre \textbf{Java 21 + Spring
Boot 3.2 + Spring Security 6 + Spring Data JPA + PostgreSQL 16 + Redis
7}. Ningún requisito de este documento fue inventado sin sustento: cada
uno proviene del SRS original, del código fuente verificado, o de la
Guía oficial de la Tercera Entrega. El detalle de cada cambio respecto
al documento original está en
\texttt{docs/\allowbreak{}requisitos/\allowbreak{}cambios/\allowbreak{}CAMBIOS-\allowbreak{}SRS.\allowbreak{}md} y, fila por requisito,
en \texttt{docs/\allowbreak{}requisitos/\allowbreak{}CHANGELOG-\allowbreak{}REQ.\allowbreak{}md}.

\begin{center}\rule{0.5\linewidth}{0.5pt}\end{center}

\subsection{1. Introducción}\label{introducciuxf3n}

\subsubsection{1.1. Propósito}\label{propuxf3sito}

Este documento especifica los requisitos funcionales y no funcionales
del sistema BIOPET en su estado \textbf{v1.0.0 (Entrega Final)}. Su
propósito es servir como fuente única de verdad para la matriz de
trazabilidad, las pruebas automatizadas y la evidencia empírica exigidas
por el bloque A.3 de la Guía, incorporando además los módulos de la
Unidad IV (citas, consultas, vacunas, gestión administrativa de usuarios
y consulta externa de especies) que no existían en la revisión v0.9.0-rc
de la Tercera Entrega.

\subsubsection{1.2. Alcance}\label{alcance}

BIOPET es un sistema web para centralizar la información clínica y
administrativa de clínicas veterinarias de pequeña y mediana escala. El
alcance completo del producto (definido en la Entrega 1A) contempla
gestión de dueños y mascotas, historial clínico, citas, telemetría IoT,
recomendación clínica asistida, facturación digital y reportes.

\textbf{Alcance implementado y verificado en v1.0.0:} autenticación
completa (registro, login, refresh, logout con revocación), control de
acceso por rol (RBAC), CRUD completo de la entidad Mascota con
verificación de propiedad, resumen agregado de mascotas por especie (vía
procedimiento almacenado), gestión administrativa de usuarios, gestión
de vacunas, gestión de citas a nivel de API y consulta externa de
información de especies con caché.

\textbf{Ampliado en v1.0.0 (Unidad IV):} gestión administrativa de
usuarios (REQ-F-023), CRUD de vacunas (REQ-F-024) y consulta externa de
información de especies con caché (REQ-F-025), con interfaz de usuario
real para vacunas
(\texttt{frontend/\allowbreak{}src/\allowbreak{}app/\allowbreak{}features/\allowbreak{}vacunas.\allowbreak{}component.\allowbreak{}ts}) y solo a nivel
de API para usuarios y especies externas (sin pantalla propia todavía).
Adicionalmente, el backend incorporó CRUD completo de citas
(\texttt{Cita\allowbreak{}Controller}) y de consultas médicas
(\texttt{Consulta\allowbreak{}Controller}), que cubren respectivamente la obligación
A de REQ-F-015 y la de REQ-F-013 (detalle de alcance real, sin
sobre-declarar, en la sección 3.1).

\textbf{Alcance pendiente, heredado de la Entrega 1A:} vista consolidada
de historial clínico cronológico (obligación B de REQ-F-013) y
calendario interactivo de citas (obligación B de REQ-F-015) --- ambos
con trabajo parcial de backend desde v1.0.0 ---, más prescripción de
medicamentos, telemetría IoT, recomendaciones asistidas, facturación
digital, reportes exportables, notificaciones por correo y recuperación
de contraseña. El modelo de datos conceptual de estos módulos ya existe
(ver sección 5 y el DER de la Entrega 1A). Se documentan como requisitos
con estado \texttt{implementado} o \texttt{pendiente} (según la
taxonomía de la sección 1.3.1) para que la matriz de trazabilidad
(bloque A.3.3 de la Guía) los declare correctamente, en vez de
omitirlos.

\subsubsection{1.3. Definiciones, acrónimos y
abreviaturas}\label{definiciones-acruxf3nimos-y-abreviaturas}

{\def\LTcaptype{none} % do not increment counter
\begin{longtable}[]{@{}
  >{\raggedright\arraybackslash}p{(\linewidth - 2\tabcolsep) * \real{0.5000}}
  >{\raggedright\arraybackslash}p{(\linewidth - 2\tabcolsep) * \real{0.5000}}@{}}
\toprule\noalign{}
\begin{minipage}[b]{\linewidth}\raggedright
Término
\end{minipage} & \begin{minipage}[b]{\linewidth}\raggedright
Definición
\end{minipage} \\
\midrule\noalign{}
\endhead
\bottomrule\noalign{}
\endlastfoot
RBAC & Role-Based Access Control: control de acceso basado en el rol del
usuario autenticado. \\
JWT & JSON Web Token (RFC 7519): token firmado que representa
afirmaciones sobre un sujeto autenticado. \\
ORM & Object-Relational Mapping: mapeo objeto-relacional (Spring Data
JPA / Hibernate en este proyecto). \\
SP & Stored Procedure / procedimiento almacenado en el motor de base de
datos. \\
TTL & Time To Live: tiempo de vida de una entrada de caché. \\
MoSCoW & Técnica de priorización: Must, Should, Could, Won't. \\
deberá & Verbo modal normativo (equivalente al \emph{shall} de
ISO/IEC/IEEE 29148) que indica obligación contractual del requisito.
Todos los enunciados de la sección 3 lo usan. \\
Aislamiento de datos & Regla por la cual un usuario con rol
\texttt{ROLE\_\allowbreak{}DUENO} solo puede acceder a los recursos asociados a su
propia cuenta (ver REQ-NF-015). \\
\end{longtable}
}

\paragraph{1.3.1. Taxonomía de estados de requisito (valores
permitidos)}\label{taxonomuxeda-de-estados-de-requisito-valores-permitidos}

Todo requisito de este documento declara su campo \textbf{Estado} con
\textbf{uno y solo uno} de los tres valores siguientes. No existe ningún
cuarto valor: los matices (cobertura parcial de pruebas, defectos
conocidos, evidencia faltante) se registran en el campo
\textbf{Observaciones} del propio requisito, no inventando un estado
nuevo.

{\def\LTcaptype{none} % do not increment counter
\begin{longtable}[]{@{}
  >{\raggedright\arraybackslash}p{(\linewidth - 2\tabcolsep) * \real{0.5000}}
  >{\raggedright\arraybackslash}p{(\linewidth - 2\tabcolsep) * \real{0.5000}}@{}}
\toprule\noalign{}
\begin{minipage}[b]{\linewidth}\raggedright
Estado
\end{minipage} & \begin{minipage}[b]{\linewidth}\raggedright
Significado exacto
\end{minipage} \\
\midrule\noalign{}
\endhead
\bottomrule\noalign{}
\endlastfoot
\texttt{verificado} & \textbf{El requisito se da por cumplido.} La
conducta exigida está realizada en el repositorio y \textbf{todos} sus
criterios de aceptación están respaldados por un artefacto reproducible
y nombrado: prueba automatizada, captura HTTP o de log, corrida de
medición archivada, o inspección de un fragmento de código/configuración
citado explícitamente (método \emph{Inspection} de ISO/IEC/IEEE
29148). \\
\texttt{implementado} & \textbf{El requisito NO se da por cumplido.}
Existe código real en el repositorio que lo realiza, total o
parcialmente, pero al menos uno de sus criterios de aceptación
\textbf{falla} por un defecto conocido, o bien \textbf{carece de
artefacto} que lo respalde. La causa exacta se declara siempre en
Observaciones. \\
\texttt{pendiente} & \textbf{El requisito NO se da por cumplido.} No
existe en el repositorio implementación específica de lo que el
requisito exige. Se conserva especificado (con criterios de aceptación y
método de verificación previstos) para que la matriz de trazabilidad lo
declare sin ocultarlo. \\
\end{longtable}
}

\begin{quote}
\textbf{Lectura obligada de \texttt{implementado}.}
\texttt{implementado} \textbf{nunca} significa ``entregado'' ni
``cumplido parcialmente de forma aceptable'': significa que el
requisito, tomado como un todo, \textbf{no está satisfecho hoy}. Se
distingue de \texttt{pendiente} únicamente en que existe código real que
lo aborda, dato que se conserva para no ocultar el trabajo hecho ni
exagerar el que falta. Un requisito \texttt{implementado} cuenta como
\textbf{no cumplido} en todos los recuentos de este documento (secciones
3.4.1 y 3.4.2).
\end{quote}

\paragraph{1.3.2. Plantilla única de
requisito}\label{plantilla-uxfanica-de-requisito}

Todos los requisitos de las secciones 3.1 y 3.2 usan exactamente esta
plantilla, sin excepción:

\begin{itemize}
\tightlist
\item
  \textbf{Tipo} --- Funcional, o la categoría no funcional (Rendimiento,
  Seguridad, Usabilidad, Mantenibilidad, Compatibilidad, Disponibilidad,
  Operación).
\item
  \textbf{Prioridad} --- MoSCoW: Must, Should o Could.
\item
  \textbf{Enunciado} --- patrón
  \texttt{{[}condición{]}\ {[}\allowbreak{}sujeto{]}\ deberá\ {[}\allowbreak{}acción{]}\ {[}obje\allowbreak{}to{]}\ {[}restricc\allowbreak{}ión{]}}
  de ISO/IEC/IEEE 29148.
\item
  \textbf{Rationale} --- origen y justificación del requisito.
\item
  \textbf{Criterios de aceptación} --- condiciones observables, en forma
  \emph{Dado / cuando / entonces}, que permiten responder CUMPLE / NO
  CUMPLE sin juicio subjetivo.
\item
  \textbf{Método de verificación} --- cómo se comprueban esos criterios
  (prueba automatizada, inspección, demostración, análisis o medición),
  con el artefacto concreto cuando existe.
\item
  \textbf{Trazabilidad} --- rutas reales del repositorio: historia de
  usuario, caso de uso, clases, endpoints, pruebas, scripts,
  configuración, decisiones de arquitectura o evidencias, según
  corresponda al tipo de requisito.
\item
  \textbf{Estado} --- uno de los tres valores de la sección 1.3.1.
\item
  \textbf{Observaciones} --- solo cuando hay un matiz, defecto o
  limitación que declarar.
\end{itemize}

\textbf{Nota sobre requisitos con más de una obligación.} Tres
requisitos heredados agrupan dos obligaciones distintas con distinto
grado de avance: \texttt{REQ-\allowbreak{}F-\allowbreak{}013} (registrar consultas médicas /
consultar el historial consolidado), \texttt{REQ-\allowbreak{}F-\allowbreak{}015} (gestionar citas
por API / presentarlas en un calendario interactivo) y
\texttt{REQ-\allowbreak{}NF-\allowbreak{}012} (TTL de caché configurable / medir la tasa de
aciertos). En lugar de crear identificadores nuevos, \textbf{cada uno
conserva su identificador único} y separa sus obligaciones dentro del
propio bloque, bajo los epígrafes \textbf{Obligación A} y
\textbf{Obligación B}, con criterios de aceptación numerados por
obligación (\texttt{A1}, \texttt{A2}, \ldots, \texttt{B1}, \texttt{B2},
\ldots). El campo Estado es \textbf{uno solo} y describe el requisito
completo: basta con que una de las dos obligaciones no esté cumplida
para que el requisito no pueda declararse \texttt{verificado}.

Se descartó deliberadamente la alternativa de crear subrequisitos con
sufijo de letra por dos motivos: rompería la correspondencia exacta de
identificadores entre este SRS y \texttt{docs/\allowbreak{}trazabilidad/\allowbreak{}matriz.\allowbreak{}csv}
---que el guion de CI \texttt{scripts/\allowbreak{}validate-\allowbreak{}traceability.\allowbreak{}sh}
comprueba en ambos sentidos--- y fragmentaría la trazabilidad histórica
hacia las historias de usuario, los casos de uso y los identificadores
\texttt{RF-\allowbreak{}NN} de la Entrega 1A, que están fijados sobre el
identificador base. \textbf{Ningún requisito se renumera y no existe
ningún identificador con sufijo de letra en este documento.}

\subsubsection{1.4. Referencias}\label{referencias}

\begin{itemize}
\tightlist
\item
  ISO/IEC/IEEE 29148:2018 --- Requirements Engineering.
\item
  INCOSE Guide to Writing Requirements v4.
\item
  RFC 7519 (JWT), RFC 7807 (ProblemDetails).
\item
  Guía de la Tercera Entrega --- PFC Aplicaciones Web 2026-2027, UTEQ.
\item
  \texttt{PFC\_\allowbreak{}Entrega1\allowbreak{}A\_\allowbreak{}BMT.\allowbreak{}pdf} --- SRS y diseño original (ASP.NET
  Core).
\item
  \texttt{docs/\allowbreak{}requisitos/\allowbreak{}cambios/\allowbreak{}CAMBIOS-\allowbreak{}SRS.\allowbreak{}md} --- registro de
  cambios respecto al documento original.
\item
  \texttt{docs/\allowbreak{}requisitos/\allowbreak{}CHANGELOG-\allowbreak{}REQ.\allowbreak{}md} --- registro formal de
  cambios fila por requisito.
\item
  \texttt{docs/\allowbreak{}trazabilidad/\allowbreak{}matriz.\allowbreak{}csv} --- matriz de trazabilidad
  (bloque A.3.3).
\item
  \texttt{docs/\allowbreak{}basedatos/\allowbreak{}CATALOGO-\allowbreak{}SP.\allowbreak{}md} --- catálogo de procedimientos
  almacenados.
\item
  \texttt{docs/\allowbreak{}adr/} --- decisiones de arquitectura (ADR-002 a ADR-007).
\end{itemize}

\subsubsection{1.5. Resumen del documento}\label{resumen-del-documento}

La sección 2 describe el producto de forma global (perspectiva,
funciones, usuarios, restricciones, interfaces externas, matriz de
permisos y estados del dominio). La sección 3 contiene el detalle de
cada requisito funcional y no funcional con la plantilla única de la
sección 1.3.2, e incluye en 3.4 la tabla de control de completitud de
los campos obligatorios. La sección 4 resume la trazabilidad end-to-end.
Las secciones 5 y 6 remiten a los artefactos de modelo de datos e
interfaz que ya existen en el repositorio. La sección 7 documenta, sin
inventar contenido, lo que aún falta por completar.

\subsubsection{1.6. Control documental e historial de
revisiones}\label{control-documental-e-historial-de-revisiones}

Este documento tiene \textbf{un solo índice}: la tabla de contenidos
generada automáticamente al inicio del PDF. La lista manual de secciones
que duplicaba ese índice en revisiones anteriores fue eliminada en la
revisión 2026-09-11.

Cada fila del historial está respaldada por un commit real del
repositorio (\texttt{git\ log\ -\/-\ docs/\allowbreak{}requisitos/\allowbreak{}SRS.\allowbreak{}md}). No se
declara ninguna revisión anterior que no pueda sustentarse con historial
Git.

{\def\LTcaptype{none} % do not increment counter
\begin{longtable}[]{@{}
  >{\raggedright\arraybackslash}p{(\linewidth - 6\tabcolsep) * \real{0.2500}}
  >{\raggedright\arraybackslash}p{(\linewidth - 6\tabcolsep) * \real{0.2500}}
  >{\raggedright\arraybackslash}p{(\linewidth - 6\tabcolsep) * \real{0.2500}}
  >{\raggedright\arraybackslash}p{(\linewidth - 6\tabcolsep) * \real{0.2500}}@{}}
\toprule\noalign{}
\begin{minipage}[b]{\linewidth}\raggedright
Revisión
\end{minipage} & \begin{minipage}[b]{\linewidth}\raggedright
Fecha
\end{minipage} & \begin{minipage}[b]{\linewidth}\raggedright
Commit
\end{minipage} & \begin{minipage}[b]{\linewidth}\raggedright
Resumen del cambio
\end{minipage} \\
\midrule\noalign{}
\endhead
\bottomrule\noalign{}
\endlastfoot
v0.7.0 & 2026-06-14 & --- & Primera consolidación de requisitos tras la
migración de pila (registro histórico en \texttt{CHANGELOG-\allowbreak{}REQ.\allowbreak{}md}). \\
v0.9.0-rc & 2026-07-30 & \texttt{2115a35} & Consolidación de historias
de usuario y casos de uso en el SRS; creación de
\texttt{CHANGELOG-\allowbreak{}REQ.\allowbreak{}md}. \\
v0.9.0-rc & 2026-07-31 & \texttt{73c41b4} & Cierre de observaciones
oficiales de las Entregas 1A y 1B (OBS-02 a OBS-05). \\
v0.9.0-rc & 2026-08-10 & \texttt{9804ed3} & Reconciliación de SRS y
matriz de trazabilidad con la Unidad IV (corrección del CI de
trazabilidad). \\
v0.9.0-rc & 2026-08-17 & \texttt{5b735e6} & Enlace de evidencia para
varios REQ-NF pendientes. \\
v0.9.0-rc & 2026-08-18 & \texttt{91b72b5} & Cierre de REQ-NF-005 con la
re-corrida Lighthouse (móvil y escritorio) y cierre de PRISMA en la
sección 7. \\
v1.0.0 & 2026-09-03 & \texttt{4f53376} & Paso a v1.0.0: sección 3.3
(INVEST sobre las 24 historias), alineación de REQ-F-013 / REQ-F-015 /
REQ-F-025 con el código de Unidad IV. \\
v1.0.0 & 2026-09-04 & \texttt{a9f0562} & Cierre de trazabilidad y
evidencia de requisitos. \\
v1.0.0 & 2026-09-07 & \texttt{70c3596}, \texttt{e9b1dc5} &
Sincronización SRS / matriz / código antes de la firma; portada
institucional y URL del repositorio en la copia para firma. \\
v1.0.0 & 2026-09-11 & \texttt{8130ee0} (base) & \textbf{Esta revisión.}
Corrección integral tras la retroalimentación del docente-director:
plantilla única con criterios de aceptación verificables en todos los
requisitos, bloque individual para REQ-F-013 a REQ-F-020, taxonomía de
tres estados, descomposición de requisitos compuestos, interfaces
externas, matriz de permisos, estados del dominio, corrección de la
sección 2.6 y separación explícita entre versión documental, tag
histórico y commit de revisión. \\
\end{longtable}
}

\begin{center}\rule{0.5\linewidth}{0.5pt}\end{center}

\subsection{2. Descripción global}\label{descripciuxf3n-global}

\subsubsection{2.1. Perspectiva del
producto}\label{perspectiva-del-producto}

BIOPET es un sistema nuevo, no una extensión de un producto previo
existente en la clínica. Reemplaza procesos manuales (hojas de cálculo,
mensajería) por una plataforma web centralizada, según el problema
documentado en la Entrega 1A (entrevistas a tres veterinarios y dos
auxiliares, encuestas a quince dueños de mascotas).

\subsubsection{2.2. Cambio de plataforma tecnológica respecto a la
Entrega
1A}\label{cambio-de-plataforma-tecnoluxf3gica-respecto-a-la-entrega-1a}

El SRS original (Entrega 1A) especificaba ASP.NET Core 8/C\# como
backend (ADR-001 original) y Bootstrap 5 + HTML/CSS/JS como frontend. El
equipo migró la implementación real a \textbf{Java 21 + Spring Boot 3.2}
en el backend y \textbf{Angular 17+} en el frontend, decisión reflejada
en \texttt{ADR-\allowbreak{}002-\allowbreak{}pila-\allowbreak{}tecnologica.\allowbreak{}md} del repositorio actual. Los
requisitos funcionales de negocio (qué hace el sistema) no cambiaron;
los requisitos no funcionales técnicos (cómo se implementa: framework,
ORM, mecanismo JWT) se actualizaron para reflejar la pila real, evitando
que el SRS describa un sistema que no es el que existe en el
repositorio.

\subsubsection{2.3. Funciones del producto
(resumen)}\label{funciones-del-producto-resumen}

\begin{itemize}
\tightlist
\item
  Gestión de usuarios y autenticación (roles: \texttt{ROLE\_\allowbreak{}ADMIN},
  \texttt{ROLE\_\allowbreak{}VETERINARIO}, \texttt{ROLE\_\allowbreak{}AUXILIAR},
  \texttt{ROLE\_\allowbreak{}DUENO}), incluida la administración de cuentas por un
  \texttt{ROLE\_\allowbreak{}ADMIN} (Unidad IV).
\item
  Gestión de mascotas (CRUD + resumen agregado por especie).
\item
  Gestión de vacunas (CRUD, con interfaz de usuario) y de citas y
  consultas médicas (CRUD a nivel de API; sin calendario interactivo ni
  vista de historial clínico consolidado todavía --- ver 3.1).
\item
  Consulta de información externa de especies (taxonomía, hábitat,
  dieta), con caché Redis.
\item
  \emph{(Pendiente)} Vista consolidada de historial clínico cronológico,
  prescripción de medicamentos, telemetría IoT, recomendaciones
  asistidas, facturación, reportes exportables, notificaciones por
  correo y recuperación de contraseña.
\end{itemize}

\subsubsection{2.4. Características de los
usuarios}\label{caracteruxedsticas-de-los-usuarios}

{\def\LTcaptype{none} % do not increment counter
\begin{longtable}[]{@{}
  >{\raggedright\arraybackslash}p{(\linewidth - 4\tabcolsep) * \real{0.3333}}
  >{\raggedright\arraybackslash}p{(\linewidth - 4\tabcolsep) * \real{0.3333}}
  >{\raggedright\arraybackslash}p{(\linewidth - 4\tabcolsep) * \real{0.3333}}@{}}
\toprule\noalign{}
\begin{minipage}[b]{\linewidth}\raggedright
Rol
\end{minipage} & \begin{minipage}[b]{\linewidth}\raggedright
Descripción
\end{minipage} & \begin{minipage}[b]{\linewidth}\raggedright
Nivel técnico esperado
\end{minipage} \\
\midrule\noalign{}
\endhead
\bottomrule\noalign{}
\endlastfoot
\texttt{ROLE\_\allowbreak{}ADMIN} & Supervisa roles, operatividad general del sistema
y tiene visibilidad global de datos. Único rol con acceso al CRUD
administrativo de usuarios. & Medio-alto (personal administrativo de la
clínica). \\
\texttt{ROLE\_\allowbreak{}VETERINARIO} & Registra consultas médicas y vacunas;
gestiona mascotas; solo puede modificar las citas en las que él es el
veterinario asignado. & Medio (personal clínico). \\
\texttt{ROLE\_\allowbreak{}AUXILIAR} & Apoya operaciones administrativas: agenda
citas, gestiona mascotas, vacunas y consultas. & Medio. \\
\texttt{ROLE\_\allowbreak{}DUENO} & Consulta únicamente la información asociada a sus
propias mascotas. No crea, actualiza ni elimina recursos clínicos. &
Básico (público general, sin capacitación previa). \\
\end{longtable}
}

\subsubsection{2.5. Restricciones}\label{restricciones}

\begin{itemize}
\tightlist
\item
  El proyecto debe completarse dentro de un semestre académico (PPA
  2026-2027), lo que limita el alcance implementado al núcleo Must Have.
\item
  La base de datos debe ser PostgreSQL 16 (decisión ya tomada, ver
  \texttt{ADR-\allowbreak{}004-\allowbreak{}postgresql.\allowbreak{}md}).
\item
  La autenticación debe usar JWT en cookies
  \texttt{Http\allowbreak{}Only\ +\ Secure\ +\ Same\allowbreak{}Site=Strict} (bloque A.1 de la
  Guía), no \texttt{local\allowbreak{}Storage} ni \texttt{session\allowbreak{}Storage}.
\item
  Toda operación de base de datos que no sea CRUD elemental debe
  implementarse como función/procedimiento almacenado (bloque A.2 de la
  Guía), no como JPQL/HQL con joins o agregaciones.
\item
  \textbf{Restricción operacional (compromiso académico, no requisito
  verificable).} El equipo se compromete a mantener el sistema
  desplegable y operativo durante las semanas de evaluación establecidas
  por la asignatura. Esta condición se registra aquí como restricción
  operacional y \textbf{no} como requisito con estado de cumplimiento,
  porque el proyecto no dispone de monitoreo continuo que permita medir
  disponibilidad histórica. La parte que sí es técnicamente verificable
  (existencia y respuesta del endpoint de verificación de estado) se
  especifica en REQ-NF-007.
\end{itemize}

\subsubsection{2.6. Supuestos y
dependencias}\label{supuestos-y-dependencias}

\begin{itemize}
\tightlist
\item
  Se asume disponibilidad de un motor PostgreSQL 16 y Redis 7 accesibles
  desde el backend (vía Docker Compose en desarrollo; servicios
  gestionados en el despliegue de Render, ver
  \texttt{docs/\allowbreak{}despliegue/\allowbreak{}DEPLOYMENT.\allowbreak{}md}).
\item
  Se asume que el operador del despliegue ejecuta el procedimiento de
  respaldo descrito en \texttt{docs/\allowbreak{}despliegue/\allowbreak{}BACKUP.\allowbreak{}md} (ver
  REQ-NF-017).
\end{itemize}

\textbf{Servicios externos realmente integrados en v1.0.0 (corrección de
esta revisión).} Una versión anterior de este documento afirmaba que
``ninguno está implementado todavía'', lo cual contradecía a REQ-F-025.
La situación real, verificada contra el código, es la siguiente:

{\def\LTcaptype{none} % do not increment counter
\begin{longtable}[]{@{}
  >{\raggedright\arraybackslash}p{(\linewidth - 6\tabcolsep) * \real{0.2500}}
  >{\raggedright\arraybackslash}p{(\linewidth - 6\tabcolsep) * \real{0.2500}}
  >{\raggedright\arraybackslash}p{(\linewidth - 6\tabcolsep) * \real{0.2500}}
  >{\raggedright\arraybackslash}p{(\linewidth - 6\tabcolsep) * \real{0.2500}}@{}}
\toprule\noalign{}
\begin{minipage}[b]{\linewidth}\raggedright
Servicio externo
\end{minipage} & \begin{minipage}[b]{\linewidth}\raggedright
Estado real
\end{minipage} & \begin{minipage}[b]{\linewidth}\raggedright
Dónde está en el código
\end{minipage} & \begin{minipage}[b]{\linewidth}\raggedright
Requisito
\end{minipage} \\
\midrule\noalign{}
\endhead
\bottomrule\noalign{}
\endlastfoot
\textbf{API Ninjas --- Animals API}
(\texttt{https:/\allowbreak{}/\allowbreak{}api.\allowbreak{}api-\allowbreak{}ninjas.\allowbreak{}com/\allowbreak{}v1/\allowbreak{}animals}) & \textbf{Integrado y
en uso.} Consulta de taxonomía, hábitat y dieta por nombre de especie,
con caché \emph{cache-aside} en Redis. Autenticación por cabecera
\texttt{X-\allowbreak{}Api-\allowbreak{}Key} (\texttt{APP\_\allowbreak{}EXTERNAL\_\allowbreak{}API\_\allowbreak{}KE\allowbreak{}Y}). &
\texttt{Backend/\allowbreak{}src/\allowbreak{}main/\allowbreak{}java/\allowbreak{}com/\allowbreak{}biopet/\allowbreak{}integration/\allowbreak{}External\allowbreak{}Api\allowbreak{}Client.\allowbreak{}java},
\texttt{External\allowbreak{}Api\allowbreak{}Service.\allowbreak{}java}, \texttt{Rest\allowbreak{}Template\allowbreak{}Config.\allowbreak{}java} &
REQ-F-025 \\
Servicio de correo electrónico & \textbf{No integrado.} Sin cliente, sin
configuración y sin proveedor elegido. & --- & REQ-F-022
(\texttt{pendiente}) \\
Servicio de IA para recomendaciones clínicas & \textbf{No integrado.}
Sin cliente y sin decisión de arquitectura formal sobre el proveedor. &
--- & REQ-F-017 (\texttt{pendiente}) \\
Plataforma de telemetría IoT & \textbf{No integrado.} Sin API de
ingesta. & --- & REQ-F-016 (\texttt{pendiente}) \\
\end{longtable}
}

Las tres integraciones no implementadas se mantienen documentadas como
requisitos con estado \texttt{pendiente}; \textbf{no} se declaran como
implementadas ni se les asigna configuración, endpoint o evidencia que
no exista.

\subsubsection{2.7. Interfaces externas}\label{interfaces-externas}

Esta sección describe las interfaces del sistema con su entorno,
conforme a ISO/IEC/IEEE 29148:2018. Solo se describen interfaces que
existen realmente en el repositorio.

\paragraph{2.7.1. Interfaces de usuario}\label{interfaces-de-usuario}

La interfaz es una aplicación web de página única (Angular 17+) servida
como estático. Las pantallas realmente implementadas están en
\texttt{frontend/\allowbreak{}src/\allowbreak{}app/\allowbreak{}features/}: \texttt{login.\allowbreak{}component.\allowbreak{}ts},
\texttt{mascotas.\allowbreak{}component.\allowbreak{}ts} y \texttt{vacunas.\allowbreak{}component.\allowbreak{}ts}. El
detalle está en la sección 6.

\paragraph{2.7.2. Interfaces de software}\label{interfaces-de-software}

{\def\LTcaptype{none} % do not increment counter
\begin{longtable}[]{@{}
  >{\raggedright\arraybackslash}p{(\linewidth - 6\tabcolsep) * \real{0.2500}}
  >{\raggedright\arraybackslash}p{(\linewidth - 6\tabcolsep) * \real{0.2500}}
  >{\raggedright\arraybackslash}p{(\linewidth - 6\tabcolsep) * \real{0.2500}}
  >{\raggedright\arraybackslash}p{(\linewidth - 6\tabcolsep) * \real{0.2500}}@{}}
\toprule\noalign{}
\begin{minipage}[b]{\linewidth}\raggedright
Interfaz
\end{minipage} & \begin{minipage}[b]{\linewidth}\raggedright
Protocolo / mecanismo
\end{minipage} & \begin{minipage}[b]{\linewidth}\raggedright
Dirección
\end{minipage} & \begin{minipage}[b]{\linewidth}\raggedright
Artefacto en el repositorio
\end{minipage} \\
\midrule\noalign{}
\endhead
\bottomrule\noalign{}
\endlastfoot
API REST de BIOPET & HTTP/JSON sobre TLS; errores en formato
\texttt{Problem\allowbreak{}Detail} (RFC 7807); autenticación JWT en cookies
\texttt{Http\allowbreak{}Only\ +\ Secure\ +\ Same\allowbreak{}Site=Strict} o cabecera
\texttt{Authorization:\ Be\allowbreak{}arer} & Entrante (frontend y clientes) &
\texttt{com.\allowbreak{}biopet.\allowbreak{}controller.*},
\texttt{com.\allowbreak{}biopet.\allowbreak{}exception.\allowbreak{}Problem\allowbreak{}Detail\allowbreak{}Factory} \\
Contrato OpenAPI & OpenAPI 3 (springdoc), expuesto en
\texttt{/\allowbreak{}api/\allowbreak{}openapi} y navegable en \texttt{/\allowbreak{}api/\allowbreak{}docs} & Saliente
(documentación) & \texttt{com.\allowbreak{}biopet.\allowbreak{}config.\allowbreak{}Open\allowbreak{}Api\allowbreak{}Config},
\texttt{application.\allowbreak{}yml} (\texttt{springdoc.*}) \\
PostgreSQL 16 & JDBC (\texttt{org.\allowbreak{}postgresql.\allowbreak{}Driver}); esquema
versionado con Flyway; rutinas almacenadas invocadas con
\texttt{@Procedure} & Saliente & \texttt{application.\allowbreak{}yml}
(\texttt{spring.\allowbreak{}datasource.*}, \texttt{spring.\allowbreak{}flyway.*}),
\texttt{Backend/\allowbreak{}src/\allowbreak{}main/\allowbreak{}resources/\allowbreak{}db/\allowbreak{}migration/}, \texttt{db/\allowbreak{}procs/} \\
Redis 7 & Protocolo RESP vía Spring Data Redis
(\texttt{String\allowbreak{}Redis\allowbreak{}Template}) y caché de Spring
(\texttt{spring.\allowbreak{}cache.\allowbreak{}type:\ redis}) & Saliente &
\texttt{application.\allowbreak{}yml} (\texttt{spring.\allowbreak{}data.\allowbreak{}redis.*},
\texttt{spring.\allowbreak{}cache.\allowbreak{}redis.*}), \texttt{Token\allowbreak{}Blacklist\allowbreak{}Service},
\texttt{External\allowbreak{}Api\allowbreak{}Service} \\
API Ninjas --- Animals API & HTTPS/JSON con \texttt{Rest\allowbreak{}Template}; clave
en cabecera \texttt{X-\allowbreak{}Api-\allowbreak{}Key} & Saliente & \texttt{External\allowbreak{}Api\allowbreak{}Client},
\texttt{Rest\allowbreak{}Template\allowbreak{}Config} \\
Endpoints de operación & HTTP/JSON; Spring Boot Actuator, exposición
limitada a \texttt{health}, \texttt{info} y \texttt{metrics} & Entrante
(operador / plataforma) & \texttt{application.\allowbreak{}yml}
(\texttt{management.\allowbreak{}endpoints.\allowbreak{}web.\allowbreak{}exposure.\allowbreak{}include}) \\
\end{longtable}
}

\paragraph{2.7.3. Interfaces de hardware}\label{interfaces-de-hardware}

BIOPET no interactúa directamente con hardware. La API de ingesta de
telemetría de dispositivos IoT prevista en la Entrega 1A corresponde a
REQ-F-016 y está en estado \texttt{pendiente}: no existe endpoint,
cliente ni protocolo definido en el repositorio.

\paragraph{2.7.4. Interfaces de
comunicación}\label{interfaces-de-comunicaciuxf3n}

{\def\LTcaptype{none} % do not increment counter
\begin{longtable}[]{@{}
  >{\raggedright\arraybackslash}p{(\linewidth - 4\tabcolsep) * \real{0.3333}}
  >{\raggedright\arraybackslash}p{(\linewidth - 4\tabcolsep) * \real{0.3333}}
  >{\raggedright\arraybackslash}p{(\linewidth - 4\tabcolsep) * \real{0.3333}}@{}}
\toprule\noalign{}
\begin{minipage}[b]{\linewidth}\raggedright
Aspecto
\end{minipage} & \begin{minipage}[b]{\linewidth}\raggedright
Valor real configurado
\end{minipage} & \begin{minipage}[b]{\linewidth}\raggedright
Artefacto
\end{minipage} \\
\midrule\noalign{}
\endhead
\bottomrule\noalign{}
\endlastfoot
Transporte & HTTPS/TLS 1.3 (conector dual HTTP/HTTPS en desarrollo) &
\texttt{Tomcat\allowbreak{}Dual\allowbreak{}Connector\allowbreak{}Config}, \texttt{application-\allowbreak{}tls.\allowbreak{}yml},
\texttt{docker-\allowbreak{}compose.\allowbreak{}tls.\allowbreak{}yml} \\
CORS & Orígenes permitidos por variable \texttt{CORS\_\allowbreak{}ALLOWED\_\allowbreak{}ORIGI\allowbreak{}NS}
(por defecto \texttt{http:/\allowbreak{}/\allowbreak{}localhost:4200}) &
\texttt{Security\allowbreak{}Config.\allowbreak{}cors\allowbreak{}Configuration\allowbreak{}Source()} \\
Cabeceras de seguridad & HSTS (\texttt{max-\allowbreak{}age=31536000},
\texttt{include\allowbreak{}Sub\allowbreak{}Domains}, \texttt{preload}), CSP
\texttt{default-\allowbreak{}src\ \textquotesingle{}self\textquotesingle{};\allowbreak{}\ frame-\allowbreak{}ancestors\ \textquotesingle{}none\textquotesingle{};\allowbreak{}\ object-\allowbreak{}src\ \textquotesingle{}none\textquotesingle{}},
\texttt{X-\allowbreak{}Content-\allowbreak{}Type-\allowbreak{}Options}, \texttt{X-\allowbreak{}Frame-\allowbreak{}Options:\ DENY},
\texttt{Referrer-\allowbreak{}Policy:\ no-\allowbreak{}referrer} & \texttt{Security\allowbreak{}Config}, prueba
\texttt{Security\allowbreak{}Headers\allowbreak{}Test} \\
Sesión & Sin estado de sesión en servidor
(\texttt{Session\allowbreak{}Creation\allowbreak{}Policy.\allowbreak{}STATELESS}) & \texttt{Security\allowbreak{}Config} \\
\end{longtable}
}

\subsubsection{2.8. Matriz de permisos por rol y
operación}\label{matriz-de-permisos-por-rol-y-operaciuxf3n}

Matriz construida \textbf{exclusivamente} a partir del código de este
repositorio: las anotaciones \texttt{@Pre\allowbreak{}Authorize} de
\texttt{com.\allowbreak{}biopet.\allowbreak{}controller.*} (control por rol) y las reglas de datos
implementadas en \texttt{com.\allowbreak{}biopet.\allowbreak{}service.*} (control por propiedad
del recurso). No se toma ninguna regla de otro proyecto.

Leyenda de las columnas de rol:

\begin{itemize}
\tightlist
\item
  \textbf{Sí} --- el rol puede ejecutar la operación sobre cualquier
  registro.
\item
  \textbf{Solo propias} --- el rol puede ejecutar la operación, pero el
  resultado debe restringirse a los recursos asociados a su propia
  cuenta.
\item
  \textbf{No} --- la solicitud se rechaza con \texttt{403\ Forbidden} y
  cuerpo \texttt{Problem\allowbreak{}Detail}.
\item
  \textbf{Público} --- no requiere autenticación
  (\texttt{Security\allowbreak{}Config}, lista \texttt{permit\allowbreak{}All}).
\end{itemize}

\paragraph{2.8.1. Autenticación y
perfil}\label{autenticaciuxf3n-y-perfil}

{\def\LTcaptype{none} % do not increment counter
\begin{longtable}[]{@{}lllll@{}}
\toprule\noalign{}
Operación (endpoint) & ADMIN & VETERINARIO & AUXILIAR & DUENO \\
\midrule\noalign{}
\endhead
\bottomrule\noalign{}
\endlastfoot
\texttt{POST\ /\allowbreak{}api/\allowbreak{}auth/\allowbreak{}registro} & Público & Público & Público &
Público \\
\texttt{POST\ /\allowbreak{}api/\allowbreak{}auth/\allowbreak{}login} & Público & Público & Público &
Público \\
\texttt{POST\ /\allowbreak{}api/\allowbreak{}auth/\allowbreak{}refresh} & Público & Público & Público &
Público \\
\texttt{POST\ /\allowbreak{}api/\allowbreak{}auth/\allowbreak{}logout} & Público & Público & Público &
Público \\
\texttt{GET\ /\allowbreak{}api/\allowbreak{}usuarios/\allowbreak{}me} & Sí & Sí & Sí & Sí \\
\end{longtable}
}

\paragraph{2.8.2. Usuarios
(administración)}\label{usuarios-administraciuxf3n}

{\def\LTcaptype{none} % do not increment counter
\begin{longtable}[]{@{}lllll@{}}
\toprule\noalign{}
Operación (endpoint) & ADMIN & VETERINARIO & AUXILIAR & DUENO \\
\midrule\noalign{}
\endhead
\bottomrule\noalign{}
\endlastfoot
\texttt{GET\ /\allowbreak{}api/\allowbreak{}usuarios} & Sí & No & No & No \\
\texttt{GET\ /\allowbreak{}api/\allowbreak{}usuarios/\allowbreak{}\{id\}} & Sí & No & No & No \\
\texttt{POST\ /\allowbreak{}api/\allowbreak{}usuarios} & Sí & No & No & No \\
\texttt{PUT\ /\allowbreak{}api/\allowbreak{}usuarios/\allowbreak{}\{id\}} & Sí & No & No & No \\
\texttt{DELETE\ /\allowbreak{}api/\allowbreak{}usuarios/\allowbreak{}\{id\}} & Sí & No & No & No \\
\end{longtable}
}

Regla adicional de datos: un \texttt{ROLE\_\allowbreak{}ADMIN} \textbf{no} puede
modificar el rol de su propia cuenta (\texttt{Usuario\allowbreak{}Service}, prueba
\texttt{admin\allowbreak{}No\allowbreak{}Puede\allowbreak{}Escalar\allowbreak{}Su\allowbreak{}Propio\allowbreak{}Rol\allowbreak{}Devuelve403}).

\paragraph{2.8.3. Mascotas}\label{mascotas}

{\def\LTcaptype{none} % do not increment counter
\begin{longtable}[]{@{}lllll@{}}
\toprule\noalign{}
Operación (endpoint) & ADMIN & VETERINARIO & AUXILIAR & DUENO \\
\midrule\noalign{}
\endhead
\bottomrule\noalign{}
\endlastfoot
\texttt{GET\ /\allowbreak{}api/\allowbreak{}mascotas} & Sí & Sí & Sí & Solo propias \\
\texttt{GET\ /\allowbreak{}api/\allowbreak{}mascotas/\allowbreak{}\{id\}} & Sí & Sí & Sí & Solo propias \\
\texttt{POST\ /\allowbreak{}api/\allowbreak{}mascotas} & Sí & Sí & Sí & No \\
\texttt{PUT\ /\allowbreak{}api/\allowbreak{}mascotas/\allowbreak{}\{id\}} & Sí & Sí & Sí & No \\
\texttt{DELETE\ /\allowbreak{}api/\allowbreak{}mascotas/\allowbreak{}\{id\}} & Sí & Sí & Sí & No \\
\end{longtable}
}

El resumen agregado \texttt{GET\ /\allowbreak{}api/\allowbreak{}mascotas/\allowbreak{}resumen-\allowbreak{}especies} está
autorizado para los cuatro roles y se rige por una regla de datos
propia: el rol \texttt{ROLE\_\allowbreak{}ADMIN} obtiene el resumen global, o el de
un dueño concreto si envía el parámetro \texttt{duenio\allowbreak{}Id}; para
cualquier otro rol el parámetro se ignora y el resumen se calcula
siempre sobre el propio usuario autenticado
(\texttt{Mascota\allowbreak{}Service.\allowbreak{}resumen\allowbreak{}Por\allowbreak{}Especie}). La consulta de especies
externas y los endpoints de operación se detallan en 2.8.7.

\paragraph{2.8.4. Citas}\label{citas}

{\def\LTcaptype{none} % do not increment counter
\begin{longtable}[]{@{}lllll@{}}
\toprule\noalign{}
Operación (endpoint) & ADMIN & VETERINARIO & AUXILIAR & DUENO \\
\midrule\noalign{}
\endhead
\bottomrule\noalign{}
\endlastfoot
\texttt{GET\ /\allowbreak{}api/\allowbreak{}citas} & Sí & Sí & Sí & Solo propias \\
\texttt{GET\ /\allowbreak{}api/\allowbreak{}citas/\allowbreak{}\{id\}} & Sí & Sí & Sí & Solo propias \\
\texttt{POST\ /\allowbreak{}api/\allowbreak{}citas} & Sí & No & Sí & No \\
\texttt{PUT\ /\allowbreak{}api/\allowbreak{}citas/\allowbreak{}\{id\}} & Sí & Solo las asignadas a él & Sí &
No \\
\texttt{DELETE\ /\allowbreak{}api/\allowbreak{}citas/\allowbreak{}\{id\}} & Sí & No & No & No \\
\end{longtable}
}

\paragraph{2.8.5. Consultas médicas}\label{consultas-muxe9dicas}

{\def\LTcaptype{none} % do not increment counter
\begin{longtable}[]{@{}lllll@{}}
\toprule\noalign{}
Operación (endpoint) & ADMIN & VETERINARIO & AUXILIAR & DUENO \\
\midrule\noalign{}
\endhead
\bottomrule\noalign{}
\endlastfoot
\texttt{GET\ /\allowbreak{}api/\allowbreak{}consultas} & Sí & Sí & Sí & Solo propias \\
\texttt{GET\ /\allowbreak{}api/\allowbreak{}consultas/\allowbreak{}\{id\}} & Sí & Sí & Sí & Solo propias \\
\texttt{POST\ /\allowbreak{}api/\allowbreak{}consultas} & Sí & Sí & Sí & No \\
\texttt{PUT\ /\allowbreak{}api/\allowbreak{}consultas/\allowbreak{}\{id\}} & Sí & Sí & Sí & No \\
\texttt{DELETE\ /\allowbreak{}api/\allowbreak{}consultas/\allowbreak{}\{id\}} & Sí & Sí & Sí & No \\
\end{longtable}
}

\begin{quote}
\textbf{Defecto abierto declarado.} La celda ``Solo propias'' de
\texttt{GET\ /\allowbreak{}api/\allowbreak{}consultas} describe la regla \textbf{especificada}, no
la implementada: \texttt{Consulta\allowbreak{}Service.\allowbreak{}listar()} devuelve
\texttt{consulta\allowbreak{}Repository.\allowbreak{}find\allowbreak{}All\allowbreak{}By\allowbreak{}Activo\allowbreak{}True(pageable)} también para
\texttt{ROLE\_\allowbreak{}DUENO}, sin filtrar por propietario. Ver REQ-F-006,
REQ-F-013, REQ-NF-015 y la sección 7. No se corrige código en esta
revisión documental.
\end{quote}

\paragraph{2.8.6. Vacunas}\label{vacunas}

{\def\LTcaptype{none} % do not increment counter
\begin{longtable}[]{@{}lllll@{}}
\toprule\noalign{}
Operación (endpoint) & ADMIN & VETERINARIO & AUXILIAR & DUENO \\
\midrule\noalign{}
\endhead
\bottomrule\noalign{}
\endlastfoot
\texttt{GET\ /\allowbreak{}api/\allowbreak{}vacunas} & Sí & Sí & Sí & Solo propias \\
\texttt{GET\ /\allowbreak{}api/\allowbreak{}vacunas/\allowbreak{}mascota/\allowbreak{}\{mascota\allowbreak{}Id\}} & Sí & Sí & Sí & Solo
propias \\
\texttt{GET\ /\allowbreak{}api/\allowbreak{}vacunas/\allowbreak{}\{id\}} & Sí & Sí & Sí & Solo propias \\
\texttt{POST\ /\allowbreak{}api/\allowbreak{}vacunas} & Sí & Sí & Sí & No \\
\texttt{PUT\ /\allowbreak{}api/\allowbreak{}vacunas/\allowbreak{}\{id\}} & Sí & Sí & Sí & No \\
\texttt{DELETE\ /\allowbreak{}api/\allowbreak{}vacunas/\allowbreak{}\{id\}} & Sí & Sí & Sí & No \\
\end{longtable}
}

\paragraph{2.8.7. Información externa de especies y
operación}\label{informaciuxf3n-externa-de-especies-y-operaciuxf3n}

{\def\LTcaptype{none} % do not increment counter
\begin{longtable}[]{@{}lllll@{}}
\toprule\noalign{}
Operación (endpoint) & ADMIN & VETERINARIO & AUXILIAR & DUENO \\
\midrule\noalign{}
\endhead
\bottomrule\noalign{}
\endlastfoot
\texttt{GET\ /\allowbreak{}api/\allowbreak{}externa/\allowbreak{}especies} & Sí & Sí & Sí & Sí \\
\texttt{GET\ /\allowbreak{}actuator/\allowbreak{}health} & Público & Público & Público &
Público \\
\texttt{GET\ /\allowbreak{}api/\allowbreak{}docs} & Público & Público & Público & Público \\
\texttt{GET\ /\allowbreak{}api/\allowbreak{}openapi} & Público & Público & Público & Público \\
\end{longtable}
}

\subsubsection{2.9. Estados principales del dominio y
transiciones}\label{estados-principales-del-dominio-y-transiciones}

Estados tomados directamente de las entidades y migraciones del
repositorio. No se documentan estados que el código no tenga.

\paragraph{\texorpdfstring{2.9.1. Cita
(\texttt{com.\allowbreak{}biopet.\allowbreak{}entity.\allowbreak{}Estado\allowbreak{}Cita})}{2.9.1. Cita (com.biopet.entity.EstadoCita)}}\label{cita-com.biopet.entity.estadocita}

Valores permitidos: \texttt{PROGRAMADA}, \texttt{CANCELADA},
\texttt{COMPLETADA} (restricción \texttt{CHECK} replicada en
\texttt{sp\_\allowbreak{}actualizar\_\allowbreak{}esta\allowbreak{}do\_\allowbreak{}citas\_\allowbreak{}masivas}).

{\def\LTcaptype{none} % do not increment counter
\begin{longtable}[]{@{}
  >{\raggedright\arraybackslash}p{(\linewidth - 4\tabcolsep) * \real{0.3333}}
  >{\raggedright\arraybackslash}p{(\linewidth - 4\tabcolsep) * \real{0.3333}}
  >{\raggedright\arraybackslash}p{(\linewidth - 4\tabcolsep) * \real{0.3333}}@{}}
\toprule\noalign{}
\begin{minipage}[b]{\linewidth}\raggedright
Desde
\end{minipage} & \begin{minipage}[b]{\linewidth}\raggedright
Hacia
\end{minipage} & \begin{minipage}[b]{\linewidth}\raggedright
Disparador real
\end{minipage} \\
\midrule\noalign{}
\endhead
\bottomrule\noalign{}
\endlastfoot
\emph{(inexistente)} & \texttt{PROGRAMADA} & \texttt{POST\ /\allowbreak{}api/\allowbreak{}citas}.
\texttt{Cita\allowbreak{}Service.\allowbreak{}crear} fuerza \texttt{Estado\allowbreak{}Cita.\allowbreak{}PROGRAMADA} e
\textbf{ignora} cualquier estado enviado por el cliente (prueba
\texttt{crear\allowbreak{}Cita\allowbreak{}Forzando\allowbreak{}Estado\allowbreak{}Distinto\allowbreak{}Ignora\allowbreak{}El\allowbreak{}Estado\allowbreak{}Enviado}). \\
\texttt{PROGRAMADA} & \texttt{COMPLETADA} &
\texttt{PUT\ /\allowbreak{}api/\allowbreak{}citas/\allowbreak{}\{id\}} con \texttt{estado=COMPLETAD\allowbreak{}A}, o
\texttt{CALL\ sp\_\allowbreak{}actualizar\_\allowbreak{}e\allowbreak{}stado\_\allowbreak{}citas\_\allowbreak{}masiva\allowbreak{}s(...)} para cierre
masivo por veterinario y fecha límite. \\
\texttt{PROGRAMADA} & \texttt{CANCELADA} &
\texttt{PUT\ /\allowbreak{}api/\allowbreak{}citas/\allowbreak{}\{id\}} con \texttt{estado=CANCELADA}, o el
mismo procedimiento masivo. \\
cualquiera & \emph{(baja lógica)} & \texttt{DELETE\ /\allowbreak{}api/\allowbreak{}citas/\allowbreak{}\{id\}}
marca \texttt{activo\ =\ false}; el valor de \texttt{estado} no
cambia. \\
\end{longtable}
}

El procedimiento masivo valida ambos estados contra la lista permitida y
lanza \texttt{RAISE\ EXCEPTION} si alguno es inválido (prueba
\texttt{actualizar\allowbreak{}Estado\allowbreak{}Citas\allowbreak{}Masivas\_\allowbreak{}estado\allowbreak{}Invalido\_\allowbreak{}lanza\allowbreak{}Excepcion}).

\paragraph{2.9.2. Ciclo de vida común de los recursos de
negocio}\label{ciclo-de-vida-comuxfan-de-los-recursos-de-negocio}

\texttt{Usuario}, \texttt{Mascota}, \texttt{Cita}, \texttt{Consulta} y
\texttt{Vacuna} comparten un único atributo de ciclo de vida:
\texttt{activo} (booleano). No existe borrado físico en ninguna de las
cinco entidades.

{\def\LTcaptype{none} % do not increment counter
\begin{longtable}[]{@{}
  >{\raggedright\arraybackslash}p{(\linewidth - 4\tabcolsep) * \real{0.3333}}
  >{\raggedright\arraybackslash}p{(\linewidth - 4\tabcolsep) * \real{0.3333}}
  >{\raggedright\arraybackslash}p{(\linewidth - 4\tabcolsep) * \real{0.3333}}@{}}
\toprule\noalign{}
\begin{minipage}[b]{\linewidth}\raggedright
Desde
\end{minipage} & \begin{minipage}[b]{\linewidth}\raggedright
Hacia
\end{minipage} & \begin{minipage}[b]{\linewidth}\raggedright
Disparador real
\end{minipage} \\
\midrule\noalign{}
\endhead
\bottomrule\noalign{}
\endlastfoot
\emph{(inexistente)} & \texttt{activo\ =\ true} & Operación
\texttt{POST} del recurso correspondiente. \\
\texttt{activo\ =\ true} & \texttt{activo\ =\ false} & Operación
\texttt{DELETE} del recurso (baja lógica). \\
\texttt{activo\ =\ false} & --- & No hay operación de reactivación
implementada en la API. \\
\end{longtable}
}

Las consultas de lectura filtran siempre por \texttt{activo\allowbreak{}True}, de
modo que un registro dado de baja deja de aparecer en listados y
búsquedas por identificador (responde \texttt{404} con
\texttt{Problem\allowbreak{}Detail}). El trigger \texttt{set\_\allowbreak{}actualizado\_\allowbreak{}en}
mantiene \texttt{actualizado\_\allowbreak{}en} en cada \texttt{UPDATE}.

\paragraph{2.9.3. Sesión y token JWT}\label{sesiuxf3n-y-token-jwt}

{\def\LTcaptype{none} % do not increment counter
\begin{longtable}[]{@{}
  >{\raggedright\arraybackslash}p{(\linewidth - 4\tabcolsep) * \real{0.3333}}
  >{\raggedright\arraybackslash}p{(\linewidth - 4\tabcolsep) * \real{0.3333}}
  >{\raggedright\arraybackslash}p{(\linewidth - 4\tabcolsep) * \real{0.3333}}@{}}
\toprule\noalign{}
\begin{minipage}[b]{\linewidth}\raggedright
Desde
\end{minipage} & \begin{minipage}[b]{\linewidth}\raggedright
Hacia
\end{minipage} & \begin{minipage}[b]{\linewidth}\raggedright
Disparador real
\end{minipage} \\
\midrule\noalign{}
\endhead
\bottomrule\noalign{}
\endlastfoot
\emph{(sin sesión)} & access token vigente + refresh token vigente &
\texttt{POST\ /\allowbreak{}api/\allowbreak{}auth/\allowbreak{}login}. \\
access token vigente & access token renovado &
\texttt{POST\ /\allowbreak{}api/\allowbreak{}auth/\allowbreak{}refresh} con un refresh token válido y no
revocado. \\
access token vigente & revocado & \texttt{POST\ /\allowbreak{}api/\allowbreak{}auth/\allowbreak{}logout}:
\texttt{Token\allowbreak{}Blacklist\allowbreak{}Service.\allowbreak{}revoke} escribe
\texttt{jwt:blacklist:\textless{}\allowbreak{}jti\textgreater{}} en Redis con TTL
igual al tiempo restante del token. \\
access token vigente & expirado & Transcurre
\texttt{JWT\_\allowbreak{}EXPIRATION\_\allowbreak{}MS} (1 h por defecto). \\
revocado o expirado & rechazo & Cualquier solicitud a recurso protegido
responde \texttt{401} con \texttt{Problem\allowbreak{}Detail}. \\
\end{longtable}
}

\begin{center}\rule{0.5\linewidth}{0.5pt}\end{center}

\subsection{3. Requisitos específicos}\label{requisitos-especuxedficos}

Cada requisito sigue el patrón
\texttt{{[}condición{]}\ {[}\allowbreak{}sujeto{]}\ deberá\ {[}\allowbreak{}acción{]}\ {[}obje\allowbreak{}to{]}\ {[}restricc\allowbreak{}ión{]}}
de ISO/IEC/IEEE 29148, con identificador persistente, rationale,
prioridad MoSCoW, criterios de aceptación verificables y método de
verificación, cumpliendo las características INCOSE C1--C9 (Necessary,
Appropriate, Unambiguous, Complete, Singular, Feasible, Verifiable,
Correct, Conforming). La plantilla de campos es la de la sección 1.3.2 y
los valores permitidos del campo Estado son exclusivamente los tres de
la sección 1.3.1.

\subsubsection{3.1. Requisitos funcionales
(REQ-F)}\label{requisitos-funcionales-req-f}

\begin{quote}
\textbf{Nota de consistencia de numeración.} La numeración de esta
sección es la que ya está fijada en
\texttt{docs/\allowbreak{}requisitos/\allowbreak{}historias/\allowbreak{}Historias\allowbreak{}Usuario.\allowbreak{}md} (HU-001 a HU-024)
y \texttt{docs/\allowbreak{}requisitos/\allowbreak{}casos-\allowbreak{}de-\allowbreak{}uso/\allowbreak{}Casos\allowbreak{}De\allowbreak{}Uso.\allowbreak{}md} (CU-01 a CU-24).
Esta versión del SRS adopta esa numeración como única fuente de verdad,
para que los tres documentos no se contradigan entre sí. \textbf{No
existe ningún identificador con sufijo de letra:} los tres requisitos
que agrupan dos obligaciones (\texttt{REQ-\allowbreak{}F-\allowbreak{}013}, \texttt{REQ-\allowbreak{}F-\allowbreak{}015} y
\texttt{REQ-\allowbreak{}NF-\allowbreak{}012}) las separan dentro de su propio bloque como
\emph{Obligación A} y \emph{Obligación B}, conservando un identificador
único cada uno (ver sección 1.3.2).
\end{quote}

\subsubsection{REQ-F-001 --- Registro de usuario dueño de
mascota}\label{req-f-001-registro-de-usuario-dueuxf1o-de-mascota}

\begin{itemize}
\tightlist
\item
  \textbf{Tipo:} Funcional
\item
  \textbf{Prioridad:} Must
\item
  \textbf{Enunciado:} El sistema deberá permitir que cualquier visitante
  no autenticado se registre proporcionando nombre, correo electrónico y
  contraseña, asignándole automáticamente el rol \texttt{ROLE\_\allowbreak{}DUENO} e
  ignorando cualquier rol enviado por el cliente.
\item
  \textbf{Rationale:} heredado de RF-01 de la Entrega 1A; el registro
  público solo aplica al rol dueño, los demás roles se crean
  administrativamente (REQ-F-023).
\item
  \textbf{Criterios de aceptación:}

  \begin{enumerate}
  \def\labelenumi{\arabic{enumi}.}
  \tightlist
  \item
    Dado un visitante no autenticado, cuando envía
    \texttt{POST\ /\allowbreak{}api/\allowbreak{}auth/\allowbreak{}registro} con nombre, correo no registrado y
    contraseña válidos, entonces la respuesta es \texttt{201\ Created} y
    el cuerpo contiene \texttt{id}, \texttt{nombre}, \texttt{email} y
    \texttt{rol}, sin ningún campo de contraseña ni hash.
  \item
    Dado un cuerpo de registro que incluye \texttt{"rol":"ROLE\_\allowbreak{}ADMIN"}
    (o cualquier otro rol), cuando se procesa el registro, entonces el
    usuario creado queda con \texttt{rol\ =\ ROLE\_\allowbreak{}DUENO} y
    \texttt{activo\ =\ true}.
  \item
    Dado un correo con mayúsculas, cuando se registra, entonces el
    correo se persiste en minúsculas.
  \item
    Dado un cuerpo con nombre vacío, correo con formato inválido o
    contraseña de menos de 8 caracteres, cuando se envía la solicitud,
    entonces la respuesta es \texttt{422\ Unprocessable\ E\allowbreak{}ntity} con
    \texttt{Problem\allowbreak{}Detail} de tipo \texttt{urn:biopet:error:\allowbreak{}validation}
    y el objeto \texttt{errors} enumera los campos inválidos.
  \end{enumerate}
\item
  \textbf{Método de verificación:} demostración con captura HTTP real
  --- \texttt{docs/\allowbreak{}mediciones/\allowbreak{}sec/\allowbreak{}raw/\allowbreak{}A01-\allowbreak{}access-\allowbreak{}control.\allowbreak{}txt} (líneas
  9--11) registra dos registros reales con \texttt{status\_\allowbreak{}registro=2\allowbreak{}01}
  y el rol forzado a \texttt{ROLE\_\allowbreak{}DUENO} pese al valor enviado,
  generados por \texttt{scripts/\allowbreak{}security-\allowbreak{}evidence.\allowbreak{}sh}. Para el criterio
  4, prueba automatizada
  \texttt{Auth\allowbreak{}Controller\allowbreak{}Test.\allowbreak{}registro\allowbreak{}Con\allowbreak{}Campos\allowbreak{}Invalidos}. Para los
  criterios 2 y 3, inspección de \texttt{Auth\allowbreak{}Service.\allowbreak{}registrar}
  (\texttt{rol(Rol.\allowbreak{}ROLE\_\allowbreak{}DUENO)} y
  \texttt{email(request.\allowbreak{}email().\allowbreak{}to\allowbreak{}Lower\allowbreak{}Case())}).
\item
  \textbf{Trazabilidad:} HU-001 → CU-01 →
  \texttt{Auth\allowbreak{}Controller.\allowbreak{}registro} → \texttt{Auth\allowbreak{}Service.\allowbreak{}registrar} →
  \texttt{POST\ /\allowbreak{}api/\allowbreak{}auth/\allowbreak{}registro}; DTO
  \texttt{Registro\allowbreak{}Request}/\texttt{Usuario\allowbreak{}Response}; evidencia
  \texttt{docs/\allowbreak{}mediciones/\allowbreak{}sec/\allowbreak{}A01-\allowbreak{}access-\allowbreak{}control.\allowbreak{}md} y
  \texttt{docs/\allowbreak{}mediciones/\allowbreak{}sec/\allowbreak{}raw/\allowbreak{}A01-\allowbreak{}access-\allowbreak{}control.\allowbreak{}txt}.
\item
  \textbf{Estado:} verificado
\item
  \textbf{Observaciones:} no existe prueba JUnit dedicada al camino
  feliz del registro; la evidencia del criterio 1 es la captura HTTP
  real citada, no una prueba automatizada. Una versión anterior de este
  documento citaba un test ``\texttt{Auth\allowbreak{}Controller\allowbreak{}Test} (registro
  exitoso)'' que no existe en
  \texttt{Backend/\allowbreak{}src/\allowbreak{}test/\allowbreak{}java/\allowbreak{}com/\allowbreak{}biopet/\allowbreak{}Auth\allowbreak{}Controller\allowbreak{}Test.\allowbreak{}java}; la
  referencia se corrige aquí.
\end{itemize}

\subsubsection{REQ-F-002 --- Rechazo de registro con correo
duplicado}\label{req-f-002-rechazo-de-registro-con-correo-duplicado}

\begin{itemize}
\tightlist
\item
  \textbf{Tipo:} Funcional
\item
  \textbf{Prioridad:} Must
\item
  \textbf{Enunciado:} Al recibir una solicitud de registro con un correo
  electrónico ya existente, el sistema deberá rechazarla con código
  \texttt{409\ Conflict} y un cuerpo \texttt{Problem\allowbreak{}Detail}, sin crear
  ningún registro nuevo.
\item
  \textbf{Rationale:} funcionalidad presente en el código
  (\texttt{Email\allowbreak{}Duplicado\allowbreak{}Exception}) no documentada como requisito
  independiente en el SRS original de la Entrega 1A; se agrega para
  cerrar el hueco (bloque A.3 de la Guía).
\item
  \textbf{Criterios de aceptación:}

  \begin{enumerate}
  \def\labelenumi{\arabic{enumi}.}
  \tightlist
  \item
    Dado un correo previamente registrado, cuando se intenta registrar
    de nuevo ese mismo correo, entonces la respuesta es
    \texttt{409\ Conflict}.
  \item
    El cuerpo de esa respuesta tiene
    \texttt{Content-\allowbreak{}Type:\ application/\allowbreak{}problem+json} y los campos
    \texttt{type\ =\ urn:biopet:\allowbreak{}error:conflict},
    \texttt{title\ =\ "Conflicto\ d\allowbreak{}e\ datos"}, \texttt{status\ =\ 409},
    \texttt{detail} no vacío e \texttt{instance\ =\ /\allowbreak{}api/\allowbreak{}auth/\allowbreak{}registro}.
  \item
    La cantidad de usuarios con ese correo permanece igual: no se crea
    un segundo usuario.
  \end{enumerate}
\item
  \textbf{Método de verificación:} prueba automatizada
  \texttt{Auth\allowbreak{}Controller\allowbreak{}Test.\allowbreak{}registro\allowbreak{}Email\allowbreak{}Duplicado} (criterios 1 y 2,
  con aserciones sobre los cinco campos del \texttt{Problem\allowbreak{}Detail}).
  Criterio 3: inspección de \texttt{Auth\allowbreak{}Service.\allowbreak{}registrar}, donde la
  comprobación \texttt{usuario\allowbreak{}Repository.\allowbreak{}exists\allowbreak{}By\allowbreak{}Email(...)} lanza
  \texttt{Email\allowbreak{}Duplicado\allowbreak{}Exception} antes de cualquier invocación a
  \texttt{usuario\allowbreak{}Repository.\allowbreak{}save(...)}.
\item
  \textbf{Trazabilidad:} HU-001 → CU-01 → \texttt{Auth\allowbreak{}Service.\allowbreak{}registrar}
  → \texttt{Email\allowbreak{}Duplicado\allowbreak{}Exception} →
  \texttt{Global\allowbreak{}Exception\allowbreak{}Handler.\allowbreak{}email\allowbreak{}Duplicado} →
  \texttt{Problem\allowbreak{}Type.\allowbreak{}CONFLICT} → \texttt{POST\ /\allowbreak{}api/\allowbreak{}auth/\allowbreak{}registro}.
\item
  \textbf{Estado:} verificado
\end{itemize}

\subsubsection{REQ-F-003 --- Autenticación mediante usuario y
contraseña}\label{req-f-003-autenticaciuxf3n-mediante-usuario-y-contraseuxf1a}

\begin{itemize}
\tightlist
\item
  \textbf{Tipo:} Funcional
\item
  \textbf{Prioridad:} Must
\item
  \textbf{Enunciado:} El sistema deberá permitir que un usuario
  registrado y activo inicie sesión mediante correo electrónico y
  contraseña, emitiendo un access token con vigencia de 1 hora y un
  refresh token, ambos con los siete claims JWT estándar del RFC 7519.
\item
  \textbf{Rationale:} heredado de RF-16/RF-WEB-01 de la Entrega 1A.
\item
  \textbf{Criterios de aceptación:}

  \begin{enumerate}
  \def\labelenumi{\arabic{enumi}.}
  \tightlist
  \item
    Dadas credenciales válidas, cuando se envía
    \texttt{POST\ /\allowbreak{}api/\allowbreak{}auth/\allowbreak{}login}, entonces la respuesta es
    \texttt{200\ OK} y establece las cookies \texttt{access\_\allowbreak{}token} y
    \texttt{refresh\_\allowbreak{}token} con atributos \texttt{Http\allowbreak{}Only},
    \texttt{Secure} y \texttt{Same\allowbreak{}Site=Strict}.
  \item
    Cada token emitido contiene exactamente los siete claims
    \texttt{iss}, \texttt{sub}, \texttt{aud}, \texttt{exp},
    \texttt{nbf}, \texttt{iat}, \texttt{jti}, con \texttt{nbf\ ==\ iat}
    y \texttt{jti} distinto entre dos tokens emitidos consecutivamente.
  \item
    Dadas credenciales inválidas, cuando se envía la solicitud, entonces
    la respuesta es \texttt{401\ Unauthorized} con
    \texttt{Problem\allowbreak{}Detail} de tipo
    \texttt{urn:biopet:error:\allowbreak{}unauthorized} y no se emite ningún token.
  \item
    Un token con \texttt{iss} o \texttt{aud} distintos de los
    configurados es rechazado.
  \end{enumerate}
\item
  \textbf{Método de verificación:} pruebas automatizadas
  \texttt{Auth\allowbreak{}Controller\allowbreak{}Test.\allowbreak{}login\allowbreak{}Exitoso} y
  \texttt{Auth\allowbreak{}Controller\allowbreak{}Test.\allowbreak{}login\allowbreak{}Clave\allowbreak{}Incorrecta} (criterios 1 y 3);
  \texttt{Jwt\allowbreak{}Service\allowbreak{}Test.\allowbreak{}access\allowbreak{}Token\allowbreak{}Contiene\allowbreak{}Los\allowbreak{}Siete\allowbreak{}Claims},
  \texttt{refresh\allowbreak{}Token\allowbreak{}Contiene\allowbreak{}Los\allowbreak{}Siete\allowbreak{}Claims}, \texttt{nbf\allowbreak{}Es\allowbreak{}Igual\allowbreak{}AIat},
  \texttt{jti\allowbreak{}Es\allowbreak{}Unico\allowbreak{}Entre\allowbreak{}Tokens} (criterio 2);
  \texttt{Jwt\allowbreak{}Service\allowbreak{}Test.\allowbreak{}issuer\allowbreak{}Incorrecto\allowbreak{}Es\allowbreak{}Rechazado} y
  \texttt{audience\allowbreak{}Incorrecta\allowbreak{}Es\allowbreak{}Rechazada} (criterio 4).
\item
  \textbf{Trazabilidad:} HU-002 → CU-02 → \texttt{Auth\allowbreak{}Controller.\allowbreak{}login}
  → \texttt{Auth\allowbreak{}Service.\allowbreak{}login} → \texttt{Jwt\allowbreak{}Service.\allowbreak{}build\allowbreak{}Token} →
  \texttt{Jwt\allowbreak{}Cookie\allowbreak{}Service} → \texttt{POST\ /\allowbreak{}api/\allowbreak{}auth/\allowbreak{}login};
  configuración \texttt{security.\allowbreak{}jwt.*} en
  \texttt{Backend/\allowbreak{}src/\allowbreak{}main/\allowbreak{}resources/\allowbreak{}application.\allowbreak{}yml}.
\item
  \textbf{Estado:} verificado
\end{itemize}

\subsubsection{REQ-F-004 --- Renovación de sesión
(refresh)}\label{req-f-004-renovaciuxf3n-de-sesiuxf3n-refresh}

\begin{itemize}
\tightlist
\item
  \textbf{Tipo:} Funcional
\item
  \textbf{Prioridad:} Must
\item
  \textbf{Enunciado:} El sistema deberá permitir renovar el access token
  utilizando un refresh token válido y no revocado, sin exigir
  nuevamente las credenciales del usuario.
\item
  \textbf{Rationale:} heredado de RNF-03/RNF-WEB-03 de la Entrega 1A
  (expiración configurable de JWT), llevado a requisito funcional
  explícito.
\item
  \textbf{Criterios de aceptación:}

  \begin{enumerate}
  \def\labelenumi{\arabic{enumi}.}
  \tightlist
  \item
    Dada una cookie \texttt{refresh\_\allowbreak{}token} válida y no revocada, cuando
    se envía \texttt{POST\ /\allowbreak{}api/\allowbreak{}auth/\allowbreak{}refresh}, entonces la respuesta es
    \texttt{200\ OK} y se emite una nueva cookie \texttt{access\_\allowbreak{}token}.
  \item
    Dada la ausencia de cookie de refresh, entonces la respuesta es
    \texttt{401} con \texttt{Problem\allowbreak{}Detail}.
  \item
    Dada una cookie de refresh inválida o manipulada, entonces la
    respuesta es \texttt{401} con \texttt{Problem\allowbreak{}Detail}.
  \item
    Un access token presentado como refresh token no sirve para renovar:
    la respuesta es \texttt{401}.
  \item
    Un refresh token previamente revocado por logout devuelve
    \texttt{401}.
  \end{enumerate}
\item
  \textbf{Método de verificación:} pruebas automatizadas
  \texttt{Auth\allowbreak{}Controller\allowbreak{}Test.\allowbreak{}refresh\allowbreak{}Cookie\allowbreak{}Valida\allowbreak{}Emite\allowbreak{}Nueva\allowbreak{}Access\allowbreak{}Cookie},
  \texttt{refresh\allowbreak{}Sin\allowbreak{}Cookie\allowbreak{}Devuelve401\allowbreak{}Problem\allowbreak{}Detail},
  \texttt{refresh\allowbreak{}Cookie\allowbreak{}Invalida\allowbreak{}Devuelve401\allowbreak{}Problem\allowbreak{}Detail},
  \texttt{access\allowbreak{}Token\allowbreak{}No\allowbreak{}Sirve\allowbreak{}Como\allowbreak{}Refresh\allowbreak{}Cookie} y
  \texttt{refresh\allowbreak{}Cookie\allowbreak{}Revocada\allowbreak{}Devuelve401} (un criterio por prueba, en
  ese orden).
\item
  \textbf{Trazabilidad:} HU-003 → CU-03 →
  \texttt{Auth\allowbreak{}Controller.\allowbreak{}refresh} → \texttt{Auth\allowbreak{}Service.\allowbreak{}refresh} →
  \texttt{Jwt\allowbreak{}Service.\allowbreak{}is\allowbreak{}Access\allowbreak{}Token} →
  \texttt{Token\allowbreak{}Blacklist\allowbreak{}Service.\allowbreak{}is\allowbreak{}Revoked} →
  \texttt{POST\ /\allowbreak{}api/\allowbreak{}auth/\allowbreak{}refresh}.
\item
  \textbf{Estado:} verificado
\end{itemize}

\subsubsection{REQ-F-005 --- Cierre de sesión con revocación de
token}\label{req-f-005-cierre-de-sesiuxf3n-con-revocaciuxf3n-de-token}

\begin{itemize}
\tightlist
\item
  \textbf{Tipo:} Funcional
\item
  \textbf{Prioridad:} Must
\item
  \textbf{Enunciado:} El sistema deberá permitir cerrar sesión
  registrando el \texttt{jti} de cada token vigente en una lista negra
  en Redis, con TTL igual a su tiempo restante de expiración, de modo
  que una solicitud posterior con ese token a un recurso protegido
  responda \texttt{401}.
\item
  \textbf{Rationale:} heredado de RF-17/RF-WEB-04 de la Entrega 1A;
  mecanismo de revocación definido en
  \texttt{docs/\allowbreak{}adr/\allowbreak{}ADR-\allowbreak{}003-\allowbreak{}jwt-\allowbreak{}redis.\allowbreak{}md}.
\item
  \textbf{Criterios de aceptación:}

  \begin{enumerate}
  \def\labelenumi{\arabic{enumi}.}
  \tightlist
  \item
    Dado un usuario autenticado con ambas cookies, cuando envía
    \texttt{POST\ /\allowbreak{}api/\allowbreak{}auth/\allowbreak{}logout}, entonces la respuesta es
    \texttt{204\ No\ Content}, se escriben las claves
    \texttt{jwt:blacklist:\textless{}\allowbreak{}jti\textgreater{}} correspondientes
    y ambas cookies se eliminan del cliente.
  \item
    Dado un token cuyo tiempo restante de expiración es positivo, cuando
    se revoca, entonces la clave se escribe con TTL igual a ese tiempo
    restante; si el tiempo restante es cero o negativo, no se escribe
    nada en Redis.
  \item
    Dado un refresh token revocado, cuando se intenta usar en
    \texttt{POST\ /\allowbreak{}api/\allowbreak{}auth/\allowbreak{}refresh}, entonces la respuesta es
    \texttt{401}.
  \item
    El logout es idempotente: sin cookies, o con cookies inválidas, la
    respuesta sigue siendo \texttt{204}.
  \end{enumerate}
\item
  \textbf{Método de verificación:} pruebas automatizadas
  \texttt{Auth\allowbreak{}Controller\allowbreak{}Test.\allowbreak{}logout\allowbreak{}Con\allowbreak{}Ambas\allowbreak{}Cookies\allowbreak{}Revoca\allowbreak{}Tokens\allowbreak{}YLas\allowbreak{}Elimina},
  \texttt{logout\allowbreak{}Solo\allowbreak{}Con\allowbreak{}Access\allowbreak{}Cookie\allowbreak{}Revoca\allowbreak{}Access\allowbreak{}YElimina\allowbreak{}Ambas},
  \texttt{logout\allowbreak{}Solo\allowbreak{}Con\allowbreak{}Refresh\allowbreak{}Cookie\allowbreak{}Revoca\allowbreak{}Refresh\allowbreak{}YElimina\allowbreak{}Ambas},
  \texttt{logout\allowbreak{}Sin\allowbreak{}Cookies\allowbreak{}Es\allowbreak{}Idempotente},
  \texttt{logout\allowbreak{}Con\allowbreak{}Cookies\allowbreak{}Invalidas\allowbreak{}Sigue\allowbreak{}Siendo204},
  \texttt{refresh\allowbreak{}Cookie\allowbreak{}Revocada\allowbreak{}Devuelve401}; y
  \texttt{Token\allowbreak{}Blacklist\allowbreak{}Service\allowbreak{}Test.\allowbreak{}ttl\allowbreak{}Positivo\allowbreak{}Escribe\allowbreak{}En\allowbreak{}Redis\allowbreak{}Con\allowbreak{}Ttl\allowbreak{}Restante},
  \texttt{ttl\allowbreak{}Negativo\allowbreak{}No\allowbreak{}Escribe\allowbreak{}En\allowbreak{}Redis}, \texttt{ttl\allowbreak{}Cero\allowbreak{}No\allowbreak{}Escribe\allowbreak{}En\allowbreak{}Redis}
  (criterio 2).
\item
  \textbf{Trazabilidad:} HU-004 → CU-04 → \texttt{Auth\allowbreak{}Controller.\allowbreak{}logout}
  → \texttt{Token\allowbreak{}Blacklist\allowbreak{}Service.\allowbreak{}revoke} → clave Redis
  \texttt{jwt:blacklist:\textless{}\allowbreak{}jti\textgreater{}} →
  \texttt{POST\ /\allowbreak{}api/\allowbreak{}auth/\allowbreak{}logout};
  \texttt{docs/\allowbreak{}adr/\allowbreak{}ADR-\allowbreak{}003-\allowbreak{}jwt-\allowbreak{}redis.\allowbreak{}md}.
\item
  \textbf{Estado:} verificado
\item
  \textbf{Observaciones:} el comportamiento del sistema cuando Redis no
  está disponible no forma parte de este requisito y \textbf{no está
  definido hoy}; se especifica por separado en REQ-NF-014, con estado
  \texttt{pendiente}.
\end{itemize}

\subsubsection{REQ-F-006 --- Control de acceso por rol y por propietario
del recurso (RBAC +
aislamiento)}\label{req-f-006-control-de-acceso-por-rol-y-por-propietario-del-recurso-rbac-aislamiento}

\begin{itemize}
\tightlist
\item
  \textbf{Tipo:} Funcional
\item
  \textbf{Prioridad:} Must
\item
  \textbf{Enunciado:} El sistema deberá restringir el acceso a cada
  endpoint protegido según el rol del usuario autenticado, rechazando
  con \texttt{403\ Forbidden} y cuerpo \texttt{Problem\allowbreak{}Detail} toda
  solicitud de un rol no autorizado para ese recurso, y deberá además
  restringir el contenido devuelto a un usuario con rol
  \texttt{ROLE\_\allowbreak{}DUENO} a los recursos asociados a su propia cuenta,
  conforme a la matriz de permisos de la sección 2.8.
\item
  \textbf{Rationale:} heredado de RF-13/RF-WEB-02 de la Entrega 1A. En
  esta revisión el enunciado se \textbf{amplía explícitamente} para
  incluir el aislamiento de datos por propietario: la retroalimentación
  del docente-director señaló que un control que solo comprueba el rol,
  sin comprobar la propiedad del dato, no satisface el propósito del
  requisito. La regla transversal completa se especifica en REQ-NF-015.
\item
  \textbf{Criterios de aceptación:}

  \begin{enumerate}
  \def\labelenumi{\arabic{enumi}.}
  \tightlist
  \item
    Dado un usuario autenticado cuyo rol no está autorizado para una
    operación según la sección 2.8, cuando la invoca, entonces la
    respuesta es \texttt{403\ Forbidden} con
    \texttt{Content-\allowbreak{}Type:\ application/\allowbreak{}problem+json},
    \texttt{type\ =\ urn:biopet:\allowbreak{}error:forbidden} e \texttt{instance}
    igual a la ruta solicitada.
  \item
    Dada una solicitud sin token, entonces la respuesta es \texttt{401}
    con \texttt{Problem\allowbreak{}Detail} de tipo
    \texttt{urn:biopet:error:\allowbreak{}unauthorized}.
  \item
    Dado un usuario con rol \texttt{ROLE\_\allowbreak{}DUENO}, cuando invoca
    \textbf{cualquier} operación de listado marcada ``Solo propias'' en
    la sección 2.8 (\texttt{GET\ /\allowbreak{}api/\allowbreak{}mascotas},
    \texttt{GET\ /\allowbreak{}api/\allowbreak{}citas}, \texttt{GET\ /\allowbreak{}api/\allowbreak{}consultas},
    \texttt{GET\ /\allowbreak{}api/\allowbreak{}vacunas}), entonces \textbf{ningún} elemento
    devuelto pertenece a otro dueño: el dueño nunca puede listar datos
    pertenecientes a otro dueño.
  \item
    Dado un usuario con rol \texttt{ROLE\_\allowbreak{}DUENO}, cuando solicita por
    identificador un recurso de otro dueño, entonces la respuesta es
    \texttt{403}.
  \end{enumerate}
\item
  \textbf{Método de verificación:} pruebas automatizadas
  \texttt{Mascota\allowbreak{}Controller\allowbreak{}Test.\allowbreak{}crear\allowbreak{}Mascota\allowbreak{}Con\allowbreak{}Rol\allowbreak{}Insuficiente\allowbreak{}Devuelve\allowbreak{}Problem\allowbreak{}Detail},
  \texttt{dueno\allowbreak{}Intenta\allowbreak{}Actualizar\allowbreak{}Mascota\allowbreak{}Propia\allowbreak{}Sigue\allowbreak{}Recibiendo403\allowbreak{}Por\allowbreak{}Rol},
  \texttt{dueno\allowbreak{}Intenta\allowbreak{}Eliminar\allowbreak{}Mascota\allowbreak{}Propia\allowbreak{}Sigue\allowbreak{}Recibiendo403\allowbreak{}Por\allowbreak{}Rol},
  \texttt{dueno\allowbreak{}Solo\allowbreak{}Ve\allowbreak{}Sus\allowbreak{}Propias\allowbreak{}Mascotas\allowbreak{}En\allowbreak{}Listado},
  \texttt{dueno\allowbreak{}Consulta\allowbreak{}Mascota\allowbreak{}De\allowbreak{}Otro\allowbreak{}Duenio\allowbreak{}Devuelve403};
  \texttt{Cita\allowbreak{}Controller\allowbreak{}Test.\allowbreak{}dueno\allowbreak{}Solo\allowbreak{}Ve\allowbreak{}Sus\allowbreak{}Propias\allowbreak{}Citas\allowbreak{}En\allowbreak{}Listado},
  \texttt{dueno\allowbreak{}No\allowbreak{}Puede\allowbreak{}Consultar\allowbreak{}Cita\allowbreak{}De\allowbreak{}Mascota\allowbreak{}Ajena};
  \texttt{Vacuna\allowbreak{}Controller\allowbreak{}Test.\allowbreak{}dueno\allowbreak{}Solo\allowbreak{}Ve\allowbreak{}Vacunas\allowbreak{}De\allowbreak{}Sus\allowbreak{}Propias\allowbreak{}Mascotas},
  \texttt{dueno\allowbreak{}Consulta\allowbreak{}Vacuna\allowbreak{}De\allowbreak{}Mascota\allowbreak{}Ajena\allowbreak{}Devuelve403};
  \texttt{Usuario\allowbreak{}Controller\allowbreak{}Test.\allowbreak{}acceso\allowbreak{}Con\allowbreak{}Rol\allowbreak{}No\allowbreak{}Admin\allowbreak{}Devuelve403},
  \texttt{acceso\allowbreak{}Sin\allowbreak{}Autenticacion\allowbreak{}Devuelve401}; y captura HTTP real
  end-to-end en \texttt{docs/\allowbreak{}mediciones/\allowbreak{}sec/\allowbreak{}A01-\allowbreak{}access-\allowbreak{}control.\allowbreak{}md} con
  \texttt{docs/\allowbreak{}mediciones/\allowbreak{}sec/\allowbreak{}raw/\allowbreak{}A01-\allowbreak{}access-\allowbreak{}control.\allowbreak{}txt} (IDOR entre
  dos dueños reales sobre TLS).
\item
  \textbf{Trazabilidad:} HU-005 → CU-05 → anotaciones
  \texttt{@Pre\allowbreak{}Authorize} de \texttt{Mascota\allowbreak{}Controller},
  \texttt{Cita\allowbreak{}Controller}, \texttt{Consulta\allowbreak{}Controller},
  \texttt{Vacuna\allowbreak{}Controller}, \texttt{Usuario\allowbreak{}Controller},
  \texttt{External\allowbreak{}Api\allowbreak{}Controller} → \texttt{Security\allowbreak{}Config}
  (\texttt{any\allowbreak{}Request().\allowbreak{}authenticated()},
  \texttt{Problem\allowbreak{}Access\allowbreak{}Denied\allowbreak{}Handler},
  \texttt{Problem\allowbreak{}Authentication\allowbreak{}Entry\allowbreak{}Point}) → reglas de propiedad en
  \texttt{Mascota\allowbreak{}Service}, \texttt{Cita\allowbreak{}Service}, \texttt{Vacuna\allowbreak{}Service},
  \texttt{Consulta\allowbreak{}Service} → matriz de permisos de la sección 2.8 →
  evidencia \texttt{docs/\allowbreak{}mediciones/\allowbreak{}sec/\allowbreak{}A01-\allowbreak{}access-\allowbreak{}control.\allowbreak{}md}.
\item
  \textbf{Estado:} implementado
\item
  \textbf{Observaciones:} el criterio 3 \textbf{no se cumple} hoy para
  \texttt{GET\ /\allowbreak{}api/\allowbreak{}consultas}. \texttt{Consulta\allowbreak{}Service.\allowbreak{}listar()}
  (\texttt{Backend/\allowbreak{}src/\allowbreak{}main/\allowbreak{}java/\allowbreak{}com/\allowbreak{}biopet/\allowbreak{}service/\allowbreak{}Consulta\allowbreak{}Service.\allowbreak{}java})
  evalúa \texttt{if\ (usuario.\allowbreak{}get\allowbreak{}Rol()\ ==\ Rol.\allowbreak{}ROLE\_\allowbreak{}DUENO)} pero ambas
  ramas devuelven
  \texttt{consulta\allowbreak{}Repository.\allowbreak{}find\allowbreak{}All\allowbreak{}By\allowbreak{}Activo\allowbreak{}True(pageable)}: no hay
  filtro por propietario, y \texttt{Consulta\allowbreak{}Repository} no declara
  ningún método equivalente a
  \texttt{find\allowbreak{}All\allowbreak{}By\allowbreak{}Mascota\_\allowbreak{}Duenio\_\allowbreak{}Id\allowbreak{}And\allowbreak{}Activo\allowbreak{}True} (que sí existe en
  \texttt{Cita\allowbreak{}Repository} y \texttt{Vacuna\allowbreak{}Repository}). En consecuencia,
  un usuario con rol \texttt{ROLE\_\allowbreak{}DUENO} autenticado puede listar
  consultas médicas de mascotas ajenas. Este documento \textbf{no
  modifica código}: el defecto se declara aquí, en REQ-F-013, en
  REQ-NF-015, en la sección 2.8.5 y en la sección 7, y su corrección
  queda como acción de seguimiento para el propietario del módulo. Los
  criterios 1, 2 y 4 sí están verificados por las pruebas citadas.
\end{itemize}

\subsubsection{REQ-F-007 --- Consulta del perfil
propio}\label{req-f-007-consulta-del-perfil-propio}

\begin{itemize}
\tightlist
\item
  \textbf{Tipo:} Funcional
\item
  \textbf{Prioridad:} Should
\item
  \textbf{Enunciado:} El sistema deberá permitir que un usuario
  autenticado, sea cual sea su rol, consulte sus propios datos de perfil
  (identificador, nombre, correo y rol), sin exponerle el listado
  completo de usuarios.
\item
  \textbf{Rationale:} funcionalidad presente en el código
  (\texttt{GET\ /\allowbreak{}api/\allowbreak{}usuarios/\allowbreak{}me}), base del \texttt{auth\allowbreak{}Guard} de
  Angular; no documentada como requisito independiente en el SRS
  original.
\item
  \textbf{Criterios de aceptación:}

  \begin{enumerate}
  \def\labelenumi{\arabic{enumi}.}
  \tightlist
  \item
    Dado un usuario autenticado con cualquier rol, cuando invoca
    \texttt{GET\ /\allowbreak{}api/\allowbreak{}usuarios/\allowbreak{}me}, entonces la respuesta es
    \texttt{200\ OK} con \texttt{id}, \texttt{nombre}, \texttt{email},
    \texttt{rol} y \texttt{activo} del propio usuario.
  \item
    La respuesta no incluye ningún campo de contraseña ni hash.
  \item
    Dado un usuario cuyo rol no es \texttt{ROLE\_\allowbreak{}ADMIN}, cuando invoca
    \texttt{GET\ /\allowbreak{}api/\allowbreak{}usuarios} (listado completo), entonces la
    respuesta es \texttt{403}.
  \item
    Dada una solicitud sin autenticación a
    \texttt{GET\ /\allowbreak{}api/\allowbreak{}usuarios/\allowbreak{}me}, entonces la respuesta es
    \texttt{401} con \texttt{Problem\allowbreak{}Detail}.
  \end{enumerate}
\item
  \textbf{Método de verificación:} pruebas automatizadas
  \texttt{Usuario\allowbreak{}Controller\allowbreak{}Test.\allowbreak{}me\allowbreak{}Sigue\allowbreak{}Funcionando} (criterio 1),
  \texttt{acceso\allowbreak{}Con\allowbreak{}Rol\allowbreak{}No\allowbreak{}Admin\allowbreak{}Devuelve403} (criterio 3),
  \texttt{Auth\allowbreak{}Controller\allowbreak{}Test.\allowbreak{}acceso\allowbreak{}Sin\allowbreak{}Token} (criterio 4). Criterio 2:
  inspección del record \texttt{Usuario\allowbreak{}Response}, que declara únicamente
  \texttt{id}, \texttt{nombre}, \texttt{email}, \texttt{rol} y
  \texttt{activo}.
\item
  \textbf{Trazabilidad:} HU-006 → CU-06 → \texttt{Usuario\allowbreak{}Controller.\allowbreak{}me}
  → \texttt{Auth\allowbreak{}Service.\allowbreak{}perfil} → \texttt{Usuario\allowbreak{}Response} →
  \texttt{GET\ /\allowbreak{}api/\allowbreak{}usuarios/\allowbreak{}me}.
\item
  \textbf{Estado:} verificado
\end{itemize}

\subsubsection{REQ-F-008 --- Creación de
mascota}\label{req-f-008-creaciuxf3n-de-mascota}

\begin{itemize}
\tightlist
\item
  \textbf{Tipo:} Funcional
\item
  \textbf{Prioridad:} Must
\item
  \textbf{Enunciado:} El sistema deberá permitir que un usuario con rol
  \texttt{ROLE\_\allowbreak{}ADMIN}, \texttt{ROLE\_\allowbreak{}VETERINARIO} o
  \texttt{ROLE\_\allowbreak{}AUXILIAR} registre una nueva mascota (nombre, especie,
  raza y fecha de nacimiento) asociada a un usuario existente, activo y
  con rol \texttt{ROLE\_\allowbreak{}DUENO}.
\item
  \textbf{Rationale:} heredado de RF-01/RF-02 de la Entrega 1A (registro
  y asociación con propietario), acotado a la entidad Mascota
  implementada en esta fase.
\item
  \textbf{Criterios de aceptación:}

  \begin{enumerate}
  \def\labelenumi{\arabic{enumi}.}
  \tightlist
  \item
    Dado un usuario \texttt{ROLE\_\allowbreak{}ADMIN} y un dueño activo con rol
    \texttt{ROLE\_\allowbreak{}DUENO}, cuando envía \texttt{POST\ /\allowbreak{}api/\allowbreak{}mascotas} con
    datos válidos, entonces la respuesta es \texttt{201\ Created} con la
    mascota creada y \texttt{activo\ =\ true}.
  \item
    Dado un \texttt{duenio\allowbreak{}Id} que corresponde a un usuario cuyo rol
    \textbf{no} es \texttt{ROLE\_\allowbreak{}DUENO} (ADMIN, VETERINARIO o AUXILIAR),
    entonces la solicitud es rechazada con \texttt{400\ Bad\ Request} y
    \texttt{Problem\allowbreak{}Detail}.
  \item
    Dado un \texttt{duenio\allowbreak{}Id} inexistente o correspondiente a un usuario
    inactivo, entonces la respuesta es \texttt{404} con
    \texttt{Problem\allowbreak{}Detail}.
  \item
    Dado un usuario con rol \texttt{ROLE\_\allowbreak{}DUENO}, cuando intenta crear
    una mascota, entonces la respuesta es \texttt{403}.
  \item
    Dado un cuerpo con campos obligatorios ausentes o inválidos,
    entonces la respuesta es \texttt{422} con \texttt{Problem\allowbreak{}Detail} de
    validación.
  \end{enumerate}
\item
  \textbf{Método de verificación:} pruebas automatizadas
  \texttt{Mascota\allowbreak{}Controller\allowbreak{}Test.\allowbreak{}admin\allowbreak{}Crea\allowbreak{}Mascota\allowbreak{}Con\allowbreak{}Duenio\allowbreak{}Activo\allowbreak{}Rol\allowbreak{}Dueno\allowbreak{}Exitosa}
  (criterio 1);
  \texttt{admin\allowbreak{}Intenta\allowbreak{}Crear\allowbreak{}Mascota\allowbreak{}Asignada\allowbreak{}AUsuario\allowbreak{}Admin\allowbreak{}Es\allowbreak{}Rechazado},
  \texttt{admin\allowbreak{}Intenta\allowbreak{}Crear\allowbreak{}Mascota\allowbreak{}Asignada\allowbreak{}AVeterinario\allowbreak{}Es\allowbreak{}Rechazado},
  \texttt{admin\allowbreak{}Intenta\allowbreak{}Crear\allowbreak{}Mascota\allowbreak{}Asignada\allowbreak{}AAuxiliar\allowbreak{}Es\allowbreak{}Rechazado}
  (criterio 2);
  \texttt{admin\allowbreak{}Intenta\allowbreak{}Crear\allowbreak{}Mascota\allowbreak{}Con\allowbreak{}Duenio\allowbreak{}Id\allowbreak{}Inexistente\allowbreak{}Devuelve404},
  \texttt{admin\allowbreak{}Intenta\allowbreak{}Crear\allowbreak{}Mascota\allowbreak{}Con\allowbreak{}Usuario\allowbreak{}Inactivo\allowbreak{}Devuelve404}
  (criterio 3);
  \texttt{crear\allowbreak{}Mascota\allowbreak{}Con\allowbreak{}Rol\allowbreak{}Insuficiente\allowbreak{}Devuelve\allowbreak{}Problem\allowbreak{}Detail} (criterio
  4); \texttt{crear\allowbreak{}Mascota\allowbreak{}Con\allowbreak{}Campos\allowbreak{}Invalidos\allowbreak{}Devuelve422} (criterio 5).
\item
  \textbf{Trazabilidad:} HU-007 → CU-07 →
  \texttt{Mascota\allowbreak{}Controller.\allowbreak{}crear} → \texttt{Mascota\allowbreak{}Service.\allowbreak{}crear} →
  \texttt{Mascota\allowbreak{}Service.\allowbreak{}resolver\allowbreak{}Duenio} → \texttt{POST\ /\allowbreak{}api/\allowbreak{}mascotas};
  DTO \texttt{Mascota\allowbreak{}Request}/\texttt{Mascota\allowbreak{}Response}.
\item
  \textbf{Estado:} verificado
\end{itemize}

\subsubsection{REQ-F-009 --- Listado paginado de mascotas activas,
filtrado por propietario según
rol}\label{req-f-009-listado-paginado-de-mascotas-activas-filtrado-por-propietario-seguxfan-rol}

\begin{itemize}
\tightlist
\item
  \textbf{Tipo:} Funcional
\item
  \textbf{Prioridad:} Must
\item
  \textbf{Enunciado:} El sistema deberá permitir consultar el listado
  paginado de mascotas activas; si el solicitante tiene rol
  \texttt{ROLE\_\allowbreak{}DUENO}, el listado deberá restringirse únicamente a sus
  propias mascotas, y para los demás roles autorizados deberá incluir
  todas las mascotas activas.
\item
  \textbf{Rationale:} heredado de RF-02 de la Entrega 1A, con la regla
  de aislamiento de datos entre dueños explícita en
  \texttt{Mascota\allowbreak{}Service.\allowbreak{}listar}.
\item
  \textbf{Criterios de aceptación:}

  \begin{enumerate}
  \def\labelenumi{\arabic{enumi}.}
  \tightlist
  \item
    Dado un usuario \texttt{ROLE\_\allowbreak{}DUENO} con mascotas propias y
    existiendo mascotas de otro dueño, cuando invoca
    \texttt{GET\ /\allowbreak{}api/\allowbreak{}mascotas}, entonces la respuesta es
    \texttt{200\ OK} y contiene exclusivamente sus mascotas: ninguna
    mascota de otro dueño aparece en la página.
  \item
    Dado un usuario \texttt{ROLE\_\allowbreak{}ADMIN} o \texttt{ROLE\_\allowbreak{}AUXILIAR},
    cuando invoca el mismo endpoint, entonces el listado incluye las
    mascotas activas de todos los dueños.
  \item
    Dadas dos cuentas \texttt{ROLE\_\allowbreak{}DUENO} distintas que solicitan la
    misma página con los mismos parámetros de paginación, entonces cada
    una recibe únicamente sus propios datos: la caché no comparte
    resultados entre usuarios.
  \item
    Las mascotas con \texttt{activo\ =\ false} no aparecen en el
    listado.
  \end{enumerate}
\item
  \textbf{Método de verificación:} pruebas automatizadas
  \texttt{Mascota\allowbreak{}Controller\allowbreak{}Test.\allowbreak{}dueno\allowbreak{}Solo\allowbreak{}Ve\allowbreak{}Sus\allowbreak{}Propias\allowbreak{}Mascotas\allowbreak{}En\allowbreak{}Listado}
  (criterio 1), \texttt{admin\allowbreak{}Conserva\allowbreak{}Listado\allowbreak{}Global\allowbreak{}De\allowbreak{}Mascotas} y
  \texttt{auxiliar\allowbreak{}Conserva\allowbreak{}Acceso\allowbreak{}Global\allowbreak{}Al\allowbreak{}Listado} (criterio 2),
  \texttt{dos\allowbreak{}Duenos\allowbreak{}Con\allowbreak{}Misma\allowbreak{}Pagina\allowbreak{}No\allowbreak{}Comparen\allowbreak{}Resultados\allowbreak{}De\allowbreak{}Cache} (criterio
  3), \texttt{admin\allowbreak{}Elimina\allowbreak{}Mascota\allowbreak{}Exitosamente} (criterio 4, que
  comprueba que la mascota dada de baja desaparece del listado activo).
\item
  \textbf{Trazabilidad:} HU-008 → CU-08 →
  \texttt{Mascota\allowbreak{}Controller.\allowbreak{}listar} → \texttt{Mascota\allowbreak{}Service.\allowbreak{}listar}
  (\texttt{@Cacheable(value\ =\ "\allowbreak{}mascotas",\allowbreak{}\ key\ =\ "\#email\ +\ …")})
  → \texttt{Mascota\allowbreak{}Repository.\allowbreak{}find\allowbreak{}All\allowbreak{}By\allowbreak{}Duenio\allowbreak{}Id\allowbreak{}And\allowbreak{}Activo\allowbreak{}True} /
  \texttt{find\allowbreak{}All\allowbreak{}By\allowbreak{}Activo\allowbreak{}True} → \texttt{GET\ /\allowbreak{}api/\allowbreak{}mascotas}; evidencia
  \texttt{docs/\allowbreak{}mediciones/\allowbreak{}sec/\allowbreak{}A01-\allowbreak{}access-\allowbreak{}control.\allowbreak{}md}.
\item
  \textbf{Estado:} verificado
\end{itemize}

\subsubsection{REQ-F-010 --- Consulta de mascota por
identificador}\label{req-f-010-consulta-de-mascota-por-identificador}

\begin{itemize}
\tightlist
\item
  \textbf{Tipo:} Funcional
\item
  \textbf{Prioridad:} Must
\item
  \textbf{Enunciado:} El sistema deberá permitir consultar el detalle
  completo de una mascota activa a partir de su identificador,
  respondiendo \texttt{404} con \texttt{Problem\allowbreak{}Detail} si no existe y
  \texttt{403} si el solicitante es un \texttt{ROLE\_\allowbreak{}DUENO} que no es su
  propietario.
\item
  \textbf{Rationale:} complementa REQ-F-009 para el caso de consulta
  puntual, necesario antes de una edición o para mostrar el detalle en
  la interfaz.
\item
  \textbf{Criterios de aceptación:}

  \begin{enumerate}
  \def\labelenumi{\arabic{enumi}.}
  \tightlist
  \item
    Dado un \texttt{ROLE\_\allowbreak{}DUENO} y una mascota de su propiedad, cuando
    invoca \texttt{GET\ /\allowbreak{}api/\allowbreak{}mascotas/\allowbreak{}\{id\}}, entonces la respuesta es
    \texttt{200\ OK} con el detalle de esa mascota.
  \item
    Dado un \texttt{ROLE\_\allowbreak{}DUENO} y una mascota de otro dueño, entonces
    la respuesta es \texttt{403} con \texttt{Problem\allowbreak{}Detail} de tipo
    \texttt{urn:biopet:error:\allowbreak{}forbidden}.
  \item
    Dado un \texttt{ROLE\_\allowbreak{}ADMIN} o un \texttt{ROLE\_\allowbreak{}VETERINARIO},
    entonces puede consultar la mascota de cualquier dueño con
    \texttt{200\ OK}.
  \item
    Dado un identificador inexistente o de una mascota inactiva,
    entonces la respuesta es \texttt{404} con \texttt{Problem\allowbreak{}Detail} de
    tipo \texttt{urn:biopet:error:\allowbreak{}not-\allowbreak{}found}.
  \end{enumerate}
\item
  \textbf{Método de verificación:} pruebas automatizadas
  \texttt{Mascota\allowbreak{}Controller\allowbreak{}Test.\allowbreak{}dueno\allowbreak{}Consulta\allowbreak{}Su\allowbreak{}Propia\allowbreak{}Mascota\allowbreak{}Por\allowbreak{}Id}
  (criterio 1), \texttt{dueno\allowbreak{}Consulta\allowbreak{}Mascota\allowbreak{}De\allowbreak{}Otro\allowbreak{}Duenio\allowbreak{}Devuelve403}
  (criterio 2), \texttt{admin\allowbreak{}Consulta\allowbreak{}Mascota\allowbreak{}De\allowbreak{}Cualquier\allowbreak{}Duenio} y
  \texttt{veterinario\allowbreak{}Consulta\allowbreak{}Mascota\allowbreak{}De\allowbreak{}Cualquier\allowbreak{}Duenio} (criterio 3),
  \texttt{buscar\allowbreak{}Mascota\allowbreak{}Inexistente\allowbreak{}Devuelve\allowbreak{}Problem\allowbreak{}Detail} (criterio 4).
  Evidencia HTTP adicional del criterio 2 en
  \texttt{docs/\allowbreak{}mediciones/\allowbreak{}sec/\allowbreak{}raw/\allowbreak{}A01-\allowbreak{}access-\allowbreak{}control.\allowbreak{}txt}.
\item
  \textbf{Trazabilidad:} HU-009 → CU-09 →
  \texttt{Mascota\allowbreak{}Controller.\allowbreak{}buscar} → \texttt{Mascota\allowbreak{}Service.\allowbreak{}buscar} →
  \texttt{Mascota\allowbreak{}Service.\allowbreak{}verificar\allowbreak{}Propiedad} →
  \texttt{Recurso\allowbreak{}No\allowbreak{}Encontrado\allowbreak{}Exception} →
  \texttt{GET\ /\allowbreak{}api/\allowbreak{}mascotas/\allowbreak{}\{id\}}.
\item
  \textbf{Estado:} verificado
\end{itemize}

\subsubsection{REQ-F-011 --- Actualización de mascota con verificación
de
propiedad}\label{req-f-011-actualizaciuxf3n-de-mascota-con-verificaciuxf3n-de-propiedad}

\begin{itemize}
\tightlist
\item
  \textbf{Tipo:} Funcional
\item
  \textbf{Prioridad:} Must
\item
  \textbf{Enunciado:} El sistema deberá permitir que un usuario con rol
  \texttt{ROLE\_\allowbreak{}ADMIN}, \texttt{ROLE\_\allowbreak{}VETERINARIO} o
  \texttt{ROLE\_\allowbreak{}AUXILIAR} actualice los datos de una mascota activa,
  registrando automáticamente la fecha de actualización
  (\texttt{actualizado\_\allowbreak{}en}) y manteniendo la restricción de que el
  propietario asignado tenga rol \texttt{ROLE\_\allowbreak{}DUENO}.
\item
  \textbf{Rationale:} heredado de RF-02 de la Entrega 1A.
\item
  \textbf{Criterios de aceptación:}

  \begin{enumerate}
  \def\labelenumi{\arabic{enumi}.}
  \tightlist
  \item
    Dado un \texttt{ROLE\_\allowbreak{}ADMIN} y una mascota activa, cuando envía
    \texttt{PUT\ /\allowbreak{}api/\allowbreak{}mascotas/\allowbreak{}\{id\}} con datos válidos, entonces la
    respuesta es \texttt{200\ OK} con los datos actualizados.
  \item
    Dado un nuevo \texttt{duenio\allowbreak{}Id} cuyo rol no es \texttt{ROLE\_\allowbreak{}DUENO},
    entonces la solicitud es rechazada con \texttt{400} y
    \texttt{Problem\allowbreak{}Detail}.
  \item
    Dado un usuario con rol \texttt{ROLE\_\allowbreak{}DUENO}, cuando intenta
    actualizar incluso su propia mascota, entonces la respuesta es
    \texttt{403} (la restricción es por rol, no por propiedad).
  \item
    Tras una actualización, el valor de \texttt{actualizado\_\allowbreak{}en} de la
    fila es posterior al que tenía antes de la operación.
  \end{enumerate}
\item
  \textbf{Método de verificación:} pruebas automatizadas
  \texttt{Mascota\allowbreak{}Controller\allowbreak{}Test.\allowbreak{}admin\allowbreak{}Actualiza\allowbreak{}Mascota\allowbreak{}Asignandola\allowbreak{}AOtro\allowbreak{}Usuario\allowbreak{}Activo\allowbreak{}Dueno\allowbreak{}Exitosa}
  (criterio 1),
  \texttt{admin\allowbreak{}Intenta\allowbreak{}Actualizar\allowbreak{}Mascota\allowbreak{}Con\allowbreak{}Rol\allowbreak{}De\allowbreak{}Duenio\allowbreak{}Incorrecto\allowbreak{}Es\allowbreak{}Rechazado}
  (criterio 2),
  \texttt{dueno\allowbreak{}Intenta\allowbreak{}Actualizar\allowbreak{}Mascota\allowbreak{}Propia\allowbreak{}Sigue\allowbreak{}Recibiendo403\allowbreak{}Por\allowbreak{}Rol}
  (criterio 3), y prueba de integración
  \texttt{Trigger\allowbreak{}Actualizado\allowbreak{}En\allowbreak{}Integration\allowbreak{}Test} sobre el trigger
  \texttt{set\_\allowbreak{}actualizado\_\allowbreak{}en} en PostgreSQL real (criterio 4).
\item
  \textbf{Trazabilidad:} HU-010 → CU-10 →
  \texttt{Mascota\allowbreak{}Controller.\allowbreak{}actualizar} →
  \texttt{Mascota\allowbreak{}Service.\allowbreak{}actualizar} → trigger
  \texttt{set\_\allowbreak{}actualizado\_\allowbreak{}en}
  (\texttt{Backend/\allowbreak{}src/\allowbreak{}main/\allowbreak{}resources/\allowbreak{}db/\allowbreak{}migration/\allowbreak{}V1\_\allowbreak{}\_\allowbreak{}schema\_\allowbreak{}inicial.\allowbreak{}sql})
  → \texttt{PUT\ /\allowbreak{}api/\allowbreak{}mascotas/\allowbreak{}\{id\}}.
\item
  \textbf{Estado:} verificado
\end{itemize}

\subsubsection{REQ-F-012 --- Baja lógica de una
mascota}\label{req-f-012-baja-luxf3gica-de-una-mascota}

\begin{itemize}
\tightlist
\item
  \textbf{Tipo:} Funcional
\item
  \textbf{Prioridad:} Must
\item
  \textbf{Enunciado:} El sistema deberá permitir dar de baja lógicamente
  una mascota existente (cambio del atributo \texttt{activo} a
  \texttt{false}), sin eliminar físicamente el registro de la base de
  datos.
\item
  \textbf{Rationale:} preservación de historial e integridad referencial
  con consultas, citas y vacunas, que referencian la mascota.
\item
  \textbf{Criterios de aceptación:}

  \begin{enumerate}
  \def\labelenumi{\arabic{enumi}.}
  \tightlist
  \item
    Dado un \texttt{ROLE\_\allowbreak{}ADMIN} y una mascota activa, cuando envía
    \texttt{DELETE\ /\allowbreak{}api/\allowbreak{}mascotas/\allowbreak{}\{id\}}, entonces la respuesta es
    \texttt{204\ No\ Content}.
  \item
    Tras la baja, la fila sigue existiendo en la tabla \texttt{mascotas}
    con \texttt{activo\ =\ false}; no se ejecuta ningún \texttt{DELETE}
    físico.
  \item
    Tras la baja, la mascota deja de aparecer en
    \texttt{GET\ /\allowbreak{}api/\allowbreak{}mascotas} y \texttt{GET\ /\allowbreak{}api/\allowbreak{}mascotas/\allowbreak{}\{id\}}
    responde \texttt{404}.
  \item
    Dado un usuario con rol \texttt{ROLE\_\allowbreak{}DUENO}, cuando intenta dar de
    baja incluso su propia mascota, entonces la respuesta es
    \texttt{403}.
  \end{enumerate}
\item
  \textbf{Método de verificación:} prueba automatizada
  \texttt{Mascota\allowbreak{}Controller\allowbreak{}Test.\allowbreak{}admin\allowbreak{}Elimina\allowbreak{}Mascota\allowbreak{}Exitosamente}
  (criterios 1 a 3, incluida la comprobación de que la mascota
  desaparece del listado activo) y
  \texttt{dueno\allowbreak{}Intenta\allowbreak{}Eliminar\allowbreak{}Mascota\allowbreak{}Propia\allowbreak{}Sigue\allowbreak{}Recibiendo403\allowbreak{}Por\allowbreak{}Rol}
  (criterio 4). Criterio 2: inspección de
  \texttt{Mascota\allowbreak{}Service.\allowbreak{}eliminar}
  (\texttt{mascota.\allowbreak{}set\allowbreak{}Activo(false);\allowbreak{}\ mascota\allowbreak{}Repository.\allowbreak{}save(mascota);},
  sin invocación a \texttt{delete}).
\item
  \textbf{Trazabilidad:} HU-011 → CU-11 →
  \texttt{Mascota\allowbreak{}Controller.\allowbreak{}eliminar} → \texttt{Mascota\allowbreak{}Service.\allowbreak{}eliminar}
  → \texttt{DELETE\ /\allowbreak{}api/\allowbreak{}mascotas/\allowbreak{}\{id\}}; sección 2.9.2 (ciclo de vida
  común).
\item
  \textbf{Estado:} verificado
\end{itemize}

\subsubsection{REQ-F-013 --- Registro de atención médica y consulta del
historial
clínico}\label{req-f-013-registro-de-atenciuxf3n-muxe9dica-y-consulta-del-historial-cluxednico}

\begin{quote}
\textbf{Requisito con dos obligaciones.} Este requisito, heredado de
RF-03 y RF-04 de la Entrega 1A, agrupa dos obligaciones distintas con
distinto grado de avance: \textbf{(A)} registrar y gestionar las
consultas médicas, y \textbf{(B)} consultar el historial clínico
consolidado de una mascota. Conserva su identificador único, su historia
HU-012 y su caso de uso CU-12; las obligaciones se separan más abajo con
criterios numerados \texttt{A1…A7} y \texttt{B1…B5}, y el Estado
describe el requisito completo.
\end{quote}

\begin{itemize}
\tightlist
\item
  \textbf{Tipo:} Funcional
\item
  \textbf{Prioridad:} Should
\item
  \textbf{Enunciado:} El sistema deberá permitir \textbf{(A)} registrar,
  consultar por identificador, actualizar y dar de baja lógica las
  consultas médicas (motivo, diagnóstico, tratamiento y observaciones)
  asociadas a una mascota activa y a un usuario con rol
  \texttt{ROLE\_\allowbreak{}VETERINARIO}, restringiendo el acceso de un usuario con
  rol \texttt{ROLE\_\allowbreak{}DUENO} a las consultas de sus propias mascotas; y
  \textbf{(B)} consultar, a través de un endpoint propio, el historial
  clínico consolidado de una mascota activa ---datos de la mascota y de
  su dueño, número de consultas y de citas, última consulta, última
  vacuna y próxima vacuna---, ordenando cronológicamente los eventos
  clínicos que lo componen.
\item
  \textbf{Rationale:} heredado de RF-03 (registro de atención médica) y
  RF-04 (consulta del historial clínico ``de forma cronológica'') de la
  Entrega 1A. La Unidad IV entregó el módulo real
  \texttt{Consulta\allowbreak{}Controller}/\texttt{Consulta\allowbreak{}Service}/\texttt{Consulta\allowbreak{}Repository},
  que cubre la obligación A, y la capa de datos de la obligación B
  (\texttt{fn\_\allowbreak{}historial\_\allowbreak{}clini\allowbreak{}co\_\allowbreak{}mascota}), por lo que el requisito
  deja de estar sin implementación.
\item
  \textbf{Criterios de aceptación:}

  \begin{enumerate}
  \def\labelenumi{\arabic{enumi}.}
  \tightlist
  \item
    \textbf{(A1)} Dado un \texttt{ROLE\_\allowbreak{}ADMIN},
    \texttt{ROLE\_\allowbreak{}VETERINARIO} o \texttt{ROLE\_\allowbreak{}AUXILIAR}, una mascota
    activa y un usuario con rol \texttt{ROLE\_\allowbreak{}VETERINARIO}, cuando envía
    \texttt{POST\ /\allowbreak{}api/\allowbreak{}consultas} con datos válidos, entonces la
    respuesta es \texttt{201\ Created} con la consulta creada y
    \texttt{activo\ =\ true}.
  \item
    \textbf{(A2)} Dado un \texttt{veterinario\allowbreak{}Id} cuyo rol no es
    \texttt{ROLE\_\allowbreak{}VETERINARIO}, entonces la solicitud es rechazada con
    \texttt{400} y \texttt{Problem\allowbreak{}Detail}.
  \item
    \textbf{(A3)} Dado un \texttt{ROLE\_\allowbreak{}DUENO}, cuando intenta crear,
    actualizar o eliminar una consulta, entonces la respuesta es
    \texttt{403}.
  \item
    \textbf{(A4)} Dado un \texttt{ROLE\_\allowbreak{}DUENO} y una consulta de una
    mascota ajena, cuando invoca \texttt{GET\ /\allowbreak{}api/\allowbreak{}consultas/\allowbreak{}\{id\}},
    entonces la respuesta es \texttt{403}.
  \item
    \textbf{(A5)} Dado un identificador inexistente o de una consulta
    inactiva, entonces \texttt{GET\ /\allowbreak{}api/\allowbreak{}consultas/\allowbreak{}\{id\}} responde
    \texttt{404} con \texttt{Problem\allowbreak{}Detail}.
  \item
    \textbf{(A6)} Dado un \texttt{ROLE\_\allowbreak{}DUENO}, cuando invoca
    \texttt{GET\ /\allowbreak{}api/\allowbreak{}consultas}, entonces \textbf{ningún} elemento de
    la página devuelta corresponde a una mascota de otro dueño.
  \item
    \textbf{(A7)} Dado un \texttt{DELETE\ /\allowbreak{}api/\allowbreak{}consultas/\allowbreak{}\{id\}},
    entonces la fila permanece en la tabla con \texttt{activo\ =\ false}
    y deja de aparecer en listados y búsquedas.
  \item
    \textbf{(B1)} Dada una mascota activa con consultas, citas y vacunas
    registradas, cuando se solicita su historial clínico consolidado,
    entonces la respuesta incluye los datos de la mascota, los de su
    dueño, el conteo de consultas y de citas, la última consulta, la
    última vacuna y la próxima vacuna.
  \item
    \textbf{(B2)} Dada una mascota sin eventos clínicos, entonces la
    respuesta es satisfactoria con los conteos en cero, no un error.
  \item
    \textbf{(B3)} Dado un identificador de mascota inexistente, entonces
    el resultado es vacío (a nivel de rutina) o \texttt{404} con
    \texttt{Problem\allowbreak{}Detail} (a nivel de API).
  \item
    \textbf{(B4)} Dado un usuario con rol \texttt{ROLE\_\allowbreak{}DUENO}, cuando
    solicita el historial de una mascota ajena, entonces la respuesta es
    \texttt{403}.
  \item
    \textbf{(B5)} La agregación se implementa como procedimiento
    almacenado versionado en \texttt{db/\allowbreak{}procs/}, no como JPQL con joins
    (ver REQ-NF-013).
  \end{enumerate}
\item
  \textbf{Método de verificación:} obligación A --- pruebas
  automatizadas
  \texttt{Consulta\allowbreak{}Controller\allowbreak{}Test.\allowbreak{}admin\allowbreak{}Crea\allowbreak{}Consulta\allowbreak{}Exitosamente} (A1),
  \texttt{crear\allowbreak{}Consulta\allowbreak{}Con\allowbreak{}Campos\allowbreak{}Invalidos\allowbreak{}Devuelve422} e inspección de
  \texttt{Consulta\allowbreak{}Service.\allowbreak{}resolver\allowbreak{}Veterinario} (A2),
  \texttt{dueno\allowbreak{}No\allowbreak{}Puede\allowbreak{}Crear\allowbreak{}Consulta\allowbreak{}Devuelve403} (A3),
  \texttt{dueno\allowbreak{}De\allowbreak{}Otra\allowbreak{}Mascota\allowbreak{}No\allowbreak{}Puede\allowbreak{}Ver\allowbreak{}Consulta\allowbreak{}Ajena} (A4),
  \texttt{buscar\allowbreak{}Consulta\allowbreak{}Inexistente\allowbreak{}Devuelve404} (A5),
  \texttt{admin\allowbreak{}Elimina\allowbreak{}Consulta\allowbreak{}Exitosamente} (A7); \textbf{A6 sin
  artefacto que lo demuestre y, además, incumplido --- ver
  Observaciones}. Obligación B --- pruebas de integración
  \texttt{Procedimientos\allowbreak{}Biopet\allowbreak{}Integration\allowbreak{}Test.\allowbreak{}historial\allowbreak{}Clinico\_\allowbreak{}mascota\allowbreak{}Con\allowbreak{}Datos\_\allowbreak{}consolida\allowbreak{}Historial}
  y \texttt{historial\allowbreak{}Clinico\_\allowbreak{}mascota\allowbreak{}Inexistente\_\allowbreak{}devue\allowbreak{}lve\allowbreak{}Lista\allowbreak{}Vacia}
  sobre PostgreSQL real con Testcontainers (B1 a B3 a nivel de rutina) y
  \texttt{docs/\allowbreak{}basedatos/\allowbreak{}CATALOGO-\allowbreak{}SP.\allowbreak{}md} (B5); \textbf{B4 y la
  exposición de B1 como endpoint: verificación prevista} --- prueba de
  controlador equivalente a
  \texttt{Mascota\allowbreak{}Controller\allowbreak{}Test.\allowbreak{}dueno\allowbreak{}Consulta\allowbreak{}Mascota\allowbreak{}De\allowbreak{}Otro\allowbreak{}Duenio\allowbreak{}Devuelve403},
  a escribir cuando el endpoint exista.
\item
  \textbf{Trazabilidad:} HU-012 → CU-12 → obligación A:
  \texttt{Consulta\allowbreak{}Controller.\allowbreak{}\{listar,\allowbreak{}buscar,\allowbreak{}crear,\allowbreak{}actualizar,\allowbreak{}eliminar\}}
  → \texttt{Consulta\allowbreak{}Service} → \texttt{Consulta\allowbreak{}Repository} →
  \texttt{GET\ /\allowbreak{}api/\allowbreak{}consultas}, \texttt{GET\ /\allowbreak{}api/\allowbreak{}consultas/\allowbreak{}\{id\}},
  \texttt{POST\ /\allowbreak{}api/\allowbreak{}consultas}, \texttt{PUT\ /\allowbreak{}api/\allowbreak{}consultas/\allowbreak{}\{id\}},
  \texttt{DELETE\ /\allowbreak{}api/\allowbreak{}consultas/\allowbreak{}\{id\}}; prueba
  \texttt{Backend/\allowbreak{}src/\allowbreak{}test/\allowbreak{}java/\allowbreak{}com/\allowbreak{}biopet/\allowbreak{}Consulta\allowbreak{}Controller\allowbreak{}Test.\allowbreak{}java};
  procedimiento de validación cruzada
  \texttt{db/\allowbreak{}procs/\allowbreak{}sp\_\allowbreak{}registrar\_\allowbreak{}consu\allowbreak{}lta\_\allowbreak{}validada.\allowbreak{}sql} con su prueba
  \texttt{Procedimientos\allowbreak{}Biopet\allowbreak{}Integration\allowbreak{}Test.\allowbreak{}registrar\allowbreak{}Consulta\allowbreak{}Validada\_\allowbreak{}caso\allowbreak{}Feliz\_\allowbreak{}retorna\allowbreak{}Id\allowbreak{}YPersiste}.
  Obligación B: \texttt{db/\allowbreak{}procs/\allowbreak{}fn\_\allowbreak{}historial\_\allowbreak{}clini\allowbreak{}co\_\allowbreak{}mascota.\allowbreak{}sql} →
  \texttt{Procedimiento\allowbreak{}Biopet\allowbreak{}Repository.\allowbreak{}historial\allowbreak{}Clinico\allowbreak{}Mascota} →
  proyección \texttt{com.\allowbreak{}biopet.\allowbreak{}repository.\allowbreak{}Historial\allowbreak{}Clinico} →
  \texttt{@Named\allowbreak{}Stored\allowbreak{}Procedure\allowbreak{}Query} declarado en
  \texttt{com.\allowbreak{}biopet.\allowbreak{}entity.\allowbreak{}Mascota} →
  \texttt{Backend/\allowbreak{}src/\allowbreak{}test/\allowbreak{}java/\allowbreak{}com/\allowbreak{}biopet/\allowbreak{}repository/\allowbreak{}Procedimientos\allowbreak{}Biopet\allowbreak{}Integration\allowbreak{}Test.\allowbreak{}java};
  migración
  \texttt{Backend/\allowbreak{}src/\allowbreak{}main/\allowbreak{}resources/\allowbreak{}db/\allowbreak{}migration/\allowbreak{}V6\_\allowbreak{}\_\allowbreak{}formalizar\_\allowbreak{}proc\allowbreak{}edimientos\_\allowbreak{}jpa.\allowbreak{}sql};
  catálogo \texttt{docs/\allowbreak{}basedatos/\allowbreak{}CATALOGO-\allowbreak{}SP.\allowbreak{}md}; matriz de permisos,
  sección 2.8.5.
\item
  \textbf{Estado:} implementado
\item
  \textbf{Observaciones:} el requisito \textbf{no está cumplido}. Tres
  limitaciones reales, verificadas contra el código actual y \textbf{no
  corregidas} en esta revisión documental:

  \begin{enumerate}
  \def\labelenumi{(\alph{enumi})}
  \tightlist
  \item
    \textbf{El criterio A6 no se cumple.}
    \texttt{Consulta\allowbreak{}Service.\allowbreak{}listar()} devuelve
    \texttt{consulta\allowbreak{}Repository.\allowbreak{}find\allowbreak{}All\allowbreak{}By\allowbreak{}Activo\allowbreak{}True(pageable)} tanto para
    \texttt{ROLE\_\allowbreak{}DUENO} como para el resto de roles, y
    \texttt{Consulta\allowbreak{}Repository} no declara ningún método de filtrado por
    propietario (a diferencia de \texttt{Cita\allowbreak{}Repository} y
    \texttt{Vacuna\allowbreak{}Repository}, que sí declaran
    \texttt{find\allowbreak{}All\allowbreak{}By\allowbreak{}Mascota\_\allowbreak{}Duenio\_\allowbreak{}Id\allowbreak{}And\allowbreak{}Activo\allowbreak{}True}). Ver REQ-F-006 y
    REQ-NF-015.
  \item
    \textbf{La obligación B no está expuesta.} Ningún controlador de
    \texttt{com.\allowbreak{}biopet.\allowbreak{}controller} invoca
    \texttt{Procedimiento\allowbreak{}Biopet\allowbreak{}Repository.\allowbreak{}historial\allowbreak{}Clinico\allowbreak{}Mascota}: no
    hay endpoint de historial consolidado ni pantalla asociada, por lo
    que B4 no puede comprobarse. Además la rutina devuelve un
    \textbf{resumen consolidado} (conteos, última consulta, última y
    próxima vacuna), no la lista cronológica completa de eventos que
    pide el enunciado original.
  \item
    \texttt{GET\ /\allowbreak{}api/\allowbreak{}consultas} (listado) y
    \texttt{PUT\ /\allowbreak{}api/\allowbreak{}consultas/\allowbreak{}\{id\}} no tienen prueba automatizada
    dedicada en \texttt{Consulta\allowbreak{}Controller\allowbreak{}Test} (6 pruebas frente a las
    25 de \texttt{Cita\allowbreak{}Controller\allowbreak{}Test}, que cubre el módulo equivalente).
  \end{enumerate}
\end{itemize}

\subsubsection{REQ-F-014 --- Registro de medicamentos
prescritos}\label{req-f-014-registro-de-medicamentos-prescritos}

\begin{itemize}
\tightlist
\item
  \textbf{Tipo:} Funcional
\item
  \textbf{Prioridad:} Could
\item
  \textbf{Enunciado:} El sistema deberá permitir registrar los
  medicamentos prescritos durante una consulta médica (nombre del
  medicamento, dosis, frecuencia y duración), asociándolos a la consulta
  que los origina.
\item
  \textbf{Rationale:} heredado de RF-05 de la Entrega 1A. Depende de la
  obligación A de REQ-F-013 (la consulta médica debe existir antes de
  poder prescribir sobre ella).
\item
  \textbf{Criterios de aceptación:}

  \begin{enumerate}
  \def\labelenumi{\arabic{enumi}.}
  \tightlist
  \item
    Dada una consulta médica existente y activa, cuando se registra una
    prescripción con medicamento, dosis, frecuencia y duración, entonces
    la prescripción queda persistida y asociada a esa consulta.
  \item
    Dada una consulta inexistente, entonces la operación responde
    \texttt{404} con \texttt{Problem\allowbreak{}Detail}.
  \item
    Dado un usuario con rol \texttt{ROLE\_\allowbreak{}DUENO}, cuando intenta
    registrar una prescripción, entonces la respuesta es \texttt{403}.
  \item
    Dado un \texttt{ROLE\_\allowbreak{}DUENO} propietario de la mascota, cuando
    consulta las prescripciones de una consulta de esa mascota, entonces
    puede leerlas; si la mascota es ajena, la respuesta es \texttt{403}.
  \item
    Dado un cuerpo con campos obligatorios ausentes, entonces la
    respuesta es \texttt{422} con \texttt{Problem\allowbreak{}Detail} de validación.
  \end{enumerate}
\item
  \textbf{Método de verificación:} previsto --- prueba de controlador
  con \texttt{Mock\allowbreak{}Mvc} siguiendo el patrón de
  \texttt{Consulta\allowbreak{}Controller\allowbreak{}Test} (casos 201, 404, 403 y 422), más la
  entidad y migración Flyway correspondientes. No se puede ejecutar hoy
  porque no existe implementación.
\item
  \textbf{Trazabilidad:} HU-013 → CU-13 → depende de REQ-F-013; entidad
  conceptual del DER de la Entrega 1A (\texttt{PFC\_\allowbreak{}Entrega1\allowbreak{}A\_\allowbreak{}BMT.\allowbreak{}pdf},
  sección 6). Sin controlador, servicio, repositorio, entidad ni
  migración en el repositorio actual.
\item
  \textbf{Estado:} pendiente
\end{itemize}

\subsubsection{REQ-F-015 --- Gestión de citas
veterinarias}\label{req-f-015-gestiuxf3n-de-citas-veterinarias}

\begin{quote}
\textbf{Requisito con dos obligaciones.} Heredado de RF-06 de la Entrega
1A (``Registrar, modificar y consultar citas veterinarias mediante
calendario interactivo''), agrupa \textbf{(A)} la gestión de las citas a
través de la API y \textbf{(B)} su presentación en un calendario
interactivo. Conserva su identificador único, HU-014 y CU-14; las
obligaciones se separan con criterios \texttt{A1…A6} y \texttt{B1…B5}, y
el Estado describe el requisito completo.
\end{quote}

\begin{itemize}
\tightlist
\item
  \textbf{Tipo:} Funcional
\item
  \textbf{Prioridad:} Should
\item
  \textbf{Enunciado:} El sistema deberá permitir \textbf{(A)} listar,
  consultar por identificador, crear, actualizar y dar de baja lógica
  citas veterinarias (mascota, veterinario asignado, fecha y hora,
  estado y motivo), restringiendo la lectura de un usuario con rol
  \texttt{ROLE\_\allowbreak{}DUENO} a las citas de sus propias mascotas y permitiendo
  a un \texttt{ROLE\_\allowbreak{}VETERINARIO} modificar únicamente las citas en las
  que él es el veterinario asignado; y \textbf{(B)} presentar esas citas
  en una vista de calendario que permita consultarlas por periodo, crear
  una cita sobre una fecha seleccionada y abrir el detalle de una cita
  existente, respetando la visibilidad por rol de la sección 2.8.4.
\item
  \textbf{Rationale:} heredado de RF-06 de la Entrega 1A, que pide
  explícitamente ``mediante calendario interactivo''. La Unidad IV
  entregó el módulo real
  \texttt{Cita\allowbreak{}Controller}/\texttt{Cita\allowbreak{}Service}/\texttt{Cita\allowbreak{}Repository}
  con CRUD completo y reglas de acceso por dato, que cubre la obligación
  A; la obligación B sigue sin iniciarse.
\item
  \textbf{Criterios de aceptación:}

  \begin{enumerate}
  \def\labelenumi{\arabic{enumi}.}
  \tightlist
  \item
    \textbf{(A1)} Dado un \texttt{ROLE\_\allowbreak{}ADMIN} o
    \texttt{ROLE\_\allowbreak{}AUXILIAR}, una mascota activa y un usuario con rol
    \texttt{ROLE\_\allowbreak{}VETERINARIO}, cuando envía \texttt{POST\ /\allowbreak{}api/\allowbreak{}citas}
    con datos válidos, entonces la respuesta es \texttt{201\ Created} y
    la cita queda con \texttt{estado\ =\ PROGRAMAD\allowbreak{}A}, aunque el cliente
    haya enviado otro estado.
  \item
    \textbf{(A2)} Dado un \texttt{ROLE\_\allowbreak{}DUENO} con citas propias y
    existiendo citas de otro dueño, cuando invoca
    \texttt{GET\ /\allowbreak{}api/\allowbreak{}citas}, entonces ningún elemento devuelto
    corresponde a otro dueño; y \texttt{GET\ /\allowbreak{}api/\allowbreak{}citas/\allowbreak{}\{id\}} de una
    cita ajena responde \texttt{403}.
  \item
    \textbf{(A3)} Dado un \texttt{ROLE\_\allowbreak{}VETERINARIO}, cuando intenta
    actualizar una cita asignada a otro veterinario, entonces la
    respuesta es \texttt{403}; cuando actualiza una cita propia, la
    respuesta es \texttt{200\ OK}.
  \item
    \textbf{(A4)} Dado un \texttt{ROLE\_\allowbreak{}DUENO} o un
    \texttt{ROLE\_\allowbreak{}VETERINARIO}, cuando intenta
    \texttt{POST\ /\allowbreak{}api/\allowbreak{}citas}, entonces la respuesta es \texttt{403};
    sólo \texttt{ROLE\_\allowbreak{}ADMIN} puede ejecutar
    \texttt{DELETE\ /\allowbreak{}api/\allowbreak{}citas/\allowbreak{}\{id\}}.
  \item
    \textbf{(A5)} Dada una mascota o un veterinario inexistentes,
    entonces la respuesta es \texttt{404}; dado un
    \texttt{veterinario\allowbreak{}Id} cuyo rol no es \texttt{ROLE\_\allowbreak{}VETERINARIO},
    entonces \texttt{400}; dados campos inválidos, \texttt{422}.
  \item
    \textbf{(A6)} Dado un \texttt{DELETE\ /\allowbreak{}api/\allowbreak{}citas/\allowbreak{}\{id\}}, entonces
    la fila permanece con \texttt{activo\ =\ false} (baja lógica) y el
    valor de \texttt{estado} no cambia.
  \item
    \textbf{(B1)} Dado un usuario autenticado, cuando abre la vista de
    citas, entonces se muestra una rejilla de calendario con las citas
    del periodo visible.
  \item
    \textbf{(B2)} Dado un \texttt{ROLE\_\allowbreak{}DUENO}, entonces el calendario
    muestra únicamente citas de sus propias mascotas.
  \item
    \textbf{(B3)} Dado un \texttt{ROLE\_\allowbreak{}ADMIN} o
    \texttt{ROLE\_\allowbreak{}AUXILIAR}, cuando selecciona una fecha libre, entonces
    puede crear una cita para esa fecha desde el calendario.
  \item
    \textbf{(B4)} Dada una cita mostrada en el calendario, cuando se
    selecciona, entonces se abre su detalle con mascota, veterinario,
    fecha, estado y motivo.
  \item
    \textbf{(B5)} La vista cumple el mismo umbral de accesibilidad
    exigido por REQ-NF-018.
  \end{enumerate}
\item
  \textbf{Método de verificación:} obligación A --- pruebas
  automatizadas de \texttt{Cita\allowbreak{}Controller\allowbreak{}Test} (25 pruebas):
  \texttt{crear\allowbreak{}Cita\allowbreak{}Valida\allowbreak{}Devuelve201},
  \texttt{crear\allowbreak{}Cita\allowbreak{}Forzando\allowbreak{}Estado\allowbreak{}Distinto\allowbreak{}Ignora\allowbreak{}El\allowbreak{}Estado\allowbreak{}Enviado},
  \texttt{auxiliar\allowbreak{}Puede\allowbreak{}Crear\allowbreak{}Cita} (A1);
  \texttt{dueno\allowbreak{}Solo\allowbreak{}Ve\allowbreak{}Sus\allowbreak{}Propias\allowbreak{}Citas\allowbreak{}En\allowbreak{}Listado},
  \texttt{dueno\allowbreak{}Consulta\allowbreak{}Cita\allowbreak{}De\allowbreak{}Su\allowbreak{}Propia\allowbreak{}Mascota},
  \texttt{dueno\allowbreak{}No\allowbreak{}Puede\allowbreak{}Consultar\allowbreak{}Cita\allowbreak{}De\allowbreak{}Mascota\allowbreak{}Ajena} (A2);
  \texttt{veterinario\allowbreak{}Actualiza\allowbreak{}Su\allowbreak{}Propia\allowbreak{}Cita\allowbreak{}Devuelve200},
  \texttt{veterinario\allowbreak{}No\allowbreak{}Puede\allowbreak{}Actualizar\allowbreak{}Cita\allowbreak{}De\allowbreak{}Otro\allowbreak{}Veterinario\allowbreak{}Devuelve403}
  (A3); \texttt{dueno\allowbreak{}No\allowbreak{}Puede\allowbreak{}Crear\allowbreak{}Cita\allowbreak{}Devuelve403},
  \texttt{veterinario\allowbreak{}No\allowbreak{}Puede\allowbreak{}Crear\allowbreak{}Cita\allowbreak{}Devuelve403},
  \texttt{auxiliar\allowbreak{}No\allowbreak{}Puede\allowbreak{}Eliminar\allowbreak{}Cita\allowbreak{}Devuelve403},
  \texttt{dueno\allowbreak{}No\allowbreak{}Puede\allowbreak{}Actualizar\allowbreak{}Cita\allowbreak{}Devuelve403} (A4);
  \texttt{crear\allowbreak{}Cita\allowbreak{}Con\allowbreak{}Mascota\allowbreak{}Inexistente\allowbreak{}Devuelve404},
  \texttt{crear\allowbreak{}Cita\allowbreak{}Con\allowbreak{}Veterinario\allowbreak{}Inexistente\allowbreak{}Devuelve404},
  \texttt{crear\allowbreak{}Cita\allowbreak{}Con\allowbreak{}Usuario\allowbreak{}No\allowbreak{}Veterinario\allowbreak{}Asignado\allowbreak{}Devuelve400},
  \texttt{crear\allowbreak{}Cita\allowbreak{}Con\allowbreak{}Campos\allowbreak{}Invalidos\allowbreak{}Devuelve422} (A5);
  \texttt{admin\allowbreak{}Elimina\allowbreak{}Cita\allowbreak{}Devuelve204\allowbreak{}YBaja\allowbreak{}Logica} (A6). Obligación B ---
  \textbf{verificación prevista}, no ejecutable hoy: componente Angular
  en \texttt{frontend/\allowbreak{}src/\allowbreak{}app/\allowbreak{}features/} consumiendo
  \texttt{GET\ /\allowbreak{}api/\allowbreak{}citas}, con corrida Lighthouse sobre la nueva ruta
  (B5) e inspección manual por rol (B1 a B4).
\item
  \textbf{Trazabilidad:} HU-014 → CU-14 → obligación A:
  \texttt{Cita\allowbreak{}Controller.\allowbreak{}\{listar,\allowbreak{}buscar,\allowbreak{}crear,\allowbreak{}actualizar,\allowbreak{}eliminar\}} →
  \texttt{Cita\allowbreak{}Service} (\texttt{verificar\allowbreak{}Acceso\allowbreak{}Lectura},
  \texttt{verificar\allowbreak{}Permiso\allowbreak{}Escritura}) →
  \texttt{Cita\allowbreak{}Repository.\allowbreak{}find\allowbreak{}All\allowbreak{}By\allowbreak{}Mascota\_\allowbreak{}Duenio\_\allowbreak{}Id\allowbreak{}And\allowbreak{}Activo\allowbreak{}True} →
  \texttt{GET\ /\allowbreak{}api/\allowbreak{}citas}, \texttt{GET\ /\allowbreak{}api/\allowbreak{}citas/\allowbreak{}\{id\}},
  \texttt{POST\ /\allowbreak{}api/\allowbreak{}citas}, \texttt{PUT\ /\allowbreak{}api/\allowbreak{}citas/\allowbreak{}\{id\}},
  \texttt{DELETE\ /\allowbreak{}api/\allowbreak{}citas/\allowbreak{}\{id\}}; procedimiento almacenado de
  actualización masiva
  \texttt{db/\allowbreak{}procs/\allowbreak{}sp\_\allowbreak{}actualizar\_\allowbreak{}esta\allowbreak{}do\_\allowbreak{}citas\_\allowbreak{}masivas.\allowbreak{}sql} con sus
  pruebas
  \texttt{Procedimientos\allowbreak{}Biopet\allowbreak{}Integration\allowbreak{}Test.\allowbreak{}actualizar\allowbreak{}Estado\allowbreak{}Citas\allowbreak{}Masivas\_\allowbreak{}solo\allowbreak{}Actualiza\allowbreak{}Las\allowbreak{}Que\allowbreak{}Cumplen\allowbreak{}Filtro}
  y
  \texttt{actualizar\allowbreak{}Estado\allowbreak{}Citas\allowbreak{}Masivas\_\allowbreak{}estado\allowbreak{}Invalido\_\allowbreak{}lanza\allowbreak{}Excepcion};
  estados y transiciones en la sección 2.9.1; matriz de permisos,
  sección 2.8.4. Obligación B: sin componente en
  \texttt{frontend/\allowbreak{}src/\allowbreak{}app/\allowbreak{}features/} (sólo existen
  \texttt{login.\allowbreak{}component.\allowbreak{}ts}, \texttt{mascotas.\allowbreak{}component.\allowbreak{}ts} y
  \texttt{vacunas.\allowbreak{}component.\allowbreak{}ts}); wireframe conceptual de la Entrega 1A
  (\texttt{PFC\_\allowbreak{}Entrega1\allowbreak{}A\_\allowbreak{}BMT.\allowbreak{}pdf}, sección 7).
\item
  \textbf{Estado:} implementado
\item
  \textbf{Observaciones:} el requisito \textbf{no está cumplido}. La
  obligación A está íntegramente realizada y evidenciada: los seis
  criterios \texttt{A1} a \texttt{A6} están cubiertos por
  \texttt{Cita\allowbreak{}Controller\allowbreak{}Test}, que es el módulo con mayor cobertura de
  pruebas del proyecto. La obligación B \textbf{no está iniciada}: no
  existe ningún componente de calendario en
  \texttt{frontend/\allowbreak{}src/\allowbreak{}app/\allowbreak{}features/}, por lo que los criterios
  \texttt{B1} a \texttt{B5} no tienen ni implementación ni artefacto. El
  requisito no puede declararse \texttt{verificado} mientras el
  calendario interactivo que exige RF-06 siga sin construirse.
\end{itemize}

\subsubsection{REQ-F-016 --- API de recepción de telemetría de
dispositivos
IoT}\label{req-f-016-api-de-recepciuxf3n-de-telemetruxeda-de-dispositivos-iot}

\begin{itemize}
\tightlist
\item
  \textbf{Tipo:} Funcional
\item
  \textbf{Prioridad:} Could
\item
  \textbf{Enunciado:} El sistema deberá exponer una API autenticada que
  reciba y almacene lecturas de telemetría (identificador de
  dispositivo, marca de tiempo y valores medidos) enviadas por
  dispositivos IoT de rastreo asociados a una mascota registrada.
\item
  \textbf{Rationale:} heredado de RF-08 y RF-09 de la Entrega 1A.
  Requiere una decisión de arquitectura previa sobre el protocolo de
  ingesta y el mecanismo de autenticación de dispositivos, que no se ha
  tomado.
\item
  \textbf{Criterios de aceptación:}

  \begin{enumerate}
  \def\labelenumi{\arabic{enumi}.}
  \tightlist
  \item
    Dado un dispositivo autenticado y asociado a una mascota activa,
    cuando envía una lectura válida, entonces la respuesta es
    \texttt{201\ Created} y la lectura queda persistida con su marca de
    tiempo.
  \item
    Dada una lectura de un dispositivo no registrado o no asociado a
    ninguna mascota, entonces la respuesta es \texttt{403} o
    \texttt{404} con \texttt{Problem\allowbreak{}Detail}, y no se persiste nada.
  \item
    Dado un cuerpo con valores fuera del rango declarado para el tipo de
    medida, entonces la respuesta es \texttt{422} con
    \texttt{Problem\allowbreak{}Detail}.
  \item
    Dado un \texttt{ROLE\_\allowbreak{}DUENO}, cuando consulta la telemetría de sus
    propias mascotas, entonces la obtiene; la de mascotas ajenas
    responde \texttt{403}.
  \end{enumerate}
\item
  \textbf{Método de verificación:} previsto --- prueba de integración
  del endpoint de ingesta con un cliente simulado de dispositivo, más un
  ADR que fije protocolo y autenticación de dispositivos antes de
  implementar. No ejecutable hoy.
\item
  \textbf{Trazabilidad:} HU-015 → CU-15 → entidad conceptual
  \texttt{Dispositivo\_\allowbreak{}Io\allowbreak{}T} del DER de la Entrega 1A; sistema externo
  identificado en \texttt{docs/\allowbreak{}diagrams/\allowbreak{}c4-\allowbreak{}contexto/\allowbreak{}C4-\allowbreak{}L1-\allowbreak{}contexto.\allowbreak{}md}.
  Sin controlador, servicio, entidad ni migración en el repositorio
  actual (ver sección 2.7.3).
\item
  \textbf{Estado:} pendiente
\end{itemize}

\subsubsection{REQ-F-017 --- Generación de recomendaciones clínicas
informativas a partir del historial
médico}\label{req-f-017-generaciuxf3n-de-recomendaciones-cluxednicas-informativas-a-partir-del-historial-muxe9dico}

\begin{itemize}
\tightlist
\item
  \textbf{Tipo:} Funcional
\item
  \textbf{Prioridad:} Could
\item
  \textbf{Enunciado:} Al recibir una solicitud de recomendaciones para
  una mascota con historial clínico registrado, el sistema deberá
  generar, mediante un servicio de IA externo, una lista de
  recomendaciones de cuidado en texto a partir de los datos del
  historial clínico de esa mascota, y deberá devolver cada recomendación
  acompañada de la advertencia explícita ``informativa, no sustituye
  diagnóstico veterinario'', enviando al servicio externo exclusivamente
  los campos clínicos mínimos enumerados en el criterio 5.
\item
  \textbf{Rationale:} heredado de RF-10 de la Entrega 1A; reformulado en
  la Tercera Entrega para cerrar OBS-04 (ambigüedad leve señalada por el
  docente en ``recomendaciones informativas''), sin ampliar el alcance
  original: sigue dependiendo de REQ-F-013 (historial clínico) y del
  mismo servicio de IA externo ya previsto desde la Entrega 1A. En esta
  revisión se añade además la restricción explícita sobre los datos que
  pueden salir del sistema (criterio 5), siguiendo el principio técnico
  de minimización de datos ya aplicado en REQ-F-025.
\item
  \textbf{Criterios de aceptación:}

  \begin{enumerate}
  \def\labelenumi{\arabic{enumi}.}
  \tightlist
  \item
    Dado un historial clínico existente para la mascota solicitada,
    cuando se invoca el endpoint de recomendaciones, entonces la
    respuesta incluye el campo \texttt{recomendaciones:\ s\allowbreak{}tring{[}{]}}
    y, por cada elemento, el campo
    \texttt{advertencia:\ "inf\allowbreak{}ormativa,\allowbreak{}\ no\ sustituye\ diagn\allowbreak{}óstico\ veterinari\allowbreak{}o"}.
  \item
    Dada una mascota sin historial clínico registrado, entonces la
    respuesta es \texttt{200\ OK} con \texttt{recomendaciones:\ {[}{]}},
    no un error.
  \item
    Dado un \texttt{ROLE\_\allowbreak{}DUENO}, cuando solicita recomendaciones para
    una mascota ajena, entonces la respuesta es \texttt{403}.
  \item
    Dado que el servicio de IA externo no responde o devuelve error,
    entonces la respuesta es \texttt{502\ Bad\ Gateway} con
    \texttt{Problem\allowbreak{}Detail}, siguiendo el mismo patrón que
    \texttt{External\allowbreak{}Api\allowbreak{}Exception}, sin exponer detalles internos del
    proveedor.
  \item
    \textbf{Restricción de datos enviados al servicio externo.} La carga
    útil enviada al proveedor de IA contiene exclusivamente: especie,
    raza, edad o fecha de nacimiento de la mascota, y los campos
    \texttt{motivo}, \texttt{diagnóstico}, \texttt{tratamiento} y
    \texttt{observaciones} de sus consultas. Queda explícitamente
    prohibido enviar: el nombre, el correo electrónico o el
    identificador del dueño; el nombre o el identificador del
    veterinario; el nombre o el identificador de la mascota; cualquier
    token JWT, cookie, clave de API propia o credencial del sistema; y
    cualquier campo no enumerado en la primera frase de este criterio.
    La clave del proveedor viaja únicamente en la cabecera de
    autenticación de la llamada saliente, nunca en el cuerpo ni en la
    URL.
  \item
    La comprobación del criterio 5 es observable: la carga útil enviada
    se puede capturar en una prueba con doble de prueba (mock) del
    cliente externo y comparar campo por campo contra la lista
    permitida.
  \end{enumerate}
\item
  \textbf{Método de verificación:} previsto --- prueba de integración
  con un doble de prueba (mock) del servicio de IA, siguiendo el patrón
  real ya usado en \texttt{External\allowbreak{}Api\allowbreak{}Service\allowbreak{}Test}, que valide: (a) el
  contrato de entrada/salida (criterios 1 y 2), (b) el control de acceso
  por propietario (criterio 3),

  (c) la traducción de fallos del proveedor a \texttt{502} (criterio 4)
  y (d) la carga útil saliente campo por campo contra la lista permitida
  (criterios 5 y 6). No ejecutable hoy: no hay implementación.
\item
  \textbf{Trazabilidad:} HU-016 → CU-16 → depende de REQ-F-013 (ambas
  obligaciones); patrón de integración externa de referencia ya
  implementado en
  \texttt{com.\allowbreak{}biopet.\allowbreak{}integration.\allowbreak{}External\allowbreak{}Api\allowbreak{}Client}/\texttt{External\allowbreak{}Api\allowbreak{}Service}
  y en \texttt{com.\allowbreak{}biopet.\allowbreak{}exception.\allowbreak{}External\allowbreak{}Api\allowbreak{}Exception} →
  \texttt{Global\allowbreak{}Exception\allowbreak{}Handler.\allowbreak{}error\allowbreak{}Api\allowbreak{}Externa} →
  \texttt{Problem\allowbreak{}Type.\allowbreak{}BAD\_\allowbreak{}GATEWAY}; sistema externo identificado en
  \texttt{docs/\allowbreak{}diagrams/\allowbreak{}c4-\allowbreak{}contexto/\allowbreak{}C4-\allowbreak{}L1-\allowbreak{}contexto.\allowbreak{}md}; sección 2.6
  (servicio de IA declarado como \textbf{no integrado}).
\item
  \textbf{Estado:} pendiente
\item
  \textbf{Observaciones:} el proveedor concreto del servicio de IA sigue
  siendo una decisión de arquitectura no tomada; no existe ADR que la
  fije. El criterio 5 define el principio técnico mínimo de minimización
  de datos que deberá cumplir la implementación futura; \textbf{no}
  declara conformidad con ninguna norma legal de protección de datos,
  porque el proyecto no ha realizado ese análisis.
\end{itemize}

\subsubsection{REQ-F-018 --- Comprobantes de pago
digitales}\label{req-f-018-comprobantes-de-pago-digitales}

\begin{itemize}
\tightlist
\item
  \textbf{Tipo:} Funcional
\item
  \textbf{Prioridad:} Should
\item
  \textbf{Enunciado:} El sistema deberá generar un comprobante de pago
  digital en formato PDF por cada servicio veterinario facturado,
  registrando el método de pago empleado y asociando el comprobante al
  dueño y a la mascota atendidos.
\item
  \textbf{Rationale:} heredado de RF-11 y RF-12 de la Entrega 1A.
\item
  \textbf{Criterios de aceptación:}

  \begin{enumerate}
  \def\labelenumi{\arabic{enumi}.}
  \tightlist
  \item
    Dado un servicio veterinario registrado y un método de pago válido,
    cuando se emite el comprobante, entonces se genera un archivo PDF
    descargable que incluye identificador del comprobante, fecha, dueño,
    mascota, detalle del servicio, importe y método de pago.
  \item
    El identificador del comprobante es único y correlativo.
  \item
    Dado un \texttt{ROLE\_\allowbreak{}DUENO}, cuando solicita sus propios
    comprobantes, entonces los obtiene; los de otro dueño responden
    \texttt{403}.
  \item
    Dado un método de pago no admitido, entonces la respuesta es
    \texttt{422} con \texttt{Problem\allowbreak{}Detail}.
  \end{enumerate}
\item
  \textbf{Método de verificación:} previsto --- prueba de integración
  que emita un comprobante y valide la estructura del PDF generado y la
  unicidad del correlativo, más prueba de control de acceso por
  propietario. No ejecutable hoy.
\item
  \textbf{Trazabilidad:} HU-017 → CU-17 → entidades conceptuales
  \texttt{Factura} y \texttt{Detalle\_\allowbreak{}Factura} del DER de la Entrega 1A.
  Para el criterio 2 existe ya en la base de datos el generador de
  códigos correlativos
  \texttt{db/\allowbreak{}procs/\allowbreak{}fn\_\allowbreak{}siguiente\_\allowbreak{}numer\allowbreak{}o\_\allowbreak{}ficha.\allowbreak{}sql}
  (\texttt{Procedimiento\allowbreak{}Biopet\allowbreak{}Repository.\allowbreak{}siguiente\allowbreak{}Numero\allowbreak{}Ficha},
  secuencia \texttt{seq\_\allowbreak{}ficha\_\allowbreak{}biopet}), reutilizable como numerador.
  Sin controlador, servicio, entidad ni migración de facturación en el
  repositorio actual.
\item
  \textbf{Estado:} pendiente
\end{itemize}

\subsubsection{REQ-F-019 --- Reportes estadísticos
exportables}\label{req-f-019-reportes-estaduxedsticos-exportables}

\begin{itemize}
\tightlist
\item
  \textbf{Tipo:} Funcional
\item
  \textbf{Prioridad:} Could
\item
  \textbf{Enunciado:} El sistema deberá permitir generar reportes
  estadísticos de actividad de la clínica para un rango de fechas dado y
  exportarlos en formato PDF y Excel.
\item
  \textbf{Rationale:} heredado de RF-14 de la Entrega 1A.
\item
  \textbf{Criterios de aceptación:}

  \begin{enumerate}
  \def\labelenumi{\arabic{enumi}.}
  \tightlist
  \item
    Dado un rango de fechas válido, cuando se solicita el reporte,
    entonces se obtienen los indicadores del periodo: mascotas activas,
    citas programadas, consultas y vacunas dentro del rango, y mascotas
    sin consulta.
  \item
    Dado un rango invertido (fecha inicial posterior a la final),
    entonces los indicadores del periodo resultan en cero, sin error.
  \item
    Dado un reporte generado, entonces puede descargarse como PDF y como
    Excel, con los mismos valores en ambos formatos.
  \item
    Dado un usuario cuyo rol no es \texttt{ROLE\_\allowbreak{}ADMIN}, cuando solicita
    el reporte global de la clínica, entonces la respuesta es
    \texttt{403}.
  \item
    La agregación multi-tabla se implementa como procedimiento
    almacenado (ver REQ-NF-013), no como JPQL con joins.
  \end{enumerate}
\item
  \textbf{Método de verificación:} pruebas de integración
  \texttt{Procedimientos\allowbreak{}Biopet\allowbreak{}Integration\allowbreak{}Test.\allowbreak{}reporte\allowbreak{}Dashboard\_\allowbreak{}con\allowbreak{}Datos\_\allowbreak{}calcula\allowbreak{}Indicadores}
  y \texttt{reporte\allowbreak{}Dashboard\_\allowbreak{}rango\allowbreak{}Invertido\_\allowbreak{}no\allowbreak{}Cuenta\allowbreak{}Dentro\allowbreak{}Del\allowbreak{}Rango}
  sobre PostgreSQL real (criterios 1, 2 y 5). Criterios 3 y 4:
  \textbf{verificación prevista} --- prueba de controlador sobre el
  endpoint de reportes y verificación del contenido de los archivos
  exportados, a escribir cuando existan.
\item
  \textbf{Trazabilidad:} HU-018 → CU-18 → procedimiento almacenado
  \texttt{db/\allowbreak{}procs/\allowbreak{}fn\_\allowbreak{}reporte\_\allowbreak{}dashboa\allowbreak{}rd.\allowbreak{}sql} →
  \texttt{Procedimiento\allowbreak{}Biopet\allowbreak{}Repository.\allowbreak{}reporte\allowbreak{}Dashboard} → proyección
  \texttt{com.\allowbreak{}biopet.\allowbreak{}repository.\allowbreak{}Reporte\allowbreak{}Dashboard} →
  \texttt{Backend/\allowbreak{}src/\allowbreak{}test/\allowbreak{}java/\allowbreak{}com/\allowbreak{}biopet/\allowbreak{}repository/\allowbreak{}Procedimientos\allowbreak{}Biopet\allowbreak{}Integration\allowbreak{}Test.\allowbreak{}java};
  catálogo \texttt{docs/\allowbreak{}basedatos/\allowbreak{}CATALOGO-\allowbreak{}SP.\allowbreak{}md} (categoría ``3)
  Reporte''). Sin controlador, servicio de exportación ni endpoint en el
  repositorio actual.
\item
  \textbf{Estado:} pendiente
\item
  \textbf{Observaciones:} la capa de datos del reporte existe y está
  probada (procedimiento almacenado, repositorio, proyección y prueba de
  integración), pero \textbf{ninguna} de las obligaciones visibles del
  requisito ---endpoint, exportación a PDF y a Excel, restricción por
  rol--- está implementada, por lo que el estado sigue siendo
  \texttt{pendiente} y no \texttt{implementado}. Esa base de datos se
  registra aquí en Trazabilidad para que el trabajo ya hecho no quede
  oculto.
\end{itemize}

\subsubsection{REQ-F-020 --- Auditoría de operaciones del
sistema}\label{req-f-020-auditoruxeda-de-operaciones-del-sistema}

\begin{itemize}
\tightlist
\item
  \textbf{Tipo:} Funcional
\item
  \textbf{Prioridad:} Should
\item
  \textbf{Enunciado:} El sistema deberá registrar, para cada operación
  relevante ---los eventos de autenticación y las operaciones de
  creación, actualización y baja sobre mascotas, citas, consultas,
  vacunas y usuarios---, el sujeto que la ejecuta, la marca de tiempo,
  la dirección IP de origen y el resultado de la operación.
\item
  \textbf{Rationale:} heredado de RF-15 de la Entrega 1A, reforzado por
  el control OWASP A09 exigido por el bloque C.2 de la Guía.
\item
  \textbf{Criterios de aceptación:}

  \begin{enumerate}
  \def\labelenumi{\arabic{enumi}.}
  \tightlist
  \item
    Dado un evento de autenticación (login exitoso, login fallido,
    bloqueo por rate limit, refresh exitoso, refresh fallido, logout,
    uso de token revocado), cuando ocurre, entonces se emite una línea
    de log estructurada con prefijo \texttt{AUTH\_\allowbreak{}AUDIT}, el tipo de
    evento, la IP, la marca de tiempo UTC y el sujeto.
  \item
    Los valores nulos se normalizan a \texttt{unknown} y los caracteres
    de control se eliminan del mensaje, de modo que un sujeto manipulado
    no pueda inyectar líneas de log falsas (\emph{log forging}).
  \item
    La línea de log no contiene contraseñas, hashes ni tokens.
  \item
    Dada una operación de creación, actualización o baja sobre mascotas,
    citas, consultas, vacunas o usuarios, cuando se ejecuta, entonces se
    registra un evento de auditoría con sujeto, marca de tiempo, IP,
    recurso afectado y resultado.
  \end{enumerate}
\item
  \textbf{Método de verificación:} pruebas automatizadas
  \texttt{Authentication\allowbreak{}Audit\allowbreak{}Service\allowbreak{}Test} (10 pruebas):
  \texttt{login\allowbreak{}Exitoso\allowbreak{}Registra\allowbreak{}Evento\allowbreak{}Estructurado},
  \texttt{login\allowbreak{}Fallido\allowbreak{}Registra\allowbreak{}Evento\allowbreak{}Warn},
  \texttt{login\allowbreak{}Bloqueado\allowbreak{}Registra\allowbreak{}Evento\allowbreak{}Warn},
  \texttt{refresh\allowbreak{}Exitoso\allowbreak{}Registra\allowbreak{}Evento\allowbreak{}Info},
  \texttt{refresh\allowbreak{}Fallido\allowbreak{}Registra\allowbreak{}Evento\allowbreak{}Warn},
  \texttt{logout\allowbreak{}Exitoso\allowbreak{}Registra\allowbreak{}Evento\allowbreak{}Info},
  \texttt{token\allowbreak{}Revocado\allowbreak{}Registra\allowbreak{}Evento\allowbreak{}Warn} (criterio 1);
  \texttt{valores\allowbreak{}Nulos\allowbreak{}Se\allowbreak{}Normalizan\allowbreak{}Como\allowbreak{}Unknown} y
  \texttt{elimina\allowbreak{}Caracteres\allowbreak{}De\allowbreak{}Control\allowbreak{}Para\allowbreak{}Evitar\allowbreak{}Log\allowbreak{}Forging} (criterio 2);
  \texttt{no\allowbreak{}Registra\allowbreak{}Datos\allowbreak{}Sensibles} (criterio 3); más la captura de log
  real en \texttt{docs/\allowbreak{}mediciones/\allowbreak{}sec/\allowbreak{}A09-\allowbreak{}logging.\allowbreak{}md} y
  \texttt{docs/\allowbreak{}mediciones/\allowbreak{}sec/\allowbreak{}raw/\allowbreak{}A09-\allowbreak{}audit-\allowbreak{}logs.\allowbreak{}txt}. Criterio 4:
  \textbf{sin artefacto} --- ver Observaciones.
\item
  \textbf{Trazabilidad:} HU-019 → CU-19 →
  \texttt{com.\allowbreak{}biopet.\allowbreak{}security.\allowbreak{}Authentication\allowbreak{}Audit\allowbreak{}Service} (métodos
  \texttt{login\allowbreak{}Exitoso}, \texttt{login\allowbreak{}Fallido}, \texttt{login\allowbreak{}Bloqueado},
  \texttt{refresh\allowbreak{}Exitoso}, \texttt{refresh\allowbreak{}Fallido},
  \texttt{logout\allowbreak{}Exitoso}, \texttt{token\allowbreak{}Revocado}) → invocado desde
  \texttt{Auth\allowbreak{}Service} y \texttt{Jwt\allowbreak{}Authentication\allowbreak{}Filter} → prueba
  \texttt{Backend/\allowbreak{}src/\allowbreak{}test/\allowbreak{}java/\allowbreak{}com/\allowbreak{}biopet/\allowbreak{}security/\allowbreak{}Authentication\allowbreak{}Audit\allowbreak{}Service\allowbreak{}Test.\allowbreak{}java}
  → evidencia \texttt{docs/\allowbreak{}mediciones/\allowbreak{}sec/\allowbreak{}A09-\allowbreak{}logging.\allowbreak{}md}; requisito no
  funcional hermano REQ-NF-009.
\item
  \textbf{Estado:} implementado
\item
  \textbf{Observaciones:} el criterio 4 \textbf{no se cumple}: sólo
  existe auditoría de eventos de \emph{autenticación} (cubierta además
  por REQ-NF-009 y operativa en producción). No hay ningún interceptor,
  aspecto ni servicio que registre las operaciones CRUD de negocio sobre
  mascotas, citas, consultas, vacunas ni usuarios. Se marca
  \texttt{implementado} y no \texttt{pendiente} para no ocultar el
  trabajo de auditoría de autenticación ya entregado y verificado.
\end{itemize}

\subsubsection{REQ-F-021 --- Resumen de mascotas activas agrupadas por
especie}\label{req-f-021-resumen-de-mascotas-activas-agrupadas-por-especie}

\begin{itemize}
\tightlist
\item
  \textbf{Tipo:} Funcional
\item
  \textbf{Prioridad:} Should
\item
  \textbf{Enunciado:} El sistema deberá permitir consultar el total de
  mascotas activas agrupadas por especie; para roles distintos de
  \texttt{ROLE\_\allowbreak{}ADMIN} el resultado deberá restringirse siempre a las
  mascotas del propio usuario autenticado, y para el rol
  \texttt{ROLE\_\allowbreak{}ADMIN} deberá poder filtrarse por un dueño específico o
  consultarse sin filtro para obtener el total global.
\item
  \textbf{Rationale:} requisito de la Tercera Entrega, exigido por el
  bloque A.2.2 de la Guía como ejemplo de operación agregada
  (\texttt{GROUP\ BY}) que debe implementarse obligatoriamente vía
  función/procedimiento almacenado, no vía JPQL. Se numera 021 (no 013)
  porque los identificadores 013 a 020 ya estaban ocupados por los
  requisitos heredados de la Entrega 1A, fijados en
  \texttt{Historias\allowbreak{}Usuario.\allowbreak{}md}/\texttt{Casos\allowbreak{}De\allowbreak{}Uso.\allowbreak{}md} antes de que este
  requisito existiera.
\item
  \textbf{Criterios de aceptación:}

  \begin{enumerate}
  \def\labelenumi{\arabic{enumi}.}
  \tightlist
  \item
    Dadas dos mascotas activas de un mismo dueño con especies distintas
    (``Perro'' y ``Gato''), cuando se solicita el resumen filtrado por
    ese dueño, entonces se obtienen exactamente dos filas, una por
    especie, cada una con total 1.
  \item
    Dado un identificador de dueño inexistente, entonces el resultado es
    una lista vacía, no un error.
  \item
    Dado un usuario cuyo rol no es \texttt{ROLE\_\allowbreak{}ADMIN}, cuando invoca
    \texttt{GET\ /\allowbreak{}api/\allowbreak{}mascotas/\allowbreak{}resumen-\allowbreak{}especies} con o sin parámetro
    \texttt{duenio\allowbreak{}Id}, entonces el resultado se calcula siempre sobre su
    propio identificador: el parámetro enviado se ignora.
  \item
    Dado un usuario \texttt{ROLE\_\allowbreak{}ADMIN}, cuando invoca el endpoint sin
    parámetro \texttt{duenio\allowbreak{}Id}, entonces el resumen es global (todas
    las mascotas activas); cuando lo invoca con \texttt{duenio\allowbreak{}Id}, el
    resumen se limita a ese dueño.
  \item
    La agregación se ejecuta mediante un procedimiento almacenado
    versionado, no mediante JPQL con \texttt{GROUP\ BY}.
  \end{enumerate}
\item
  \textbf{Método de verificación:} prueba de integración
  \texttt{Resumen\allowbreak{}Especies\allowbreak{}Integration\allowbreak{}Test.\allowbreak{}resumen\allowbreak{}Agrupa\allowbreak{}Por\allowbreak{}Especie\allowbreak{}YFiltra\allowbreak{}Por\allowbreak{}Duenio}
  sobre PostgreSQL real con Testcontainers (criterio 1) y
  \texttt{Procedimientos\allowbreak{}Biopet\allowbreak{}Integration\allowbreak{}Test.\allowbreak{}resumen\allowbreak{}Por\allowbreak{}Especie\_\allowbreak{}duenio\allowbreak{}Inexistente\_\allowbreak{}devue\allowbreak{}lve\allowbreak{}Lista\allowbreak{}Vacia}
  (criterio 2). Criterios 3 y 4: inspección de
  \texttt{Mascota\allowbreak{}Service.\allowbreak{}resumen\allowbreak{}Por\allowbreak{}Especie}, cuya única regla es
  \texttt{Long\ duenio\allowbreak{}Id\allowbreak{}Efectivo\ =\ (usuari\allowbreak{}o\allowbreak{}Autenticado.\allowbreak{}get\allowbreak{}Rol()\ ==\ Rol.\allowbreak{}ROLE\_\allowbreak{}ADMIN)\ ?\ dueni\allowbreak{}o\allowbreak{}Id\allowbreak{}Solicitado\ :\ usuar\allowbreak{}io\allowbreak{}Autenticado.\allowbreak{}get\allowbreak{}Id();}.
  Criterio 5: \texttt{docs/\allowbreak{}basedatos/\allowbreak{}CATALOGO-\allowbreak{}SP.\allowbreak{}md} y
  \texttt{scripts/\allowbreak{}audit-\allowbreak{}sql-\allowbreak{}dynamic.\allowbreak{}sh} (ver REQ-NF-013).
\item
  \textbf{Trazabilidad:} HU-020 → CU-20 →
  \texttt{Mascota\allowbreak{}Controller.\allowbreak{}resumen\allowbreak{}Por\allowbreak{}Especies} →
  \texttt{Mascota\allowbreak{}Service.\allowbreak{}resumen\allowbreak{}Por\allowbreak{}Especie} →
  \texttt{Procedimiento\allowbreak{}Biopet\allowbreak{}Repository.\allowbreak{}resumen\allowbreak{}Por\allowbreak{}Especie} →
  procedimiento
  \texttt{db/\allowbreak{}procs/\allowbreak{}fn\_\allowbreak{}resumen\_\allowbreak{}mascota\allowbreak{}s\_\allowbreak{}por\_\allowbreak{}especie.\allowbreak{}sql} →
  \texttt{GET\ /\allowbreak{}api/\allowbreak{}mascotas/\allowbreak{}resumen-\allowbreak{}especies}; proyección
  \texttt{com.\allowbreak{}biopet.\allowbreak{}repository.\allowbreak{}Resumen\allowbreak{}Especie}; DTO
  \texttt{Resumen\allowbreak{}Especie\allowbreak{}Response}; pruebas
  \texttt{Resumen\allowbreak{}Especies\allowbreak{}Integration\allowbreak{}Test} y
  \texttt{Procedimientos\allowbreak{}Biopet\allowbreak{}Integration\allowbreak{}Test}; catálogo
  \texttt{docs/\allowbreak{}basedatos/\allowbreak{}CATALOGO-\allowbreak{}SP.\allowbreak{}md} (categoría ``2) Agregado'').
\item
  \textbf{Estado:} verificado
\item
  \textbf{Observaciones:} los criterios 3 y 4 se verifican por
  inspección de las cuatro líneas citadas de
  \texttt{Mascota\allowbreak{}Service.\allowbreak{}resumen\allowbreak{}Por\allowbreak{}Especie}; no existe prueba
  automatizada de controlador que ejercite el endpoint por rol. Queda
  registrado como mejora de cobertura, no como defecto funcional.
\end{itemize}

\subsubsection{REQ-F-022 --- Notificaciones al usuario por correo
electrónico (RF-07
recuperado)}\label{req-f-022-notificaciones-al-usuario-por-correo-electruxf3nico-rf-07-recuperado}

\begin{itemize}
\tightlist
\item
  \textbf{Tipo:} Funcional
\item
  \textbf{Prioridad:} Could
\item
  \textbf{Enunciado:} Al ocurrir un evento relevante para la cuenta de
  un usuario (por ejemplo, un registro exitoso), el sistema deberá
  enviar una notificación por correo electrónico al usuario afectado a
  través de un servicio de correo externo, de forma asíncrona, sin
  bloquear ni condicionar la respuesta HTTP de la operación que originó
  el evento.
\item
  \textbf{Rationale:} corresponde al RF-07 original de la Entrega 1A
  (``Servicio de Correos''), documentado como sistema externo en
  \texttt{docs/\allowbreak{}diagrams/\allowbreak{}c4-\allowbreak{}contexto/\allowbreak{}C4-\allowbreak{}L1-\allowbreak{}contexto.\allowbreak{}md} pero sin
  requisito \texttt{REQ-\allowbreak{}F} formal hasta la Tercera Entrega. Se incorporó
  para cerrar OBS-02 sin duplicar \texttt{REQ-\allowbreak{}F-\allowbreak{}007} (funcionalidad
  distinta ya implementada). Por eso se numera 022, el siguiente
  identificador libre después de \texttt{REQ-\allowbreak{}F-\allowbreak{}021}.
\item
  \textbf{Criterios de aceptación:}

  \begin{enumerate}
  \def\labelenumi{\arabic{enumi}.}
  \tightlist
  \item
    Dado un registro de usuario exitoso y un servicio de correo
    disponible, cuando se completa el registro, entonces se envía un
    correo de bienvenida a la dirección registrada.
  \item
    Dado que el servicio de correo externo no está disponible, cuando
    ocurre un evento notificable, entonces la operación de dominio que
    lo originó (por ejemplo, el registro) completa exitosamente con su
    código HTTP normal, y el fallo de envío queda registrado en el log
    del sistema sin propagarse como error al cliente.
  \item
    El envío es asíncrono: el tiempo de respuesta de la operación de
    dominio no depende del tiempo de respuesta del servicio de correo.
  \item
    El cuerpo del correo no incluye contraseñas, hashes ni tokens.
  \end{enumerate}
\item
  \textbf{Método de verificación:} previsto --- prueba de integración
  con un servidor SMTP de pruebas (por ejemplo, Mailhog o GreenMail) que
  confirme el envío en el caso exitoso (criterio 1), y prueba con doble
  de prueba que fuerce el fallo del proveedor y verifique que la
  operación de dominio sigue devolviendo su código normal (criterio 2),
  siguiendo el patrón real ya usado en
  \texttt{External\allowbreak{}Api\allowbreak{}Service\allowbreak{}Test.\allowbreak{}falla\allowbreak{}Al\allowbreak{}Escribir\allowbreak{}En\allowbreak{}Redis\allowbreak{}No\allowbreak{}Rompe\allowbreak{}La\allowbreak{}Respuesta}.
  No ejecutable hoy: no hay implementación.
\item
  \textbf{Trazabilidad:} HU-021 → CU-21 → sistema externo ``Servicio de
  Correos'' identificado en
  \texttt{docs/\allowbreak{}diagrams/\allowbreak{}c4-\allowbreak{}contexto/\allowbreak{}C4-\allowbreak{}L1-\allowbreak{}contexto.\allowbreak{}md}, sección
  ``Trazabilidad''; sección 2.6 de este documento (declarado \textbf{no
  integrado}). Sin controlador, servicio, configuración SMTP ni
  dependencia de correo en \texttt{Backend/\allowbreak{}pom.\allowbreak{}xml}.
\item
  \textbf{Estado:} pendiente
\item
  \textbf{Observaciones:} el proveedor de correo concreto (SMTP propio,
  SES, SendGrid u otro) es una decisión de arquitectura no tomada; no
  existe ADR que la fije.
\end{itemize}

\subsubsection{REQ-F-023 --- Gestión administrativa de
usuarios}\label{req-f-023-gestiuxf3n-administrativa-de-usuarios}

\begin{itemize}
\tightlist
\item
  \textbf{Tipo:} Funcional
\item
  \textbf{Prioridad:} Should
\item
  \textbf{Enunciado:} El sistema deberá permitir que un usuario con rol
  \texttt{ROLE\_\allowbreak{}ADMIN} liste, consulte por identificador, cree,
  actualice y dé de baja lógica cuentas de usuario de cualquier rol, sin
  permitir que un administrador modifique el rol de su propia cuenta.
\item
  \textbf{Rationale:} requisito de la actividad de Unidad IV; formaliza
  el CRUD administrativo de
  \texttt{Usuario\allowbreak{}Controller}/\texttt{Usuario\allowbreak{}Service}, funcionalidad
  distinta de REQ-F-007 (consulta del perfil propio,
  \texttt{/\allowbreak{}api/\allowbreak{}usuarios/\allowbreak{}me}), que no se modifica ni se duplica. Se
  numera 023, el siguiente identificador libre después de
  \texttt{REQ-\allowbreak{}F-\allowbreak{}022}.
\item
  \textbf{Criterios de aceptación:}

  \begin{enumerate}
  \def\labelenumi{\arabic{enumi}.}
  \tightlist
  \item
    Dado un \texttt{ROLE\_\allowbreak{}ADMIN}, cuando invoca
    \texttt{GET\ /\allowbreak{}api/\allowbreak{}usuarios}, entonces obtiene la página de usuarios;
    cuando invoca \texttt{GET\ /\allowbreak{}api/\allowbreak{}usuarios/\allowbreak{}\{id\}} con un
    identificador existente, obtiene ese usuario; con uno inexistente,
    \texttt{404} con \texttt{Problem\allowbreak{}Detail}.
  \item
    Dado un \texttt{ROLE\_\allowbreak{}ADMIN} y datos válidos, cuando invoca
    \texttt{POST\ /\allowbreak{}api/\allowbreak{}usuarios}, entonces la respuesta es
    \texttt{201\ Created} con el usuario creado; con un correo ya
    registrado, \texttt{409\ Conflict}; sin contraseña, \texttt{400};
    con campos inválidos, \texttt{422}.
  \item
    Dado un \texttt{ROLE\_\allowbreak{}ADMIN}, cuando intenta cambiar el rol de su
    \textbf{propia} cuenta mediante \texttt{PUT\ /\allowbreak{}api/\allowbreak{}usuarios/\allowbreak{}\{id\}},
    entonces la respuesta es \texttt{403}.
  \item
    Dado un \texttt{ROLE\_\allowbreak{}ADMIN}, cuando invoca
    \texttt{DELETE\ /\allowbreak{}api/\allowbreak{}usuarios/\allowbreak{}\{id\}}, entonces la respuesta es
    \texttt{204} y el usuario queda con \texttt{activo\ =\ false} (baja
    lógica, sin borrado físico).
  \item
    Dado un usuario cuyo rol no es \texttt{ROLE\_\allowbreak{}ADMIN} (incluido
    \texttt{ROLE\_\allowbreak{}VETERINARIO}), cuando invoca cualquiera de estas cinco
    operaciones, entonces la respuesta es \texttt{403}; sin
    autenticación, \texttt{401}.
  \end{enumerate}
\item
  \textbf{Método de verificación:} pruebas automatizadas de
  \texttt{Usuario\allowbreak{}Controller\allowbreak{}Test} (16 pruebas):
  \texttt{listar\allowbreak{}Usuarios\allowbreak{}Como\allowbreak{}Admin}, \texttt{buscar\allowbreak{}Usuario\allowbreak{}Por\allowbreak{}Id},
  \texttt{buscar\allowbreak{}Usuario\allowbreak{}Inexistente\allowbreak{}Devuelve404} (criterio 1);
  \texttt{crear\allowbreak{}Usuario\allowbreak{}Devuelve201},
  \texttt{crear\allowbreak{}Usuario\allowbreak{}Con\allowbreak{}Email\allowbreak{}Duplicado\allowbreak{}Devuelve409},
  \texttt{crear\allowbreak{}Usuario\allowbreak{}Sin\allowbreak{}Password\allowbreak{}Devuelve400},
  \texttt{crear\allowbreak{}Usuario\allowbreak{}Con\allowbreak{}Campos\allowbreak{}Invalidos\allowbreak{}Devuelve422},
  \texttt{actualizar\allowbreak{}Usuario\allowbreak{}Devuelve200},
  \texttt{actualizar\allowbreak{}Usuario\allowbreak{}Inexistente\allowbreak{}Devuelve404} (criterio 2);
  \texttt{admin\allowbreak{}No\allowbreak{}Puede\allowbreak{}Escalar\allowbreak{}Su\allowbreak{}Propio\allowbreak{}Rol\allowbreak{}Devuelve403} (criterio 3);
  \texttt{eliminar\allowbreak{}Usuario\allowbreak{}Devuelve204\allowbreak{}YBaja\allowbreak{}Logica},
  \texttt{eliminar\allowbreak{}Usuario\allowbreak{}Inexistente\allowbreak{}Devuelve404} (criterio 4);
  \texttt{acceso\allowbreak{}Con\allowbreak{}Rol\allowbreak{}No\allowbreak{}Admin\allowbreak{}Devuelve403},
  \texttt{veterinario\allowbreak{}No\allowbreak{}Puede\allowbreak{}Crear\allowbreak{}Usuarios\allowbreak{}Devuelve403},
  \texttt{acceso\allowbreak{}Sin\allowbreak{}Autenticacion\allowbreak{}Devuelve401} (criterio 5).
\item
  \textbf{Trazabilidad:} HU-022 → CU-22 →
  \texttt{Usuario\allowbreak{}Controller.\allowbreak{}\{listar,\allowbreak{}buscar,\allowbreak{}crear,\allowbreak{}actualizar,\allowbreak{}eliminar\}}
  → \texttt{Usuario\allowbreak{}Service} → \texttt{Usuario\allowbreak{}Repository} →
  \texttt{GET\ /\allowbreak{}api/\allowbreak{}usuarios}, \texttt{GET\ /\allowbreak{}api/\allowbreak{}usuarios/\allowbreak{}\{id\}},
  \texttt{POST\ /\allowbreak{}api/\allowbreak{}usuarios}, \texttt{PUT\ /\allowbreak{}api/\allowbreak{}usuarios/\allowbreak{}\{id\}},
  \texttt{DELETE\ /\allowbreak{}api/\allowbreak{}usuarios/\allowbreak{}\{id\}}; DTO
  \texttt{Usuario\allowbreak{}Request}/\texttt{Usuario\allowbreak{}Response}; prueba
  \texttt{Backend/\allowbreak{}src/\allowbreak{}test/\allowbreak{}java/\allowbreak{}com/\allowbreak{}biopet/\allowbreak{}Usuario\allowbreak{}Controller\allowbreak{}Test.\allowbreak{}java};
  matriz de permisos, sección 2.8.2.
\item
  \textbf{Estado:} verificado
\end{itemize}

\subsubsection{REQ-F-024 --- Gestión de
vacunas}\label{req-f-024-gestiuxf3n-de-vacunas}

\begin{itemize}
\tightlist
\item
  \textbf{Tipo:} Funcional
\item
  \textbf{Prioridad:} Should
\item
  \textbf{Enunciado:} El sistema deberá permitir registrar, consultar
  (de forma global, filtrada por mascota, o por identificador),
  actualizar y dar de baja lógica las vacunas aplicadas a una mascota,
  restringiendo la consulta de un usuario con rol \texttt{ROLE\_\allowbreak{}DUENO} a
  las vacunas de sus propias mascotas.
\item
  \textbf{Rationale:} requisito de la actividad de Unidad IV; formaliza
  el módulo \texttt{Vacuna\allowbreak{}Controller}/\texttt{Vacuna\allowbreak{}Service}, sin
  requisito formal previo en el SRS.
\item
  \textbf{Criterios de aceptación:}

  \begin{enumerate}
  \def\labelenumi{\arabic{enumi}.}
  \tightlist
  \item
    Dado un \texttt{ROLE\_\allowbreak{}ADMIN}, \texttt{ROLE\_\allowbreak{}VETERINARIO} o
    \texttt{ROLE\_\allowbreak{}AUXILIAR} y una mascota activa, cuando envía
    \texttt{POST\ /\allowbreak{}api/\allowbreak{}vacunas} con datos válidos, entonces la respuesta
    es \texttt{201\ Created}; con una mascota inexistente, \texttt{404};
    con campos inválidos, \texttt{422}.
  \item
    Dado un \texttt{ROLE\_\allowbreak{}DUENO}, cuando intenta crear una vacuna,
    entonces la respuesta es \texttt{403}.
  \item
    Dado un \texttt{ROLE\_\allowbreak{}DUENO} con vacunas propias y existiendo
    vacunas de mascotas ajenas, cuando invoca
    \texttt{GET\ /\allowbreak{}api/\allowbreak{}vacunas}, entonces ningún elemento devuelto
    corresponde a una mascota de otro dueño.
  \item
    Dado un \texttt{ROLE\_\allowbreak{}DUENO} y una vacuna de una mascota ajena,
    cuando invoca \texttt{GET\ /\allowbreak{}api/\allowbreak{}vacunas/\allowbreak{}\{id\}}, entonces la
    respuesta es \texttt{403}; con un identificador inexistente,
    \texttt{404}.
  \item
    Dado un \texttt{ROLE\_\allowbreak{}DUENO} y un \texttt{mascota\allowbreak{}Id} ajeno, cuando
    invoca \texttt{GET\ /\allowbreak{}api/\allowbreak{}vacunas/\allowbreak{}mascota/\allowbreak{}\{mascota\allowbreak{}Id\}}, entonces la
    respuesta es \texttt{403}; con una mascota propia, \texttt{200\ OK}
    con las vacunas de esa mascota.
  \item
    Dado un \texttt{DELETE\ /\allowbreak{}api/\allowbreak{}vacunas/\allowbreak{}\{id\}} autorizado, entonces la
    vacuna queda con \texttt{activo\ =\ false} y deja de aparecer en los
    listados.
  \end{enumerate}
\item
  \textbf{Método de verificación:} pruebas automatizadas de
  \texttt{Vacuna\allowbreak{}Controller\allowbreak{}Test} (9 pruebas):
  \texttt{admin\allowbreak{}Crea\allowbreak{}Vacuna\allowbreak{}Exitosamente},
  \texttt{crear\allowbreak{}Vacuna\allowbreak{}Con\allowbreak{}Mascota\allowbreak{}Inexistente\allowbreak{}Devuelve404},
  \texttt{crear\allowbreak{}Vacuna\allowbreak{}Con\allowbreak{}Campos\allowbreak{}Invalidos\allowbreak{}Devuelve422} (criterio 1);
  \texttt{dueno\allowbreak{}Intenta\allowbreak{}Crear\allowbreak{}Vacuna\allowbreak{}Es\allowbreak{}Rechazado\allowbreak{}Por\allowbreak{}Rol} (criterio 2);
  \texttt{dueno\allowbreak{}Solo\allowbreak{}Ve\allowbreak{}Vacunas\allowbreak{}De\allowbreak{}Sus\allowbreak{}Propias\allowbreak{}Mascotas} (criterio 3);
  \texttt{dueno\allowbreak{}Consulta\allowbreak{}Vacuna\allowbreak{}De\allowbreak{}Mascota\allowbreak{}Ajena\allowbreak{}Devuelve403},
  \texttt{buscar\allowbreak{}Vacuna\allowbreak{}Inexistente\allowbreak{}Devuelve404} (criterio 4);
  \texttt{admin\allowbreak{}Actualiza\allowbreak{}Vacuna\allowbreak{}Exitosamente},
  \texttt{admin\allowbreak{}Elimina\allowbreak{}Vacuna\allowbreak{}Exitosamente} (criterio 6). Criterio 5:
  inspección de \texttt{Vacuna\allowbreak{}Service.\allowbreak{}listar\allowbreak{}Por\allowbreak{}Mascota}, que invoca
  \texttt{verificar\allowbreak{}Acceso(usuario,\allowbreak{}\ mascota)} antes de consultar el
  repositorio, aplicando la misma regla que los criterios 3 y 4; y
  colección Postman
  \texttt{docs/\allowbreak{}postman/\allowbreak{}BIOPET-\allowbreak{}Vacunas.\allowbreak{}postman\_\allowbreak{}collectio\allowbreak{}n.\allowbreak{}json} (12
  requests, ejecutable con Newman).
\item
  \textbf{Trazabilidad:} HU-023 → CU-23 →
  \texttt{Vacuna\allowbreak{}Controller.\allowbreak{}\{listar,\allowbreak{}listar\allowbreak{}Por\allowbreak{}Mascota,\allowbreak{}buscar,\allowbreak{}crear,\allowbreak{}actualizar,\allowbreak{}eliminar\}}
  → \texttt{Vacuna\allowbreak{}Service} (\texttt{verificar\allowbreak{}Acceso}) →
  \texttt{Vacuna\allowbreak{}Repository.\allowbreak{}find\allowbreak{}All\allowbreak{}By\allowbreak{}Mascota\_\allowbreak{}Duenio\_\allowbreak{}Id\allowbreak{}And\allowbreak{}Activo\allowbreak{}True} →
  \texttt{GET\ /\allowbreak{}api/\allowbreak{}vacunas},
  \texttt{GET\ /\allowbreak{}api/\allowbreak{}vacunas/\allowbreak{}mascota/\allowbreak{}\{mascota\allowbreak{}Id\}},
  \texttt{GET\ /\allowbreak{}api/\allowbreak{}vacunas/\allowbreak{}\{id\}}, \texttt{POST\ /\allowbreak{}api/\allowbreak{}vacunas},
  \texttt{PUT\ /\allowbreak{}api/\allowbreak{}vacunas/\allowbreak{}\{id\}},
  \texttt{DELETE\ /\allowbreak{}api/\allowbreak{}vacunas/\allowbreak{}\{id\}}; interfaz de usuario
  \texttt{frontend/\allowbreak{}src/\allowbreak{}app/\allowbreak{}features/\allowbreak{}vacunas.\allowbreak{}component.\allowbreak{}ts}; matriz de
  permisos, sección 2.8.6.
\item
  \textbf{Estado:} verificado
\item
  \textbf{Observaciones:} la operación
  \texttt{GET\ /\allowbreak{}api/\allowbreak{}vacunas/\allowbreak{}mascota/\allowbreak{}\{mascota\allowbreak{}Id\}} no tiene prueba
  automatizada dedicada en \texttt{Vacuna\allowbreak{}Controller\allowbreak{}Test}; su criterio 5
  se sostiene por inspección de código y por la colección Postman, no
  por una prueba JUnit. La corrida Newman archivada
  (\texttt{docs/\allowbreak{}mediciones/\allowbreak{}postman/\allowbreak{}newman-\allowbreak{}report.\allowbreak{}json}, 2026-07-31) es
  \textbf{anterior} a la colección de vacunas y no puede citarse como
  evidencia de ejecución de esta última. Queda registrado como mejora de
  cobertura pendiente.
\end{itemize}

\subsubsection{REQ-F-025 --- Consulta de información externa de
especies}\label{req-f-025-consulta-de-informaciuxf3n-externa-de-especies}

\begin{itemize}
\tightlist
\item
  \textbf{Tipo:} Funcional
\item
  \textbf{Prioridad:} Could
\item
  \textbf{Enunciado:} El sistema deberá permitir a un usuario
  autenticado consultar información de una especie (nombre científico,
  hábitat y dieta) desde un servicio externo, utilizando una caché Redis
  para evitar llamadas repetidas al servicio externo dentro de una
  ventana de tiempo configurable, y enviando al servicio externo
  únicamente el nombre de la especie consultada.
\item
  \textbf{Rationale:} requisito de la actividad de Unidad IV;
  integración con un proveedor externo (API Ninjas, \emph{Animals API})
  vía \texttt{External\allowbreak{}Api\allowbreak{}Client}, con patrón \emph{cache-aside} en
  \texttt{External\allowbreak{}Api\allowbreak{}Service} (TTL configurable mediante
  \texttt{app.\allowbreak{}external-\allowbreak{}api.\allowbreak{}cache-\allowbreak{}ttl-\allowbreak{}seconds}, valor por defecto 600 s).
\item
  \textbf{Criterios de aceptación:}

  \begin{enumerate}
  \def\labelenumi{\arabic{enumi}.}
  \tightlist
  \item
    Dada una especie ya consultada dentro de la ventana de TTL, cuando
    se vuelve a consultar, entonces la respuesta se sirve desde Redis
    con \texttt{fuente\ =\ "cache"} y \textbf{no} se invoca al servicio
    externo.
  \item
    Dada una especie no cacheada, cuando se consulta, entonces se invoca
    al servicio externo, la respuesta se devuelve con
    \texttt{fuente\ =\ "api-\allowbreak{}ninjas"} y se almacena en Redis bajo la
    clave
    \texttt{external-\allowbreak{}api:animal:\textless{}esp\allowbreak{}ecie\ normalizada\textgreater{}}
    con el TTL configurado.
  \item
    Dado que el servicio externo no devuelve resultados para la especie,
    entonces la respuesta es \texttt{502\ Bad\ Gateway} con
    \texttt{Problem\allowbreak{}Detail} de tipo
    \texttt{urn:biopet:error:\allowbreak{}external-\allowbreak{}service}.
  \item
    Dado que la escritura en la caché falla, entonces la respuesta al
    cliente se entrega igualmente: el fallo de caché no rompe la
    operación.
  \item
    Dado un valor cacheado corrupto (JSON no deserializable), entonces
    se señala el error de forma controlada en lugar de devolver datos
    inválidos.
  \item
    \textbf{Restricción de datos enviados al servicio externo.} La
    llamada saliente contiene únicamente el nombre de la especie
    normalizado (\texttt{especie.\allowbreak{}trim().\allowbreak{}to\allowbreak{}Lower\allowbreak{}Case()}) como parámetro
    \texttt{name}; no se envía ningún dato de usuario, mascota, dueño ni
    veterinario. La clave del proveedor viaja solo en la cabecera
    \texttt{X-\allowbreak{}Api-\allowbreak{}Key}, nunca en la URL ni en el cuerpo.
  \end{enumerate}
\item
  \textbf{Método de verificación:} pruebas automatizadas de
  \texttt{External\allowbreak{}Api\allowbreak{}Service\allowbreak{}Test} (6 pruebas,
  \texttt{Backend/\allowbreak{}src/\allowbreak{}test/\allowbreak{}java/\allowbreak{}com/\allowbreak{}biopet/\allowbreak{}integration/}):
  \texttt{cache\allowbreak{}Hit\allowbreak{}Devuelve\allowbreak{}Datos\allowbreak{}Desde\allowbreak{}Redis\allowbreak{}Sin\allowbreak{}Llamar\allowbreak{}Api\allowbreak{}Externa} (criterio
  1),
  \texttt{cache\allowbreak{}Miss\allowbreak{}Con\allowbreak{}Taxonomia\allowbreak{}YCaracteristicas\allowbreak{}Completas\allowbreak{}Guarda\allowbreak{}En\allowbreak{}Cache} y
  \texttt{cache\allowbreak{}Miss\allowbreak{}Sin\allowbreak{}Taxonomia\allowbreak{}Ni\allowbreak{}Caracteristicas\allowbreak{}Devuelve\allowbreak{}Campos\allowbreak{}Nulos}
  (criterio 2), \texttt{cache\allowbreak{}Miss\allowbreak{}Con\allowbreak{}Resultados\allowbreak{}Vacios\allowbreak{}Lanza\allowbreak{}Excepcion}
  (criterio 3), \texttt{falla\allowbreak{}Al\allowbreak{}Escribir\allowbreak{}En\allowbreak{}Redis\allowbreak{}No\allowbreak{}Rompe\allowbreak{}La\allowbreak{}Respuesta}
  (criterio 4), \texttt{cache\allowbreak{}Con\allowbreak{}Json\allowbreak{}Corrupto\allowbreak{}Lanza\allowbreak{}Excepcion} (criterio
  5). Criterio 6: inspección de
  \texttt{External\allowbreak{}Api\allowbreak{}Client.\allowbreak{}buscar\allowbreak{}Por\allowbreak{}Especie}, cuya URL se compone
  únicamente con \texttt{query\allowbreak{}Param("name",\allowbreak{}\ especie)} y cuya cabecera
  es \texttt{X-\allowbreak{}Api-\allowbreak{}Key}. Medición empírica adicional en
  \texttt{docs/\allowbreak{}u4/\allowbreak{}evidencias/\allowbreak{}fred/\allowbreak{}redis-\allowbreak{}cache-\allowbreak{}comparacion.\allowbreak{}md}
  (comparación real de tiempos de respuesta con caché fría vs.~caché
  caliente).
\item
  \textbf{Trazabilidad:} HU-024 → CU-24 →
  \texttt{External\allowbreak{}Api\allowbreak{}Controller.\allowbreak{}info\allowbreak{}Especie} →
  \texttt{External\allowbreak{}Api\allowbreak{}Service.\allowbreak{}obtener\allowbreak{}Info\allowbreak{}Especie} →
  \texttt{External\allowbreak{}Api\allowbreak{}Client.\allowbreak{}buscar\allowbreak{}Por\allowbreak{}Especie} →
  \texttt{GET\ /\allowbreak{}api/\allowbreak{}externa/\allowbreak{}especies?especie=}; DTO
  \texttt{External\allowbreak{}Api\allowbreak{}Response}; excepción \texttt{External\allowbreak{}Api\allowbreak{}Exception}
  → \texttt{Global\allowbreak{}Exception\allowbreak{}Handler.\allowbreak{}error\allowbreak{}Api\allowbreak{}Externa} →
  \texttt{Problem\allowbreak{}Type.\allowbreak{}BAD\_\allowbreak{}GATEWAY}; configuración
  \texttt{app.\allowbreak{}external-\allowbreak{}api.\allowbreak{}base-\allowbreak{}url}, \texttt{app.\allowbreak{}external-\allowbreak{}api.\allowbreak{}key}
  (\texttt{APP\_\allowbreak{}EXTERNAL\_\allowbreak{}API\_\allowbreak{}KE\allowbreak{}Y} en \texttt{.\allowbreak{}env.\allowbreak{}example}),
  \texttt{app.\allowbreak{}external-\allowbreak{}api.\allowbreak{}cache-\allowbreak{}ttl-\allowbreak{}seconds}; interfaz externa
  declarada en la sección 2.6 y en la sección 2.7.2.
\item
  \textbf{Estado:} verificado
\end{itemize}

\subsubsection{REQ-F-026 --- Recuperación de
contraseña}\label{req-f-026-recuperaciuxf3n-de-contraseuxf1a}

\begin{itemize}
\tightlist
\item
  \textbf{Tipo:} Funcional
\item
  \textbf{Prioridad:} Should
\item
  \textbf{Enunciado:} El sistema deberá permitir que un usuario
  registrado que ha olvidado su contraseña solicite su restablecimiento
  indicando su correo electrónico, recibiendo un enlace con un token de
  un solo uso y vigencia limitada, mediante el cual deberá poder
  establecer una contraseña nueva sin intervención de un administrador.
\item
  \textbf{Rationale:} requisito nuevo incorporado en esta revisión. La
  retroalimentación del docente-director señaló su ausencia: hoy un
  usuario que olvida su contraseña \textbf{no tiene ninguna vía de
  recuperación} --- la única forma de restablecerla es que un
  \texttt{ROLE\_\allowbreak{}ADMIN} la sobrescriba mediante
  \texttt{PUT\ /\allowbreak{}api/\allowbreak{}usuarios/\allowbreak{}\{id\}} (REQ-F-023), lo que obliga a que un
  tercero conozca la contraseña nueva. Se documenta como requisito
  planificado, no como funcionalidad entregada.
\item
  \textbf{Criterios de aceptación:}

  \begin{enumerate}
  \def\labelenumi{\arabic{enumi}.}
  \tightlist
  \item
    Dado un correo correspondiente a una cuenta activa, cuando se
    solicita el restablecimiento, entonces se genera un token de un solo
    uso con vigencia limitada y se envía al correo registrado.
  \item
    Dado un correo \textbf{no} registrado, entonces la respuesta al
    cliente es idéntica a la del caso anterior (mismo código y mismo
    cuerpo), para no revelar qué correos existen en el sistema
    (enumeración de cuentas).
  \item
    Dado un token válido y no usado, cuando se envía junto con una
    contraseña nueva que cumple REQ-NF-016, entonces la contraseña se
    actualiza, se almacena cifrada con BCrypt y el token queda
    invalidado.
  \item
    Dado un token ya usado, expirado o inexistente, entonces la
    respuesta es \texttt{400} con \texttt{Problem\allowbreak{}Detail}, y la
    contraseña no se modifica.
  \item
    Tras un restablecimiento exitoso, las sesiones vigentes de ese
    usuario quedan revocadas: los access y refresh tokens emitidos antes
    del cambio responden \texttt{401}.
  \item
    Los endpoints de solicitud y de confirmación están sujetos a la
    misma limitación de intentos por IP que el login (REQ-NF-010).
  \end{enumerate}
\item
  \textbf{Método de verificación:} previsto --- pruebas de controlador
  con \texttt{Mock\allowbreak{}Mvc} siguiendo el patrón de
  \texttt{Auth\allowbreak{}Controller\allowbreak{}Test}, que cubran: respuesta uniforme para
  correo existente e inexistente (criterio 2), restablecimiento exitoso
  y revocación de sesiones previas (criterios 3 y 5), token
  usado/expirado (criterio 4) y bloqueo por rate limit (criterio 6); más
  la prueba de envío de correo descrita en REQ-F-022. No ejecutable hoy:
  no hay implementación.
\item
  \textbf{Trazabilidad:} depende de REQ-F-022 (envío de correo,
  \texttt{pendiente}) y de REQ-NF-016 (política de contraseñas).
  Reutilizará \texttt{Token\allowbreak{}Blacklist\allowbreak{}Service.\allowbreak{}revoke} para el criterio 5 y
  \texttt{Login\allowbreak{}Rate\allowbreak{}Limiter\allowbreak{}Service} para el criterio 6, ambos ya
  implementados. Sin historia de usuario ni caso de uso propios todavía:
  se deberán crear la historia y el caso de uso correspondientes, con el
  siguiente identificador libre de cada serie, al planificar su
  implementación. Sin endpoint, servicio, entidad de token ni migración
  en el repositorio actual: la búsqueda de
  \texttt{recuperar\textbar{}reset\textbar{}\allowbreak{}forgot\textbar{}olvid} sobre
  \texttt{Backend/\allowbreak{}src} y \texttt{frontend/\allowbreak{}src} no devuelve ninguna
  coincidencia.
\item
  \textbf{Estado:} pendiente
\item
  \textbf{Observaciones:} este requisito \textbf{no está implementado} y
  no debe presentarse como entregado. Se incorpora al SRS para que la
  ausencia quede declarada explícitamente en lugar de omitida. Su fila
  en \texttt{docs/\allowbreak{}trazabilidad/\allowbreak{}matriz.\allowbreak{}csv} no referencia HU ni CU porque
  ninguna existe; el respaldo declarado es el método de verificación
  previsto.
\end{itemize}

\subsubsection{3.2. Requisitos no funcionales
(REQ-NF)}\label{requisitos-no-funcionales-req-nf}

Los requisitos no funcionales usan la misma plantilla de la sección
1.3.2. En el campo \textbf{Tipo} se indica la categoría de calidad. En
el campo \textbf{Trazabilidad} no se declaran historias de usuario salvo
que existan realmente: un requisito no funcional traza legítimamente
hacia decisiones de arquitectura, configuración, clases, pruebas,
scripts, mediciones, evidencias, normas y documentación técnica.

\subsubsection{REQ-NF-001 --- Rendimiento del listado de
mascotas}\label{req-nf-001-rendimiento-del-listado-de-mascotas}

\begin{itemize}
\tightlist
\item
  \textbf{Tipo:} Rendimiento
\item
  \textbf{Prioridad:} Must
\item
  \textbf{Enunciado:} El tiempo de respuesta del listado de mascotas
  (\texttt{GET\ /\allowbreak{}api/\allowbreak{}mascotas}) no deberá superar 200 ms en el percentil
  95 con caché caliente, ni 500 ms en el percentil 95 con caché fría.
\item
  \textbf{Rationale:} heredado de RNF-01/RNF-WEB-04 de la Entrega 1A,
  con umbrales cuantitativos añadidos por el bloque C.1 de la Guía.
\item
  \textbf{Criterios de aceptación:}

  \begin{enumerate}
  \def\labelenumi{\arabic{enumi}.}
  \tightlist
  \item
    En una corrida de carga con caché caliente sobre
    \texttt{GET\ /\allowbreak{}api/\allowbreak{}mascotas}, el p95 medido es ≤ 200 ms.
  \item
    En una corrida de carga con caché fría sobre el mismo endpoint, el
    p95 medido es ≤ 500 ms.
  \item
    La tasa de error HTTP en ambas corridas es 0,0 \%.
  \item
    Las mediciones son reproducibles: los archivos crudos de cada
    corrida quedan archivados con fecha y con el identificador de la
    versión medida.
  \end{enumerate}
\item
  \textbf{Método de verificación:} medición con k6 --- 10 corridas
  reales (5 en caliente y 5 en frío) ejecutadas el 2026-09-03 sobre el
  código del tag \texttt{v1.\allowbreak{}0.\allowbreak{}0}. Resultado registrado en
  \texttt{docs/\allowbreak{}mediciones/\allowbreak{}perf/\allowbreak{}REPORT.\allowbreak{}md}: p95 entre 6,04 ms y 18,56 ms
  según la corrida (máximo 15,06 ms en caché caliente y 18,56 ms en
  caché fría), muy por debajo de los umbrales de 200 ms y 500 ms, con
  0,0 \% de error en todas las corridas.
\item
  \textbf{Trazabilidad:} archivos crudos
  \texttt{docs/\allowbreak{}mediciones/\allowbreak{}perf/\allowbreak{}k6-\allowbreak{}20260903\allowbreak{}T*-\allowbreak{}local-\allowbreak{}tls-\allowbreak{}v1.\allowbreak{}0.\allowbreak{}0-*.\allowbreak{}json}
  (10 archivos) y reporte \texttt{docs/\allowbreak{}mediciones/\allowbreak{}perf/\allowbreak{}REPORT.\allowbreak{}md};
  guiones de carga en \texttt{k6/}; endpoint medido
  \texttt{GET\ /\allowbreak{}api/\allowbreak{}mascotas} (\texttt{Mascota\allowbreak{}Controller.\allowbreak{}listar}); caché
  \texttt{@Cacheable(value\ =\ "\allowbreak{}mascotas")} en
  \texttt{Mascota\allowbreak{}Service.\allowbreak{}listar} con TTL de REQ-NF-012; procedencia de
  los datos en \texttt{docs/\allowbreak{}mediciones/\allowbreak{}DATA-\allowbreak{}PROVENANCE.\allowbreak{}md}.
\item
  \textbf{Estado:} verificado
\end{itemize}

\subsubsection{REQ-NF-002 --- Comunicación cifrada
obligatoria}\label{req-nf-002-comunicaciuxf3n-cifrada-obligatoria}

\begin{itemize}
\tightlist
\item
  \textbf{Tipo:} Seguridad
\item
  \textbf{Prioridad:} Must
\item
  \textbf{Enunciado:} Toda comunicación entre cliente y servidor deberá
  realizarse mediante HTTPS/TLS, y las respuestas servidas sobre HTTPS
  deberán incluir la cabecera \texttt{Strict-\allowbreak{}Transport-\allowbreak{}Security}.
\item
  \textbf{Rationale:} heredado de RNF-02/RNF-WEB-02 de la Entrega 1A;
  reforzado por el control OWASP A02 del bloque C.2 de la Guía.
\item
  \textbf{Criterios de aceptación:}

  \begin{enumerate}
  \def\labelenumi{\arabic{enumi}.}
  \tightlist
  \item
    Una conexión al puerto TLS negocia TLS 1.3 con una suite AEAD; la
    captura \texttt{curl\ -\allowbreak{}v} lo muestra explícitamente.
  \item
    Las respuestas servidas sobre HTTPS incluyen
    \texttt{Strict-\allowbreak{}Transport-\allowbreak{}Security:\ max-\allowbreak{}age=31536000;\allowbreak{}\ include\allowbreak{}Sub\allowbreak{}Domains;\allowbreak{}\ preload}.
  \item
    Las respuestas servidas sobre HTTP plano \textbf{no} incluyen la
    cabecera HSTS (evitando anunciarla sobre un canal no cifrado).
  \item
    Las respuestas incluyen además
    \texttt{X-\allowbreak{}Content-\allowbreak{}Type-\allowbreak{}Options:\ nosniff},
    \texttt{X-\allowbreak{}Frame-\allowbreak{}Options:\ DENY},
    \texttt{Referrer-\allowbreak{}Policy:\ no-\allowbreak{}referrer} y la política
    \texttt{Content-\allowbreak{}Security-\allowbreak{}Policy} configurada.
  \end{enumerate}
\item
  \textbf{Método de verificación:} captura real \texttt{curl\ -\allowbreak{}v}
  archivada en \texttt{docs/\allowbreak{}mediciones/\allowbreak{}sec/\allowbreak{}raw/\allowbreak{}A02-\allowbreak{}tls.\allowbreak{}txt} (TLSv1.3,
  \texttt{TLS\_\allowbreak{}AES\_\allowbreak{}256\_\allowbreak{}GCM\_\allowbreak{}SHA3\allowbreak{}84}) y documentada en
  \texttt{docs/\allowbreak{}mediciones/\allowbreak{}sec/\allowbreak{}A02-\allowbreak{}cryptography-\allowbreak{}tls.\allowbreak{}md} (criterio 1);
  pruebas automatizadas
  \texttt{Security\allowbreak{}Headers\allowbreak{}Test.\allowbreak{}respuesta\allowbreak{}Protegida\allowbreak{}Incluye\allowbreak{}Cabeceras\allowbreak{}De\allowbreak{}Seguridad}
  (criterios 2 y 4),
  \texttt{Security\allowbreak{}Headers\allowbreak{}Test.\allowbreak{}peticion\allowbreak{}Http\allowbreak{}No\allowbreak{}Incluye\allowbreak{}Hsts} (criterio 3) y
  \texttt{respuestas\allowbreak{}Publicas\allowbreak{}Tambien\allowbreak{}Incluyen\allowbreak{}Cabeceras\allowbreak{}Basicas}; capturas
  de cabeceras en
  \texttt{docs/\allowbreak{}mediciones/\allowbreak{}sec/\allowbreak{}raw/\allowbreak{}https-\allowbreak{}8443-\allowbreak{}headers.\allowbreak{}txt} y
  \texttt{http-\allowbreak{}8080-\allowbreak{}headers.\allowbreak{}txt}.
\item
  \textbf{Trazabilidad:} \texttt{Security\allowbreak{}Config} (bloque
  \texttt{.\allowbreak{}headers(...)}), \texttt{Tomcat\allowbreak{}Dual\allowbreak{}Connector\allowbreak{}Config} (conector
  dual HTTP/HTTPS),
  \texttt{Backend/\allowbreak{}src/\allowbreak{}main/\allowbreak{}resources/\allowbreak{}application-\allowbreak{}tls.\allowbreak{}yml},
  \texttt{docker-\allowbreak{}compose.\allowbreak{}tls.\allowbreak{}yml},
  \texttt{scripts/\allowbreak{}generate-\allowbreak{}dev-\allowbreak{}keystore.\allowbreak{}sh},
  \texttt{docs/\allowbreak{}despliegue/\allowbreak{}Caddyfile} y \texttt{nginx-\allowbreak{}render.\allowbreak{}conf};
  validación documentada en
  \texttt{docs/\allowbreak{}despliegue/\allowbreak{}VALIDACION-\allowbreak{}HEADERS-\allowbreak{}Z4.\allowbreak{}md}; prueba
  \texttt{Backend/\allowbreak{}src/\allowbreak{}test/\allowbreak{}java/\allowbreak{}com/\allowbreak{}biopet/\allowbreak{}Security\allowbreak{}Headers\allowbreak{}Test.\allowbreak{}java}.
\item
  \textbf{Estado:} verificado
\item
  \textbf{Observaciones:} la evidencia TLS se generó sobre el entorno de
  TLS de desarrollo (\texttt{docker-\allowbreak{}compose.\allowbreak{}tls.\allowbreak{}yml} con keystore
  generado por \texttt{scripts/\allowbreak{}generate-\allowbreak{}dev-\allowbreak{}keystore.\allowbreak{}sh}), no sobre el
  certificado del despliegue en Render. La configuración y las cabeceras
  verificadas son las mismas; el certificado y su cadena, no.
\end{itemize}

\subsubsection{REQ-NF-003 --- Expiración configurable de tokens
JWT}\label{req-nf-003-expiraciuxf3n-configurable-de-tokens-jwt}

\begin{itemize}
\tightlist
\item
  \textbf{Tipo:} Seguridad
\item
  \textbf{Prioridad:} Must
\item
  \textbf{Enunciado:} El sistema deberá emitir access tokens con
  expiración de 1 hora y refresh tokens con expiración de 7 días, ambos
  configurables mediante variables de entorno
  (\texttt{JWT\_\allowbreak{}EXPIRATION\_\allowbreak{}MS}, \texttt{JWT\_\allowbreak{}REFRESH\_\allowbreak{}EXPIRA\allowbreak{}TION\_\allowbreak{}MS}),
  sin valores fijos en el código.
\item
  \textbf{Rationale:} heredado de RNF-03/RNF-WEB-03 de la Entrega 1A.
\item
  \textbf{Criterios de aceptación:}

  \begin{enumerate}
  \def\labelenumi{\arabic{enumi}.}
  \tightlist
  \item
    Los valores por defecto configurados son 3 600 000 ms (1 hora) para
    el access token y 604 800 000 ms (7 días) para el refresh token.
  \item
    Ambos valores se leen de variables de entorno y pueden
    sobrescribirse sin recompilar.
  \item
    Un token cuyo \texttt{exp} ya pasó es rechazado por el servicio de
    validación.
  \item
    El código fuente no contiene ninguno de esos dos valores
    literalmente escrito en una clase Java.
  \end{enumerate}
\item
  \textbf{Método de verificación:} inspección de
  \texttt{Backend/\allowbreak{}src/\allowbreak{}main/\allowbreak{}resources/\allowbreak{}application.\allowbreak{}yml}, bloque
  \texttt{security.\allowbreak{}jwt}:
  \texttt{expiration-\allowbreak{}ms:\ \$\{JWT\_\allowbreak{}EXPIRATION\_\allowbreak{}M\allowbreak{}S:3600000\}} y
  \texttt{refresh-\allowbreak{}expiration-\allowbreak{}ms:\ \$\{JWT\_\allowbreak{}REFRESH\_\allowbreak{}EXP\allowbreak{}IRATION\_\allowbreak{}MS:604800\allowbreak{}000\}}
  (criterios 1, 2 y 4); prueba automatizada
  \texttt{Jwt\allowbreak{}Service\allowbreak{}Test.\allowbreak{}token\allowbreak{}Expirado\allowbreak{}Es\allowbreak{}Rechazado} (criterio 3).
\item
  \textbf{Trazabilidad:} \texttt{com.\allowbreak{}biopet.\allowbreak{}security.\allowbreak{}Jwt\allowbreak{}Service} (campos
  \texttt{@Value} de \texttt{security.\allowbreak{}jwt.\allowbreak{}expiration-\allowbreak{}ms} y
  \texttt{security.\allowbreak{}jwt.\allowbreak{}refresh-\allowbreak{}expiration-\allowbreak{}ms}),
  \texttt{Backend/\allowbreak{}src/\allowbreak{}main/\allowbreak{}resources/\allowbreak{}application.\allowbreak{}yml},
  \texttt{.\allowbreak{}env.\allowbreak{}example} (\texttt{JWT\_\allowbreak{}EXPIRATION\_\allowbreak{}MS},
  \texttt{JWT\_\allowbreak{}REFRESH\_\allowbreak{}EXPIRA\allowbreak{}TION\_\allowbreak{}MS}), prueba
  \texttt{Backend/\allowbreak{}src/\allowbreak{}test/\allowbreak{}java/\allowbreak{}com/\allowbreak{}biopet/\allowbreak{}security/\allowbreak{}Jwt\allowbreak{}Service\allowbreak{}Test.\allowbreak{}java},
  decisión \texttt{docs/\allowbreak{}adr/\allowbreak{}ADR-\allowbreak{}003-\allowbreak{}jwt-\allowbreak{}redis.\allowbreak{}md}; requisito funcional
  relacionado REQ-F-003.
\item
  \textbf{Estado:} verificado
\end{itemize}

\subsubsection{REQ-NF-004 --- Claims estándar del
JWT}\label{req-nf-004-claims-estuxe1ndar-del-jwt}

\begin{itemize}
\tightlist
\item
  \textbf{Tipo:} Seguridad
\item
  \textbf{Prioridad:} Must
\item
  \textbf{Enunciado:} Cada JWT emitido por el sistema deberá incluir los
  siete claims estándar \texttt{iss}, \texttt{sub}, \texttt{aud},
  \texttt{exp}, \texttt{nbf}, \texttt{iat} y \texttt{jti} conforme al
  RFC 7519, con \texttt{iss} y \texttt{aud} configurables por variable
  de entorno.
\item
  \textbf{Rationale:} refinamiento de RNF-03 exigido por el bloque A.1
  de la Guía.
\item
  \textbf{Criterios de aceptación:}

  \begin{enumerate}
  \def\labelenumi{\arabic{enumi}.}
  \tightlist
  \item
    Un access token emitido contiene los siete claims enumerados.
  \item
    Un refresh token emitido contiene los siete claims enumerados.
  \item
    El valor de \texttt{nbf} es igual al de \texttt{iat}.
  \item
    El valor de \texttt{jti} es distinto entre dos tokens emitidos
    consecutivamente.
  \item
    Un token cuyo \texttt{iss} no coincide con el configurado es
    rechazado; lo mismo para \texttt{aud}.
  \item
    Access token y refresh token se distinguen entre sí de forma
    inequívoca.
  \end{enumerate}
\item
  \textbf{Método de verificación:} pruebas automatizadas de
  \texttt{Jwt\allowbreak{}Service\allowbreak{}Test}: \texttt{access\allowbreak{}Token\allowbreak{}Contiene\allowbreak{}Los\allowbreak{}Siete\allowbreak{}Claims}
  (criterio 1), \texttt{refresh\allowbreak{}Token\allowbreak{}Contiene\allowbreak{}Los\allowbreak{}Siete\allowbreak{}Claims} (criterio
  2), \texttt{nbf\allowbreak{}Es\allowbreak{}Igual\allowbreak{}AIat} (criterio 3),
  \texttt{jti\allowbreak{}Es\allowbreak{}Unico\allowbreak{}Entre\allowbreak{}Tokens} (criterio 4),
  \texttt{issuer\allowbreak{}Incorrecto\allowbreak{}Es\allowbreak{}Rechazado} y
  \texttt{audience\allowbreak{}Incorrecta\allowbreak{}Es\allowbreak{}Rechazada} (criterio 5),
  \texttt{access\allowbreak{}YRefresh\allowbreak{}Se\allowbreak{}Distinguen\allowbreak{}Correctamente} (criterio 6).
\item
  \textbf{Trazabilidad:} \texttt{com.\allowbreak{}biopet.\allowbreak{}security.\allowbreak{}Jwt\allowbreak{}Service} (método
  \texttt{build\allowbreak{}Token}), configuración
  \texttt{security.\allowbreak{}jwt.\allowbreak{}issuer:\ \$\{JWT\_\allowbreak{}ISSUER:\allowbreak{}biopet-\allowbreak{}api\}} y
  \texttt{security.\allowbreak{}jwt.\allowbreak{}audience:\ \$\{JWT\_\allowbreak{}AUDI\allowbreak{}ENCE:biopet-\allowbreak{}frontend\}}
  en \texttt{application.\allowbreak{}yml}; prueba
  \texttt{Backend/\allowbreak{}src/\allowbreak{}test/\allowbreak{}java/\allowbreak{}com/\allowbreak{}biopet/\allowbreak{}security/\allowbreak{}Jwt\allowbreak{}Service\allowbreak{}Test.\allowbreak{}java};
  norma RFC 7519; decisión \texttt{docs/\allowbreak{}adr/\allowbreak{}ADR-\allowbreak{}003-\allowbreak{}jwt-\allowbreak{}redis.\allowbreak{}md}.
\item
  \textbf{Estado:} verificado
\end{itemize}

\subsubsection{REQ-NF-005 --- Interfaz
responsiva}\label{req-nf-005-interfaz-responsiva}

\begin{itemize}
\tightlist
\item
  \textbf{Tipo:} Usabilidad
\item
  \textbf{Prioridad:} Must
\item
  \textbf{Enunciado:} La interfaz deberá adaptarse correctamente a
  resoluciones comprendidas entre 320 px y 1440 px de ancho (móvil,
  tablet y escritorio), sin desbordamiento horizontal del contenido.
\item
  \textbf{Rationale:} heredado de RNF-04/RNF-WEB-01 de la Entrega 1A.
\item
  \textbf{Criterios de aceptación:}

  \begin{enumerate}
  \def\labelenumi{\arabic{enumi}.}
  \tightlist
  \item
    El documento servido declara
    \texttt{\textless{}meta\ name="view\allowbreak{}port"\ content="wi\allowbreak{}dth=device-\allowbreak{}width,\allowbreak{}\ ..."\textgreater{}}.
  \item
    La hoja de estilos publicada contiene reglas \texttt{@media} que
    adaptan el diseño a anchos móviles.
  \item
    En corridas Lighthouse con perfil móvil y con perfil de escritorio
    sobre las rutas \texttt{/\allowbreak{}login} y \texttt{/\allowbreak{}mascotas}, la puntuación
    de Performance es ≥ 90 en todos los casos.
  \item
    La hoja de estilos publicada define al menos un punto de corte para
    anchos móviles y el diseño de las rejillas de contenido usa unidades
    relativas, de modo que a 320 px el contenido refluye en una sola
    columna en lugar de desbordarse.
  \end{enumerate}
\item
  \textbf{Método de verificación:} pruebas automatizadas de
  \texttt{Frontend\allowbreak{}Responsiveness\allowbreak{}Test} (7 pruebas):
  \texttt{index\allowbreak{}Html\allowbreak{}Contiene\allowbreak{}Viewport\allowbreak{}Meta\allowbreak{}Tag} (criterio 1),
  \texttt{styles\allowbreak{}Css\allowbreak{}Contiene\allowbreak{}Media\allowbreak{}Queries\allowbreak{}Responsive},
  \texttt{build\allowbreak{}Incluye\allowbreak{}Archivo\allowbreak{}Styles\allowbreak{}Css},
  \texttt{index\allowbreak{}Html\allowbreak{}Referencia\allowbreak{}Styles\allowbreak{}Css} (criterio 2), más
  \texttt{frontend\allowbreak{}Build\allowbreak{}Existe\allowbreak{}Index\allowbreak{}Html},
  \texttt{index\allowbreak{}Html\allowbreak{}Contiene\allowbreak{}Doctype\allowbreak{}YHtml\allowbreak{}Valido} e
  \texttt{index\allowbreak{}Html\allowbreak{}Incluye\allowbreak{}Angular\allowbreak{}Bootstrap}. Criterio 3: 12 corridas
  Lighthouse del 2026-08-18
  (\texttt{docs/\allowbreak{}mediciones/\allowbreak{}lighthouse/\allowbreak{}lhci-\allowbreak{}20260818-\allowbreak{}0538-*.\allowbreak{}json}:
  perfiles móvil y desktop, rutas \texttt{/\allowbreak{}login} y \texttt{/\allowbreak{}mascotas},
  3 corridas cada combinación) con Performance 100 en desktop y 90--94
  en móvil, con el perfil móvil simulando Slow 4G + CPU 4x
  (\texttt{lighthouserc.\allowbreak{}js}). Criterio 4: inspección de
  \texttt{frontend/\allowbreak{}src/\allowbreak{}styles.\allowbreak{}css}, que declara
  \texttt{@media\ (max-\allowbreak{}width:\ 600px)} (línea 45), y de la regla
  \texttt{.\allowbreak{}grid-\allowbreak{}mascotas} de
  \texttt{frontend/\allowbreak{}src/\allowbreak{}app/\allowbreak{}features/\allowbreak{}mascotas.\allowbreak{}component.\allowbreak{}ts}.
\item
  \textbf{Trazabilidad:} \texttt{frontend/\allowbreak{}src/\allowbreak{}index.\allowbreak{}html} (meta
  viewport), \texttt{frontend/\allowbreak{}src/\allowbreak{}styles.\allowbreak{}css} y
  \texttt{frontend/\allowbreak{}src/\allowbreak{}app/\allowbreak{}features/\allowbreak{}mascotas.\allowbreak{}component.\allowbreak{}ts}
  (\texttt{.\allowbreak{}grid-\allowbreak{}mascotas} con \texttt{@media\ (max-\allowbreak{}width:\ 600px)});
  prueba
  \texttt{Backend/\allowbreak{}src/\allowbreak{}test/\allowbreak{}java/\allowbreak{}com/\allowbreak{}biopet/\allowbreak{}Frontend\allowbreak{}Responsiveness\allowbreak{}Test.\allowbreak{}java};
  configuración de medición \texttt{lighthouserc.\allowbreak{}js} y
  \texttt{lighthouserc.\allowbreak{}desktop.\allowbreak{}js}, guion
  \texttt{scripts/\allowbreak{}run-\allowbreak{}lighthouse.\allowbreak{}sh}; evidencias
  \texttt{docs/\allowbreak{}mediciones/\allowbreak{}lighthouse/} con sus sumas
  \texttt{SHA256\allowbreak{}SUMS.\allowbreak{}txt}.
\item
  \textbf{Estado:} verificado
\end{itemize}

\subsubsection{REQ-NF-006 --- Compatibilidad de
navegadores}\label{req-nf-006-compatibilidad-de-navegadores}

\begin{itemize}
\tightlist
\item
  \textbf{Tipo:} Compatibilidad
\item
  \textbf{Prioridad:} Should
\item
  \textbf{Enunciado:} La aplicación deberá ejecutar correctamente los
  flujos críticos (inicio de sesión, listado y gestión de mascotas,
  gestión de vacunas) en las versiones estables actuales de Google
  Chrome, Mozilla Firefox y Microsoft Edge.
\item
  \textbf{Rationale:} heredado de RNF-05/RNF-WEB-05 de la Entrega 1A.
\item
  \textbf{Criterios de aceptación:}

  \begin{enumerate}
  \def\labelenumi{\arabic{enumi}.}
  \tightlist
  \item
    En cada uno de los tres navegadores, el inicio de sesión completa
    correctamente y la sesión persiste mediante las cookies
    \texttt{Http\allowbreak{}Only}.
  \item
    En cada uno de los tres navegadores, el listado, la creación, la
    edición y la baja de una mascota se completan sin errores de
    consola.
  \item
    En cada uno de los tres navegadores, la pantalla de vacunas carga y
    opera sin errores de consola.
  \item
    La ejecución queda archivada con navegador, versión, fecha y
    resultado por flujo, de modo que un tercero pueda repetirla y
    comparar.
  \end{enumerate}
\item
  \textbf{Método de verificación:} previsto --- ejecución manual
  exploratoria de los tres flujos críticos en Chrome, Firefox y Edge,
  con registro en \texttt{docs/\allowbreak{}mediciones/} del navegador, la versión,
  la fecha y el resultado por flujo (criterios 1 a 4). \textbf{No
  ejecutada todavía.}
\item
  \textbf{Trazabilidad:} aplicación Angular en \texttt{frontend/} (build
  de producción reproducible con
  \texttt{npm\ ci\ \&\&\ npm\ run\ build}, objetivo \texttt{frontend}
  del \texttt{Makefile}); flujos críticos en
  \texttt{frontend/\allowbreak{}src/\allowbreak{}app/\allowbreak{}features/\allowbreak{}login.\allowbreak{}component.\allowbreak{}ts},
  \texttt{mascotas.\allowbreak{}component.\allowbreak{}ts} y \texttt{vacunas.\allowbreak{}component.\allowbreak{}ts}; el
  proyecto no declara un archivo \texttt{browserslist} propio, por lo
  que aplica el objetivo por defecto de Angular 17.
\item
  \textbf{Estado:} pendiente
\item
  \textbf{Observaciones:} \textbf{ninguno} de los cuatro criterios de
  aceptación tiene artefacto: no existen capturas, ni matriz de
  resultados por navegador, ni corrida automatizada, ni registro de
  versiones probadas. El repositorio tampoco contiene ninguna decisión
  ni configuración específica de compatibilidad entre navegadores: no
  hay archivo \texttt{browserslist} propio en \texttt{frontend/}, por lo
  que aplica el objetivo por defecto de Angular 17. Que la aplicación
  compile y se despliegue no constituye implementación de \textbf{este}
  requisito, cuyo objeto es la ejecución comprobada de los flujos
  críticos en Chrome, Firefox y Edge. Se conserva el estado
  \texttt{pendiente} que ya declaraba
  \texttt{docs/\allowbreak{}trazabilidad/\allowbreak{}matriz.\allowbreak{}csv} antes de esta revisión; una
  versión intermedia de esta misma revisión lo había subido a
  \texttt{implementado}, lo cual no estaba justificado y se corrige
  aquí. Sigue como limitación declarada en la sección 7.3.
\end{itemize}

\subsubsection{REQ-NF-007 --- Verificación de estado y medición de
disponibilidad}\label{req-nf-007-verificaciuxf3n-de-estado-y-mediciuxf3n-de-disponibilidad}

\begin{itemize}
\tightlist
\item
  \textbf{Tipo:} Disponibilidad
\item
  \textbf{Prioridad:} Must
\item
  \textbf{Enunciado:} El sistema deberá exponer un endpoint de
  verificación de estado en \texttt{/\allowbreak{}actuator/\allowbreak{}health}, accesible sin
  autenticación, que responda \texttt{200\ OK} con
  \texttt{status:\ "UP"} cuando la aplicación esté operativa; y la
  disponibilidad del sistema durante una ventana de evaluación declarada
  deberá medirse como el porcentaje de sondeos exitosos a ese endpoint,
  con un objetivo de al menos 99 \% y registro fechado de cada sondeo.
\item
  \textbf{Rationale:} heredado de RNF-06 de la Entrega 1A. El enunciado
  anterior ---``el sistema deberá estar operativo durante las semanas de
  evaluación establecidas por la asignatura''--- \textbf{no era
  objetivamente verificable}: no define qué se mide, sobre qué ventana,
  con qué umbral ni con qué artefacto. En esta revisión se reformula en
  dos partes medibles (endpoint observable y porcentaje de sondeos
  exitosos), y el compromiso académico de mantener el sistema operativo
  durante las semanas de evaluación se traslada a la sección 2.5 como
  \textbf{restricción operacional}, donde no se le exige un estado de
  cumplimiento. No se declara ninguna disponibilidad histórica, porque
  nunca se midió.
\item
  \textbf{Criterios de aceptación:}

  \begin{enumerate}
  \def\labelenumi{\arabic{enumi}.}
  \tightlist
  \item
    \texttt{GET\ /\allowbreak{}actuator/\allowbreak{}health} responde \texttt{200\ OK} sin
    necesidad de autenticación.
  \item
    El cuerpo de la respuesta es JSON con
    \texttt{Content-\allowbreak{}Type:\ application/\allowbreak{}json} (o el tipo específico de
    Actuator) e incluye el campo \texttt{status} con valor
    \texttt{"UP"}.
  \item
    La respuesta del endpoint de salud incluye las mismas cabeceras de
    seguridad que el resto de respuestas públicas (REQ-NF-002).
  \item
    Para una ventana de evaluación declarada (fecha de inicio y fecha de
    fin), existe un registro fechado de sondeos periódicos a
    \texttt{/\allowbreak{}actuator/\allowbreak{}health} y el porcentaje de sondeos con respuesta
    \texttt{200\ OK} y \texttt{status:\ "UP"} es ≥ 99 \%.
  \end{enumerate}
\item
  \textbf{Método de verificación:} pruebas automatizadas de
  \texttt{Health\allowbreak{}Check\allowbreak{}Test} (5 pruebas):
  \texttt{health\allowbreak{}Endpoint\allowbreak{}Responde\allowbreak{}Up},
  \texttt{health\allowbreak{}Endpoint\allowbreak{}Respuesta\allowbreak{}Valida\allowbreak{}Sin\allowbreak{}Autenticacion} (criterio 1),
  \texttt{health\allowbreak{}Endpoint\allowbreak{}Estructura\allowbreak{}Valida},
  \texttt{health\allowbreak{}Endpoint\allowbreak{}Content\allowbreak{}Type\allowbreak{}Correcto} (criterio 2),
  \texttt{health\allowbreak{}Endpoint\allowbreak{}Incluye\allowbreak{}Cabeceras\allowbreak{}Seguridad} (criterio 3).
  Criterio 4: \textbf{verificación prevista} --- registro de sondeos
  periódicos (monitoreo externo o guion programado) archivado en
  \texttt{docs/\allowbreak{}mediciones/}, con fecha de inicio y fin de la ventana y
  el cómputo del porcentaje. \textbf{No existe hoy.}
\item
  \textbf{Trazabilidad:} configuración
  \texttt{management.\allowbreak{}endpoints.\allowbreak{}web.\allowbreak{}exposure.\allowbreak{}include:\ health,\allowbreak{}info,\allowbreak{}metrics}
  en \texttt{Backend/\allowbreak{}src/\allowbreak{}main/\allowbreak{}resources/\allowbreak{}application.\allowbreak{}yml}; lista
  \texttt{permit\allowbreak{}All} de \texttt{Security\allowbreak{}Config} (entrada
  \texttt{/\allowbreak{}actuator/\allowbreak{}health}); prueba
  \texttt{Backend/\allowbreak{}src/\allowbreak{}test/\allowbreak{}java/\allowbreak{}com/\allowbreak{}biopet/\allowbreak{}Health\allowbreak{}Check\allowbreak{}Test.\allowbreak{}java};
  healthcheck del despliegue definido en \texttt{render.\allowbreak{}yaml} y
  documentado en \texttt{docs/\allowbreak{}despliegue/\allowbreak{}DEPLOYMENT.\allowbreak{}md} y
  \texttt{docs/\allowbreak{}despliegue/\allowbreak{}RUNBOOK.\allowbreak{}md}; objetivo \texttt{up} del
  \texttt{Makefile} y \texttt{docker-\allowbreak{}compose.\allowbreak{}yml}; restricción
  operacional correspondiente en la sección 2.5.
\item
  \textbf{Estado:} implementado
\item
  \textbf{Observaciones:} el requisito \textbf{no está cumplido}.
  \textbf{Qué cambió respecto del enunciado anterior y por qué.} El
  enunciado que traía este requisito era \emph{``El sistema deberá estar
  operativo durante las semanas de evaluación establecidas por la
  asignatura''}, con estado \emph{``pendiente de confirmación
  explícita''} en el SRS e \emph{``implementado''} en
  \texttt{docs/\allowbreak{}trazabilidad/\allowbreak{}matriz.\allowbreak{}csv}. Ese enunciado no es
  verificable: no dice qué se mide, sobre qué ventana, con qué umbral ni
  con qué artefacto, de modo que ningún resultado podría contradecirlo.
  La reformulación \textbf{no rebaja la intención original}: conserva la
  disponibilidad como obligación medible en el criterio 4 ---que es
  precisamente la parte que \textbf{no} está cumplida--- y añade delante
  la parte observable del mecanismo (criterios 1 a 3). El compromiso
  académico de mantener el sistema operativo durante las semanas de
  evaluación se conserva íntegro en la sección 2.5 como restricción
  operacional, donde no se le atribuye un estado de cumplimiento que
  nadie ha medido. \textbf{Evidencia por criterio.} Criterios 1 a 3:
  verificados por \texttt{Health\allowbreak{}Check\allowbreak{}Test} (5 pruebas) más la
  configuración
  \texttt{management.\allowbreak{}endpoints.\allowbreak{}web.\allowbreak{}exposure.\allowbreak{}include:\ health,\allowbreak{}info,\allowbreak{}metrics}
  y la entrada \texttt{/\allowbreak{}actuator/\allowbreak{}health} en la lista \texttt{permit\allowbreak{}All}
  de \texttt{Security\allowbreak{}Config}, y el
  \texttt{health\allowbreak{}Check\allowbreak{}Path:\ /\allowbreak{}actuator/\allowbreak{}health} de \texttt{render.\allowbreak{}yaml}.
  Criterio 4: \textbf{sin ningún artefacto} --- no existe monitoreo
  continuo, ni log de uptime, ni captura fechada de sondeos sostenidos
  durante una ventana de evaluación. Por eso el requisito queda
  \texttt{implementado} y no \texttt{verificado}, y por eso no se afirma
  ninguna cifra de disponibilidad histórica. Es uno de los requisitos
  \texttt{Must} que no alcanzan \texttt{verificado} (ver sección 3.4.2).
\end{itemize}

\subsubsection{REQ-NF-008 --- Cifrado de
contraseñas}\label{req-nf-008-cifrado-de-contraseuxf1as}

\begin{itemize}
\tightlist
\item
  \textbf{Tipo:} Seguridad
\item
  \textbf{Prioridad:} Must
\item
  \textbf{Enunciado:} Las contraseñas de los usuarios deberán
  almacenarse siempre cifradas mediante el algoritmo de hash adaptativo
  BCrypt con factor de coste 12, nunca en texto plano, y no deberán
  devolverse en ninguna respuesta de la API.
\item
  \textbf{Rationale:} heredado de RNF-07 de la Entrega 1A; reforzado por
  el control OWASP A02 del bloque C.2 de la Guía.
\item
  \textbf{Criterios de aceptación:}

  \begin{enumerate}
  \def\labelenumi{\arabic{enumi}.}
  \tightlist
  \item
    El \texttt{Password\allowbreak{}Encoder} configurado es
    \texttt{BCrypt\allowbreak{}Password\allowbreak{}Encoder} con factor de coste 12.
  \item
    Toda contraseña persistida se almacena como hash BCrypt (prefijo
    \texttt{\$2b\$12\$}), incluida la del usuario administrador inicial.
  \item
    Un inicio de sesión con la contraseña correcta tiene éxito y con una
    incorrecta falla, lo que demuestra que el hash almacenado se
    verifica y no se compara texto plano.
  \item
    Ningún DTO de respuesta de la API expone la contraseña ni su hash.
  \end{enumerate}
\item
  \textbf{Método de verificación:} inspección de
  \texttt{Security\allowbreak{}Config.\allowbreak{}password\allowbreak{}Encoder()}
  (\texttt{new\ BCrypt\allowbreak{}Password\allowbreak{}Encoder(12)}, criterio 1); inspección de
  \texttt{db/\allowbreak{}seed.\allowbreak{}sql}, donde el hash del usuario admin tiene prefijo
  \texttt{\$2b\$12\$}, y de \texttt{Auth\allowbreak{}Service.\allowbreak{}registrar} /
  \texttt{Usuario\allowbreak{}Service.\allowbreak{}crear}, que invocan
  \texttt{password\allowbreak{}Encoder.\allowbreak{}encode(...)} (criterio 2); pruebas
  automatizadas \texttt{Auth\allowbreak{}Controller\allowbreak{}Test.\allowbreak{}login\allowbreak{}Exitoso} y
  \texttt{login\allowbreak{}Clave\allowbreak{}Incorrecta} (criterio 3); inspección de los records
  \texttt{Usuario\allowbreak{}Response}, \texttt{Auth\allowbreak{}Response} y
  \texttt{Auth\allowbreak{}Session\allowbreak{}Response}, ninguno de los cuales declara un campo
  de contraseña (criterio 4).
\item
  \textbf{Trazabilidad:} \texttt{com.\allowbreak{}biopet.\allowbreak{}config.\allowbreak{}Security\allowbreak{}Config} (bean
  \texttt{password\allowbreak{}Encoder}), \texttt{com.\allowbreak{}biopet.\allowbreak{}service.\allowbreak{}Auth\allowbreak{}Service},
  \texttt{com.\allowbreak{}biopet.\allowbreak{}service.\allowbreak{}Usuario\allowbreak{}Service}, entidad
  \texttt{com.\allowbreak{}biopet.\allowbreak{}entity.\allowbreak{}Usuario} (campo \texttt{password\allowbreak{}Hash}),
  \texttt{db/\allowbreak{}seed.\allowbreak{}sql},
  \texttt{Backend/\allowbreak{}src/\allowbreak{}main/\allowbreak{}resources/\allowbreak{}db/\allowbreak{}migration/\allowbreak{}V1\_\allowbreak{}\_\allowbreak{}schema\_\allowbreak{}inicial.\allowbreak{}sql};
  evidencia \texttt{docs/\allowbreak{}mediciones/\allowbreak{}sec/\allowbreak{}A02-\allowbreak{}cryptography-\allowbreak{}tls.\allowbreak{}md};
  política de contraseñas asociada en REQ-NF-016.
\item
  \textbf{Estado:} verificado
\end{itemize}

\subsubsection{REQ-NF-009 --- Registro de eventos de autenticación
(OWASP
A09)}\label{req-nf-009-registro-de-eventos-de-autenticaciuxf3n-owasp-a09}

\begin{itemize}
\tightlist
\item
  \textbf{Tipo:} Seguridad
\item
  \textbf{Prioridad:} Must
\item
  \textbf{Enunciado:} El sistema deberá registrar cada evento de
  autenticación (login exitoso, login fallido, bloqueo por límite de
  intentos, refresh exitoso, refresh fallido, logout y uso de token
  revocado) con dirección IP, marca de tiempo y sujeto, sin exceder 200
  caracteres por campo registrado y sin incluir credenciales ni tokens.
\item
  \textbf{Rationale:} requisito derivado del control OWASP A09, exigido
  por el bloque C.2 de la Guía; ausente en el SRS original de la Entrega
  1A.
\item
  \textbf{Criterios de aceptación:}

  \begin{enumerate}
  \def\labelenumi{\arabic{enumi}.}
  \tightlist
  \item
    Cada uno de los siete eventos enumerados produce una línea de log
    con prefijo \texttt{AUTH\_\allowbreak{}AUDIT}, el tipo de evento, la IP, la marca
    de tiempo UTC y el sujeto.
  \item
    Los eventos de fallo o bloqueo se registran con nivel \texttt{WARN}
    y los de éxito con nivel \texttt{INFO}.
  \item
    Un valor nulo de IP o de sujeto se normaliza a \texttt{unknown},
    nunca produce una línea incompleta.
  \item
    Los caracteres de control se eliminan del sujeto antes de
    registrarlo, de modo que no se puedan inyectar saltos de línea ni
    líneas falsas (\emph{log forging}).
  \item
    Ninguna línea contiene contraseñas, hashes ni tokens.
  \item
    Ningún campo registrado supera 200 caracteres.
  \end{enumerate}
\item
  \textbf{Método de verificación:} pruebas automatizadas de
  \texttt{Authentication\allowbreak{}Audit\allowbreak{}Service\allowbreak{}Test} (10 pruebas):
  \texttt{login\allowbreak{}Exitoso\allowbreak{}Registra\allowbreak{}Evento\allowbreak{}Estructurado},
  \texttt{login\allowbreak{}Fallido\allowbreak{}Registra\allowbreak{}Evento\allowbreak{}Warn},
  \texttt{login\allowbreak{}Bloqueado\allowbreak{}Registra\allowbreak{}Evento\allowbreak{}Warn},
  \texttt{refresh\allowbreak{}Exitoso\allowbreak{}Registra\allowbreak{}Evento\allowbreak{}Info},
  \texttt{refresh\allowbreak{}Fallido\allowbreak{}Registra\allowbreak{}Evento\allowbreak{}Warn},
  \texttt{logout\allowbreak{}Exitoso\allowbreak{}Registra\allowbreak{}Evento\allowbreak{}Info},
  \texttt{token\allowbreak{}Revocado\allowbreak{}Registra\allowbreak{}Evento\allowbreak{}Warn} (criterios 1 y 2);
  \texttt{valores\allowbreak{}Nulos\allowbreak{}Se\allowbreak{}Normalizan\allowbreak{}Como\allowbreak{}Unknown} (criterio 3);
  \texttt{elimina\allowbreak{}Caracteres\allowbreak{}De\allowbreak{}Control\allowbreak{}Para\allowbreak{}Evitar\allowbreak{}Log\allowbreak{}Forging} (criterio 4);
  \texttt{no\allowbreak{}Registra\allowbreak{}Datos\allowbreak{}Sensibles} (criterio 5); inspección de
  \texttt{Authentication\allowbreak{}Audit\allowbreak{}Service}, que declara
  \texttt{LONGITUD\_\allowbreak{}MAXIMA\ =\ 2\allowbreak{}00} y trunca cada campo con
  \texttt{substring(0,\allowbreak{}\ LONGITUD\_\allowbreak{}MAXIMA)} (criterio 6). Evidencia de
  log real en \texttt{docs/\allowbreak{}mediciones/\allowbreak{}sec/\allowbreak{}A09-\allowbreak{}logging.\allowbreak{}md} y
  \texttt{docs/\allowbreak{}mediciones/\allowbreak{}sec/\allowbreak{}raw/\allowbreak{}A09-\allowbreak{}audit-\allowbreak{}logs.\allowbreak{}txt}, con líneas
  \texttt{AUTH\_\allowbreak{}AUDIT\ event=L\allowbreak{}OGIN\_\allowbreak{}SUCCESS}, \texttt{LOGIN\_\allowbreak{}FAILURE} y
  \texttt{LOGIN\_\allowbreak{}RATE\_\allowbreak{}LIMITED} con IP, timestamp UTC y sujeto reales.
\item
  \textbf{Trazabilidad:}
  \texttt{com.\allowbreak{}biopet.\allowbreak{}security.\allowbreak{}Authentication\allowbreak{}Audit\allowbreak{}Service} (invocado
  desde \texttt{Auth\allowbreak{}Service} y \texttt{Jwt\allowbreak{}Authentication\allowbreak{}Filter}); prueba
  \texttt{Backend/\allowbreak{}src/\allowbreak{}test/\allowbreak{}java/\allowbreak{}com/\allowbreak{}biopet/\allowbreak{}security/\allowbreak{}Authentication\allowbreak{}Audit\allowbreak{}Service\allowbreak{}Test.\allowbreak{}java};
  evidencia \texttt{docs/\allowbreak{}mediciones/\allowbreak{}sec/\allowbreak{}A09-\allowbreak{}logging.\allowbreak{}md} con
  \texttt{docs/\allowbreak{}mediciones/\allowbreak{}sec/\allowbreak{}raw/\allowbreak{}A09-\allowbreak{}audit-\allowbreak{}logs.\allowbreak{}txt}; control OWASP
  A09; requisito funcional hermano REQ-F-020.
\item
  \textbf{Estado:} verificado
\item
  \textbf{Observaciones:} no existe integración con un SIEM
  centralizado: los eventos quedan en el log de la aplicación. La
  auditoría de operaciones CRUD de negocio (distinta de la de
  autenticación) corresponde a REQ-F-020, que está \texttt{implementado}
  precisamente por esa carencia.
\end{itemize}

\subsubsection{REQ-NF-010 --- Limitación de intentos de login (OWASP
A07)}\label{req-nf-010-limitaciuxf3n-de-intentos-de-login-owasp-a07}

\begin{itemize}
\tightlist
\item
  \textbf{Tipo:} Seguridad
\item
  \textbf{Prioridad:} Must
\item
  \textbf{Enunciado:} El sistema deberá bloquear temporalmente los
  intentos de inicio de sesión provenientes de una misma dirección IP
  tras 6 intentos fallidos consecutivos dentro de una ventana de 15
  minutos, respondiendo \texttt{429\ Too\ Many\ Reques\allowbreak{}ts} con cabecera
  \texttt{Retry-\allowbreak{}After} durante el bloqueo, con los tres parámetros
  configurables mediante propiedades externas.
\item
  \textbf{Rationale:} requisito derivado del control OWASP A07, exigido
  por el bloque C.2 de la Guía; ausente en el SRS original de la Entrega
  1A.
\item
  \textbf{Criterios de aceptación:}

  \begin{enumerate}
  \def\labelenumi{\arabic{enumi}.}
  \tightlist
  \item
    Los cinco primeros intentos fallidos desde una misma IP responden
    \texttt{401}.
  \item
    El sexto intento fallido responde \texttt{429} con
    \texttt{Problem\allowbreak{}Detail} de tipo
    \texttt{urn:biopet:error:\allowbreak{}rate-\allowbreak{}limited} y cabecera
    \texttt{Retry-\allowbreak{}After:\ 900}.
  \item
    Durante el bloqueo, incluso un intento con credenciales
    \textbf{correctas} desde esa IP es rechazado con \texttt{429}.
  \item
    Los contadores son independientes por IP: los fallos desde IP
    distintas no se acumulan entre sí.
  \item
    Un inicio de sesión exitoso reinicia el contador de fallos de esa
    IP.
  \item
    Al expirar la ventana sin alcanzar el umbral, el contador se
    reinicia; al expirar el bloqueo, se admite un nuevo intento.
  \item
    Los parámetros de umbral, ventana y duración de bloqueo se leen de
    propiedades externas
    (\texttt{security.\allowbreak{}rate-\allowbreak{}limit.\allowbreak{}login.\allowbreak{}max-\allowbreak{}attempts}, \texttt{.\allowbreak{}window},
    \texttt{.\allowbreak{}block-\allowbreak{}duration}), no de literales en el código.
  \end{enumerate}
\item
  \textbf{Método de verificación:} pruebas automatizadas de
  \texttt{Login\allowbreak{}Rate\allowbreak{}Limiter\allowbreak{}Service\allowbreak{}Test} (7 pruebas):
  \texttt{cinco\allowbreak{}Fallos\allowbreak{}No\allowbreak{}Bloquean} (criterio 1),
  \texttt{sexto\allowbreak{}Fallo\allowbreak{}Bloquea\allowbreak{}YLanza\allowbreak{}Excepcion} (criterio 2),
  \texttt{ip\allowbreak{}Bloqueada\allowbreak{}Sigue\allowbreak{}Rechazada\allowbreak{}Con\allowbreak{}Tiempo\allowbreak{}Restante} (criterio 3),
  \texttt{ips\allowbreak{}Diferentes\allowbreak{}Mantienen\allowbreak{}Contadores\allowbreak{}Separados} (criterio 4),
  \texttt{reiniciar\allowbreak{}Elimina\allowbreak{}Fallos\allowbreak{}YBloqueo} (criterio 5),
  \texttt{ventana\allowbreak{}Expirada\allowbreak{}Reinicia\allowbreak{}El\allowbreak{}Contador} y
  \texttt{bloqueo\allowbreak{}Expirado\allowbreak{}Permite\allowbreak{}Nuevo\allowbreak{}Intento} (criterio 6); y de
  \texttt{Auth\allowbreak{}Controller\allowbreak{}Test}:
  \texttt{quinto\allowbreak{}Intento\allowbreak{}Fallido\allowbreak{}Sigue\allowbreak{}Respondiendo401},
  \texttt{sexto\allowbreak{}Intento\allowbreak{}Fallido\allowbreak{}Devuelve429\allowbreak{}Problem\allowbreak{}Detail},
  \texttt{ip\allowbreak{}Bloqueada\allowbreak{}Rechaza\allowbreak{}Incluso\allowbreak{}Credenciales\allowbreak{}Correctas},
  \texttt{intentos\allowbreak{}Desde\allowbreak{}Ips\allowbreak{}Distintas\allowbreak{}No\allowbreak{}Se\allowbreak{}Acumulan},
  \texttt{login\allowbreak{}Exitoso\allowbreak{}Reinicia\allowbreak{}Contador\allowbreak{}De\allowbreak{}La\allowbreak{}Ip}. Criterio 7: inspección de
  las anotaciones \texttt{@Value} del constructor de
  \texttt{Login\allowbreak{}Rate\allowbreak{}Limiter\allowbreak{}Service}. Evidencia HTTP real (\texttt{curl})
  en \texttt{docs/\allowbreak{}mediciones/\allowbreak{}sec/\allowbreak{}A07-\allowbreak{}authentication.\allowbreak{}md} con
  \texttt{docs/\allowbreak{}mediciones/\allowbreak{}sec/\allowbreak{}raw/\allowbreak{}A07-\allowbreak{}auth-\allowbreak{}rate-\allowbreak{}limit.\allowbreak{}txt}: cinco
  \texttt{401}, un \texttt{429} con
  \texttt{type=urn:biopet:\allowbreak{}error:rate-\allowbreak{}limited} y
  \texttt{Retry-\allowbreak{}After:\ 900}.
\item
  \textbf{Trazabilidad:}
  \texttt{com.\allowbreak{}biopet.\allowbreak{}security.\allowbreak{}Login\allowbreak{}Rate\allowbreak{}Limiter\allowbreak{}Service} (invocado desde
  \texttt{Auth\allowbreak{}Service.\allowbreak{}login}), excepción
  \texttt{Rate\allowbreak{}Limit\allowbreak{}Excedido\allowbreak{}Exception} →
  \texttt{Global\allowbreak{}Exception\allowbreak{}Handler.\allowbreak{}demasiados\allowbreak{}Intentos} →
  \texttt{Problem\allowbreak{}Type.\allowbreak{}RATE\_\allowbreak{}LIMITED}; pruebas
  \texttt{Login\allowbreak{}Rate\allowbreak{}Limiter\allowbreak{}Service\allowbreak{}Test} y \texttt{Auth\allowbreak{}Controller\allowbreak{}Test};
  evidencia \texttt{docs/\allowbreak{}mediciones/\allowbreak{}sec/\allowbreak{}A07-\allowbreak{}authentication.\allowbreak{}md}; control
  OWASP A07.
\item
  \textbf{Estado:} verificado
\item
  \textbf{Observaciones:} el contador de intentos se mantiene en memoria
  del proceso (\texttt{Concurrent\allowbreak{}Hash\allowbreak{}Map} en
  \texttt{Login\allowbreak{}Rate\allowbreak{}Limiter\allowbreak{}Service}), no en Redis. En un despliegue con
  más de una instancia el límite se aplicaría por instancia, no de forma
  global; el despliegue actual es de instancia única. Como efecto
  secundario positivo, este mecanismo \textbf{no} depende de Redis y por
  tanto no se ve afectado por el escenario descrito en REQ-NF-014.
\end{itemize}

\subsubsection{REQ-NF-011 --- Arquitectura en
capas}\label{req-nf-011-arquitectura-en-capas}

\begin{itemize}
\tightlist
\item
  \textbf{Tipo:} Mantenibilidad
\item
  \textbf{Prioridad:} Must
\item
  \textbf{Enunciado:} La arquitectura del backend deberá organizarse en
  capas separadas de presentación (\texttt{controller}), lógica de
  negocio (\texttt{service}) y acceso a datos
  (\texttt{repository}/\texttt{entity}), sin que la capa de presentación
  acceda directamente a los repositorios.
\item
  \textbf{Rationale:} heredado de RNF-08 de la Entrega 1A; decisión de
  acceso a datos registrada en
  \texttt{docs/\allowbreak{}adr/\allowbreak{}ADR-\allowbreak{}007-\allowbreak{}acceso-\allowbreak{}datos.\allowbreak{}md}.
\item
  \textbf{Criterios de aceptación:}

  \begin{enumerate}
  \def\labelenumi{\arabic{enumi}.}
  \tightlist
  \item
    Existen los paquetes \texttt{com.\allowbreak{}biopet.\allowbreak{}controller},
    \texttt{com.\allowbreak{}biopet.\allowbreak{}service}, \texttt{com.\allowbreak{}biopet.\allowbreak{}repository},
    \texttt{com.\allowbreak{}biopet.\allowbreak{}entity} y \texttt{com.\allowbreak{}biopet.\allowbreak{}dto}, cada uno con
    las clases que le corresponden.
  \item
    Ninguna clase de \texttt{com.\allowbreak{}biopet.\allowbreak{}controller} declara una
    dependencia de un tipo del paquete \texttt{com.\allowbreak{}biopet.\allowbreak{}repository}:
    los seis controladores inyectan únicamente servicios
    (\texttt{Auth\allowbreak{}Service}, \texttt{Usuario\allowbreak{}Service},
    \texttt{Mascota\allowbreak{}Service}, \texttt{Cita\allowbreak{}Service},
    \texttt{Consulta\allowbreak{}Service}, \texttt{Vacuna\allowbreak{}Service},
    \texttt{External\allowbreak{}Api\allowbreak{}Service}).
  \item
    Las entidades JPA no se exponen en las firmas públicas de los
    controladores: las respuestas se construyen con los records de
    \texttt{com.\allowbreak{}biopet.\allowbreak{}dto}.
  \item
    La estructura de capas está reflejada en el diagrama de componentes
    C4 de nivel 3 del repositorio.
  \end{enumerate}
\item
  \textbf{Método de verificación:} inspección de la estructura de
  paquetes de \texttt{Backend/\allowbreak{}src/\allowbreak{}main/\allowbreak{}java/\allowbreak{}com/\allowbreak{}biopet/} (criterio 1);
  inspección de los constructores de los seis controladores de
  \texttt{com.\allowbreak{}biopet.\allowbreak{}controller} (criterio 2); inspección de las firmas
  de los métodos anotados con
  \texttt{@Get\allowbreak{}Mapping}/\texttt{@Post\allowbreak{}Mapping}/\texttt{@Put\allowbreak{}Mapping}/\texttt{@Delete\allowbreak{}Mapping},
  que devuelven \texttt{Mascota\allowbreak{}Response}, \texttt{Cita\allowbreak{}Response},
  \texttt{Consulta\allowbreak{}Response}, \texttt{Vacuna\allowbreak{}Response},
  \texttt{Usuario\allowbreak{}Response}, \texttt{External\allowbreak{}Api\allowbreak{}Response},
  \texttt{Resumen\allowbreak{}Especie\allowbreak{}Response} o \texttt{Auth\allowbreak{}Session\allowbreak{}Response}
  (criterio 3); diagrama
  \texttt{docs/\allowbreak{}diagrams/\allowbreak{}c4-\allowbreak{}componentes-\allowbreak{}backend/\allowbreak{}C4-\allowbreak{}L3-\allowbreak{}backend.\allowbreak{}md} y su
  render \texttt{c4-\allowbreak{}componentes-\allowbreak{}backend.\allowbreak{}png} (criterio 4).
\item
  \textbf{Trazabilidad:} paquetes
  \texttt{com.\allowbreak{}biopet.\allowbreak{}\{controller,\allowbreak{}service,\allowbreak{}repository,\allowbreak{}entity,\allowbreak{}dto\}} en
  \texttt{Backend/\allowbreak{}src/\allowbreak{}main/\allowbreak{}java/\allowbreak{}com/\allowbreak{}biopet/}; diagramas C4 en
  \texttt{docs/\allowbreak{}diagrams/\allowbreak{}c4-\allowbreak{}componentes-\allowbreak{}backend/} y
  \texttt{docs/\allowbreak{}diagrams/\allowbreak{}c4-\allowbreak{}contenedores/}; decisiones
  \texttt{docs/\allowbreak{}adr/\allowbreak{}ADR-\allowbreak{}002-\allowbreak{}pila-\allowbreak{}tecnologica.\allowbreak{}md} y
  \texttt{docs/\allowbreak{}adr/\allowbreak{}ADR-\allowbreak{}007-\allowbreak{}acceso-\allowbreak{}datos.\allowbreak{}md}; diagrama de clases
  \texttt{docs/\allowbreak{}diagrams/\allowbreak{}diagrama-\allowbreak{}clases/\allowbreak{}diagrama-\allowbreak{}clases.\allowbreak{}png}.
\item
  \textbf{Estado:} verificado
\item
  \textbf{Observaciones:} no existe un control automatizado de
  dependencias entre capas (por ejemplo, una regla de ArchUnit) que
  impida introducir en el futuro una dependencia de controlador a
  repositorio; hoy la separación se comprueba por inspección. Registrado
  como mejora, no como defecto.
\end{itemize}

\subsubsection{REQ-NF-012 --- Caché del listado de mascotas: TTL
configurable y medición de
aciertos}\label{req-nf-012-cachuxe9-del-listado-de-mascotas-ttl-configurable-y-mediciuxf3n-de-aciertos}

\begin{quote}
\textbf{Requisito con dos obligaciones.} Agrupa \textbf{(A)} que el
listado se almacene en caché con un TTL configurable por variable de
entorno y \textbf{(B)} que la tasa de aciertos de esa caché se mida y se
reporte empíricamente. La primera está cumplida y la segunda no, por lo
que en revisiones anteriores el requisito arrastraba el valor inválido
``verificado parcialmente''. Conserva su identificador único; las
obligaciones se separan con criterios \texttt{A1…A4} y \texttt{B1…B4}, y
el Estado describe el requisito completo.
\end{quote}

\begin{itemize}
\tightlist
\item
  \textbf{Tipo:} Rendimiento
\item
  \textbf{Prioridad:} Should
\item
  \textbf{Enunciado:} El listado paginado de mascotas deberá
  \textbf{(A)} almacenarse en caché Redis con un tiempo de vida (TTL)
  configurable mediante variable de entorno (\texttt{CACHE\_\allowbreak{}TTL\_\allowbreak{}MS},
  valor por defecto 300 000 ms), sin valores fijos en el código y con
  una entrada de caché distinta por usuario y por parámetros de
  paginación; y \textbf{(B)} su tasa de aciertos (\emph{hit ratio}),
  definida como
  \texttt{keyspace\_\allowbreak{}hits\ /\allowbreak{}\ (keyspace\_\allowbreak{}hits\ +\ ke\allowbreak{}yspace\_\allowbreak{}misses)},
  deberá medirse y reportarse empíricamente sobre una corrida de carga
  declarada, con un objetivo de al menos 80 \% en régimen de caché
  caliente.
\item
  \textbf{Rationale:} heredado de la estrategia de caché definida en
  \texttt{docs/\allowbreak{}adr/\allowbreak{}ADR-\allowbreak{}003-\allowbreak{}jwt-\allowbreak{}redis.\allowbreak{}md}; llevado a requisito no
  funcional explícito por el bloque A.1 de la Guía, que pide evidencia
  empírica de la estrategia de caché y no sólo de su configuración. La
  obligación A es la condición que hace alcanzables los umbrales de
  REQ-NF-001 con caché caliente.
\item
  \textbf{Criterios de aceptación:}

  \begin{enumerate}
  \def\labelenumi{\arabic{enumi}.}
  \tightlist
  \item
    \textbf{(A1)} El TTL de la caché se lee de la variable de entorno
    \texttt{CACHE\_\allowbreak{}TTL\_\allowbreak{}MS}, con valor por defecto 300 000 ms, y no
    aparece como literal en ninguna clase Java.
  \item
    \textbf{(A2)} Bajo carga sobre \texttt{GET\ /\allowbreak{}api/\allowbreak{}mascotas}, existe
    al menos una clave \texttt{mascotas::*} en Redis y el comando
    \texttt{TTL} sobre ella devuelve un valor positivo y coherente con
    el configurado.
  \item
    \textbf{(A3)} La clave de caché incluye el correo del usuario
    autenticado y los parámetros de paginación, de modo que dos dueños
    distintos que piden la misma página no comparten entrada.
  \item
    \textbf{(A4)} Toda operación de escritura sobre mascotas invalida la
    caché completa del espacio \texttt{mascotas}.
  \item
    \textbf{(B1)} Existe un archivo de evidencia con la salida de
    \texttt{redis-\allowbreak{}cli\ INFO\ stats} (o del endpoint
    \texttt{/\allowbreak{}actuator/\allowbreak{}metrics} equivalente) tomada \textbf{antes} y
    \textbf{después} de una corrida de carga declarada sobre
    \texttt{GET\ /\allowbreak{}api/\allowbreak{}mascotas}.
  \item
    \textbf{(B2)} A partir de esos dos puntos se calcula el \emph{hit
    ratio} de la corrida y el valor obtenido queda registrado junto con
    el número de peticiones, la duración y la fecha de la corrida.
  \item
    \textbf{(B3)} El \emph{hit ratio} medido en régimen de caché
    caliente es ≥ 80 \%.
  \item
    \textbf{(B4)} La medición es reproducible: el guion o los comandos
    exactos que la generan están versionados en el repositorio.
  \end{enumerate}
\item
  \textbf{Método de verificación:} obligación A --- inspección de
  \texttt{Backend/\allowbreak{}src/\allowbreak{}main/\allowbreak{}resources/\allowbreak{}application.\allowbreak{}yml}
  (\texttt{spring.\allowbreak{}cache.\allowbreak{}redis.\allowbreak{}time-\allowbreak{}to-\allowbreak{}live:\ \$\{CACHE\_\allowbreak{}TTL\_\allowbreak{}MS:\allowbreak{}300000\}},
  \texttt{cache-\allowbreak{}null-\allowbreak{}values:\ false}) y de \texttt{.\allowbreak{}env.\allowbreak{}example}
  (\texttt{CACHE\_\allowbreak{}TTL\_\allowbreak{}MS=30000\allowbreak{}0}) (A1); evidencia cruda archivada en
  \texttt{docs/\allowbreak{}mediciones/\allowbreak{}redis/}: \texttt{ttl-\allowbreak{}cache-\allowbreak{}mascotas.\allowbreak{}txt} (TTL
  real observado de 195 s sobre la clave de caché),
  \texttt{keys-\allowbreak{}mascotas.\allowbreak{}txt} (\texttt{KEYS\ mascotas::*}),
  \texttt{dbsize-\allowbreak{}durante-\allowbreak{}carga.\allowbreak{}txt} (\texttt{DBSIZE} bajo carga) y
  \texttt{config-\allowbreak{}maxmemory.\allowbreak{}txt} (A2); prueba automatizada
  \texttt{Mascota\allowbreak{}Controller\allowbreak{}Test.\allowbreak{}dos\allowbreak{}Duenos\allowbreak{}Con\allowbreak{}Misma\allowbreak{}Pagina\allowbreak{}No\allowbreak{}Comparen\allowbreak{}Resultados\allowbreak{}De\allowbreak{}Cache}
  e inspección de la anotación \texttt{@Cacheable} de
  \texttt{Mascota\allowbreak{}Service.\allowbreak{}listar}, cuya clave se compone del correo del
  usuario autenticado y de los parámetros de paginación (A3); inspección
  de las tres anotaciones
  \texttt{@Cache\allowbreak{}Evict(value\ =\ "mas\allowbreak{}cotas",\allowbreak{}\ all\allowbreak{}Entries\ =\ true)} sobre
  \texttt{crear}, \texttt{actualizar} y \texttt{eliminar} de
  \texttt{Mascota\allowbreak{}Service} (A4). Obligación B --- \textbf{verificación
  prevista, no ejecutada}: captura de
  \texttt{redis-\allowbreak{}cli\ INFO\ stats\ \textbar{}\ gr\allowbreak{}ep\ keyspace\_} antes y
  después de una corrida k6 equivalente a las de REQ-NF-001, archivada
  en \texttt{docs/\allowbreak{}mediciones/\allowbreak{}redis/} con el cálculo del cociente;
  alternativamente, las métricas de caché de Micrometer expuestas por
  \texttt{/\allowbreak{}actuator/\allowbreak{}metrics} (\texttt{cache.\allowbreak{}gets\{result=hit\}} y
  \texttt{cache.\allowbreak{}gets\{result=miss\}}), vía ya habilitada en
  \texttt{management.\allowbreak{}endpoints.\allowbreak{}web.\allowbreak{}exposure.\allowbreak{}include}.
\item
  \textbf{Trazabilidad:} \texttt{com.\allowbreak{}biopet.\allowbreak{}service.\allowbreak{}Mascota\allowbreak{}Service}
  (anotaciones \texttt{@Cacheable}/\texttt{@Cache\allowbreak{}Evict}),
  \texttt{Backend/\allowbreak{}src/\allowbreak{}main/\allowbreak{}resources/\allowbreak{}application.\allowbreak{}yml}
  (\texttt{spring.\allowbreak{}cache.\allowbreak{}type:\ redis},
  \texttt{spring.\allowbreak{}cache.\allowbreak{}redis.\allowbreak{}time-\allowbreak{}to-\allowbreak{}live},
  \texttt{management.\allowbreak{}endpoints.\allowbreak{}web.\allowbreak{}exposure.\allowbreak{}include}),
  \texttt{.\allowbreak{}env.\allowbreak{}example} (\texttt{CACHE\_\allowbreak{}TTL\_\allowbreak{}MS}), decisión
  \texttt{docs/\allowbreak{}adr/\allowbreak{}ADR-\allowbreak{}003-\allowbreak{}jwt-\allowbreak{}redis.\allowbreak{}md}, evidencias
  \texttt{docs/\allowbreak{}mediciones/\allowbreak{}redis/} (TTL, \texttt{DBSIZE}, \texttt{KEYS},
  \texttt{maxmemory} --- que documentan la \textbf{configuración y la
  existencia} de la clave, no su tasa de aciertos), guiones de carga en
  \texttt{k6/} y reporte \texttt{docs/\allowbreak{}mediciones/\allowbreak{}perf/\allowbreak{}REPORT.\allowbreak{}md}, prueba
  \texttt{Backend/\allowbreak{}src/\allowbreak{}test/\allowbreak{}java/\allowbreak{}com/\allowbreak{}biopet/\allowbreak{}Mascota\allowbreak{}Controller\allowbreak{}Test.\allowbreak{}java};
  requisito de rendimiento dependiente REQ-NF-001; comportamiento ante
  fallo de Redis, REQ-NF-014.
\item
  \textbf{Estado:} implementado
\item
  \textbf{Observaciones:} el requisito \textbf{no está cumplido}. La
  obligación A está realizada y evidenciada en sus cuatro criterios. La
  obligación B \textbf{no tiene ninguna medición}: no existe en el
  repositorio ningún registro de \texttt{keyspace\_\allowbreak{}hits} /
  \texttt{keyspace\_\allowbreak{}misses}, por lo que \texttt{B1} a \texttt{B4}
  carecen de artefacto. Una versión anterior de la matriz de
  trazabilidad afirmaba un ``hit ratio cercano al 100 \% documentado en
  \texttt{docs/\allowbreak{}mediciones/\allowbreak{}perf/\allowbreak{}REPORT.\allowbreak{}md}''; esa afirmación \textbf{no
  es cierta} ---ese reporte no contiene tal medición--- y fue retirada
  en la revisión del 2026-09-07. No se declara aquí ninguna cifra de
  \emph{hit ratio}.
\end{itemize}

\subsubsection{REQ-NF-013 --- Estrategia híbrida de acceso a datos (ORM
+ procedimientos
almacenados)}\label{req-nf-013-estrategia-huxedbrida-de-acceso-a-datos-orm-procedimientos-almacenados}

\begin{itemize}
\tightlist
\item
  \textbf{Tipo:} Mantenibilidad / Seguridad
\item
  \textbf{Prioridad:} Must
\item
  \textbf{Enunciado:} Toda operación de base de datos que no sea un CRUD
  elemental sobre una única tabla (según la definición del bloque A.2.1
  de la Guía) deberá implementarse como función o procedimiento
  almacenado versionado en \texttt{db/\allowbreak{}procs/} e invocado desde un
  repositorio Spring Data mediante el mecanismo formal
  \texttt{@Procedure}, quedando prohibida la concatenación de entrada de
  usuario en cualquier fragmento JPQL, HQL o SQL nativo.
\item
  \textbf{Rationale:} requisito exigido explícitamente por el bloque A.2
  de la Guía. Una versión anterior de este documento lo marcaba
  ``cumplida por alcance actual'' porque entonces no existía ninguna
  operación no elemental; desde v0.9.0-rc existen seis rutinas reales.
\item
  \textbf{Criterios de aceptación:}

  \begin{enumerate}
  \def\labelenumi{\arabic{enumi}.}
  \tightlist
  \item
    Las seis rutinas de \texttt{db/\allowbreak{}procs/} son objetos PostgreSQL
    \texttt{PROCEDURE} versionados en el repositorio y aplicados por
    migraciones Flyway.
  \item
    Las seis se invocan desde Java exclusivamente con
    \texttt{@Procedure} (no con \texttt{@Query(native\allowbreak{}Query\ =\ true)} ni
    JDBC directo desde los servicios).
  \item
    El análisis \texttt{scripts/\allowbreak{}audit-\allowbreak{}sql-\allowbreak{}dynamic.\allowbreak{}sh} sobre los archivos
    de \texttt{db/\allowbreak{}procs/} devuelve 0 hallazgos de SQL dinámico
    construido por concatenación, con código de salida 0.
  \item
    Cada rutina tiene al menos una prueba de integración que la ejercita
    contra un PostgreSQL real.
  \item
    El catálogo documental de rutinas está completo: cada fila es
    trazable a un archivo \texttt{.\allowbreak{}sql} real del repositorio.
  \end{enumerate}
\item
  \textbf{Método de verificación:} inspección de \texttt{db/\allowbreak{}procs/} y de
  las migraciones
  \texttt{Backend/\allowbreak{}src/\allowbreak{}main/\allowbreak{}resources/\allowbreak{}db/\allowbreak{}migration/\allowbreak{}V5\_\allowbreak{}\_\allowbreak{}procedimientos\_\allowbreak{}b\allowbreak{}iopet.\allowbreak{}sql}
  y \texttt{V6\_\allowbreak{}\_\allowbreak{}formalizar\_\allowbreak{}proc\allowbreak{}edimientos\_\allowbreak{}jpa.\allowbreak{}sql} (criterio 1);
  inspección de
  \texttt{com.\allowbreak{}biopet.\allowbreak{}repository.\allowbreak{}Procedimiento\allowbreak{}Biopet\allowbreak{}Repository}, cuyos
  seis métodos están anotados con \texttt{@Procedure} (criterio 2);
  ejecución archivada de \texttt{scripts/\allowbreak{}audit-\allowbreak{}sql-\allowbreak{}dynamic.\allowbreak{}sh} con
  resultado
  \texttt{audit-\allowbreak{}sql-\allowbreak{}dynamic:\ 0\ hallazg\allowbreak{}os\ en\ 7\ archivo(s).\allowbreak{}\ PASA}
  (código de salida 0), evidencia en
  \texttt{docs/\allowbreak{}mediciones/\allowbreak{}sec/\allowbreak{}audit-\allowbreak{}sql-\allowbreak{}dynamic-\allowbreak{}v1.\allowbreak{}0.\allowbreak{}0.\allowbreak{}txt} con su suma
  \texttt{audit-\allowbreak{}sql-\allowbreak{}dynamic-\allowbreak{}v1.\allowbreak{}0.\allowbreak{}0.\allowbreak{}sha256} (criterio 3); pruebas de
  integración \texttt{Procedimientos\allowbreak{}Biopet\allowbreak{}Integration\allowbreak{}Test} (13 pruebas
  sobre las seis rutinas) y \texttt{Resumen\allowbreak{}Especies\allowbreak{}Integration\allowbreak{}Test},
  ambas con Testcontainers sobre PostgreSQL real (criterio 4);
  \texttt{docs/\allowbreak{}basedatos/\allowbreak{}CATALOGO-\allowbreak{}SP.\allowbreak{}md} (criterio 5).
\item
  \textbf{Trazabilidad:} rutinas
  \texttt{db/\allowbreak{}procs/\allowbreak{}fn\_\allowbreak{}resumen\_\allowbreak{}mascota\allowbreak{}s\_\allowbreak{}por\_\allowbreak{}especie.\allowbreak{}sql},
  \texttt{fn\_\allowbreak{}historial\_\allowbreak{}clini\allowbreak{}co\_\allowbreak{}mascota.\allowbreak{}sql},
  \texttt{fn\_\allowbreak{}reporte\_\allowbreak{}dashboa\allowbreak{}rd.\allowbreak{}sql},
  \texttt{fn\_\allowbreak{}siguiente\_\allowbreak{}numer\allowbreak{}o\_\allowbreak{}ficha.\allowbreak{}sql},
  \texttt{sp\_\allowbreak{}actualizar\_\allowbreak{}esta\allowbreak{}do\_\allowbreak{}citas\_\allowbreak{}masivas.\allowbreak{}sql},
  \texttt{sp\_\allowbreak{}registrar\_\allowbreak{}consu\allowbreak{}lta\_\allowbreak{}validada.\allowbreak{}sql} y los permisos
  \texttt{zz\_\allowbreak{}grants\_\allowbreak{}biopet\_\allowbreak{}ap\allowbreak{}p.\allowbreak{}sql}; invocación formal en
  \texttt{com.\allowbreak{}biopet.\allowbreak{}repository.\allowbreak{}Procedimiento\allowbreak{}Biopet\allowbreak{}Repository} con los
  \texttt{@Named\allowbreak{}Stored\allowbreak{}Procedure\allowbreak{}Query} declarados en
  \texttt{com.\allowbreak{}biopet.\allowbreak{}entity.\allowbreak{}Mascota}; proyecciones
  \texttt{Resumen\allowbreak{}Especie}, \texttt{Historial\allowbreak{}Clinico},
  \texttt{Reporte\allowbreak{}Dashboard}; migraciones
  \texttt{V5\_\allowbreak{}\_\allowbreak{}procedimientos\_\allowbreak{}b\allowbreak{}iopet.\allowbreak{}sql} y
  \texttt{V6\_\allowbreak{}\_\allowbreak{}formalizar\_\allowbreak{}proc\allowbreak{}edimientos\_\allowbreak{}jpa.\allowbreak{}sql}; pruebas
  \texttt{Backend/\allowbreak{}src/\allowbreak{}test/\allowbreak{}java/\allowbreak{}com/\allowbreak{}biopet/\allowbreak{}repository/\allowbreak{}Procedimientos\allowbreak{}Biopet\allowbreak{}Integration\allowbreak{}Test.\allowbreak{}java},
  \texttt{Resumen\allowbreak{}Especies\allowbreak{}Integration\allowbreak{}Test.\allowbreak{}java},
  \texttt{Biopet\allowbreak{}App\allowbreak{}Rol\allowbreak{}Minimo\allowbreak{}Privilegios\allowbreak{}Integration\allowbreak{}Test.\allowbreak{}java} y
  \texttt{Backend/\allowbreak{}src/\allowbreak{}test/\allowbreak{}java/\allowbreak{}com/\allowbreak{}biopet/\allowbreak{}Sql\allowbreak{}Injection\allowbreak{}Security\allowbreak{}Test.\allowbreak{}java};
  guion \texttt{scripts/\allowbreak{}audit-\allowbreak{}sql-\allowbreak{}dynamic.\allowbreak{}sh} (objetivo
  \texttt{sql-\allowbreak{}audit} del \texttt{Makefile} y job \texttt{sql-\allowbreak{}audit} del
  CI); catálogo \texttt{docs/\allowbreak{}basedatos/\allowbreak{}CATALOGO-\allowbreak{}SP.\allowbreak{}md}; decisión
  \texttt{docs/\allowbreak{}adr/\allowbreak{}ADR-\allowbreak{}007-\allowbreak{}acceso-\allowbreak{}datos.\allowbreak{}md}; requisitos que consumen
  estas rutinas: REQ-F-021, REQ-F-013, REQ-F-015, REQ-F-018, REQ-F-019.
\item
  \textbf{Estado:} verificado
\item
  \textbf{Observaciones:} la reproducción de
  \texttt{scripts/\allowbreak{}audit-\allowbreak{}sql-\allowbreak{}dynamic.\allowbreak{}sh} se generó \textbf{después} del
  tag \texttt{v1.\allowbreak{}0.\allowbreak{}0}, sobre el commit exacto de ese tag
  (\texttt{0d5cd525ce648cca\allowbreak{}7219da204e16fa62\allowbreak{}2e671a87}), en un worktree
  aislado que no modificó el árbol de trabajo principal: el código
  auditado es el del tag; el log es evidencia posterior, no un artefacto
  que existiera dentro de ese commit. La rutina
  \texttt{fn\_\allowbreak{}siguiente\_\allowbreak{}numer\allowbreak{}o\_\allowbreak{}ficha} existe, está probada
  (\texttt{siguiente\allowbreak{}Numero\allowbreak{}Ficha\_\allowbreak{}formato\allowbreak{}YSecuencia},
  \texttt{siguiente\allowbreak{}Numero\allowbreak{}Ficha\_\allowbreak{}prefijo\allowbreak{}Vacio\allowbreak{}Usa\allowbreak{}Default}) y cumple la
  categoría ``6) Código secuencial'' de la rúbrica, pero \textbf{ningún
  servicio de la aplicación la invoca todavía}: su consumidor natural es
  REQ-F-018 (numeración correlativa de comprobantes), que está
  \texttt{pendiente}. Se declara aquí para no presentar como
  funcionalidad de negocio lo que hoy es sólo un objeto de base de datos
  disponible.
\end{itemize}

\subsubsection{REQ-NF-014 --- Comportamiento del sistema ante
indisponibilidad de
Redis}\label{req-nf-014-comportamiento-del-sistema-ante-indisponibilidad-de-redis}

\begin{itemize}
\tightlist
\item
  \textbf{Tipo:} Seguridad / Disponibilidad
\item
  \textbf{Prioridad:} Must
\item
  \textbf{Enunciado:} Cuando el servicio Redis no esté disponible, el
  sistema deberá aplicar una política explícita y documentada de
  degradación: las verificaciones de revocación de token deberán
  \textbf{fallar de forma cerrada} (\emph{fail-closed}), rechazando la
  solicitud con \texttt{503\ Service\ Unavai\allowbreak{}lable} y cuerpo
  \texttt{Problem\allowbreak{}Detail} en lugar de conceder acceso sin comprobar la
  lista negra; y las operaciones cuya única dependencia de Redis sea la
  caché de lectura deberán \textbf{degradarse a consulta directa a la
  base de datos} (\emph{fail-open} limitado a la caché), sin exponer
  trazas internas al cliente.
\item
  \textbf{Rationale:} requisito nuevo incorporado en esta revisión. La
  retroalimentación del docente-director señaló que el SRS no
  especificaba qué debe ocurrir si Redis cae, pese a que Redis sostiene
  tres mecanismos distintos: la lista negra de tokens (REQ-F-005), la
  caché de listados (REQ-NF-012) y la caché de la API externa
  (REQ-F-025). La revisión del código confirma que \textbf{hoy no existe
  ninguna política}: ni \emph{fail-open} ni \emph{fail-closed}
  deliberados, sino un fallo no controlado (ver Observaciones). Una
  caída de Redis no debe poder convertirse en una omisión silenciosa del
  control de revocación de tokens.
\item
  \textbf{Criterios de aceptación:}

  \begin{enumerate}
  \def\labelenumi{\arabic{enumi}.}
  \tightlist
  \item
    Dado Redis no disponible y una solicitud a un recurso protegido con
    un access token, cuando \texttt{Token\allowbreak{}Blacklist\allowbreak{}Service.\allowbreak{}is\allowbreak{}Revoked} no
    puede consultarse, entonces la respuesta es
    \texttt{503\ Service\ Unavai\allowbreak{}lable} con
    \texttt{Content-\allowbreak{}Type:\ application/\allowbreak{}problem+json}, y \textbf{no} se
    establece autenticación en el contexto de seguridad.
  \item
    Dado Redis no disponible, la respuesta del criterio 1 no contiene
    traza de pila, nombre de clase ni mensaje de la excepción de
    conexión.
  \item
    Dado Redis no disponible y una solicitud a
    \texttt{POST\ /\allowbreak{}api/\allowbreak{}auth/\allowbreak{}logout}, entonces la respuesta es
    \texttt{503} con \texttt{Problem\allowbreak{}Detail}, sin error 500 sin
    controlar.
  \item
    Dado Redis no disponible y una solicitud a
    \texttt{GET\ /\allowbreak{}api/\allowbreak{}mascotas} de un usuario ya autenticado por otra
    vía, entonces la operación se resuelve consultando directamente
    PostgreSQL y responde \texttt{200\ OK}: el fallo de la caché no
    impide la lectura.
  \item
    Cada caída y cada recuperación de la conexión con Redis queda
    registrada en el log del sistema con marca de tiempo, siguiendo el
    formato de REQ-NF-009.
  \item
    Existe una prueba automatizada que fuerza el fallo de conexión con
    Redis y comprueba los criterios 1, 3 y 4.
  \end{enumerate}
\item
  \textbf{Método de verificación:} previsto --- prueba de integración
  que arranque la aplicación con un \texttt{String\allowbreak{}Redis\allowbreak{}Template} que
  lance \texttt{Redis\allowbreak{}Connection\allowbreak{}Failure\allowbreak{}Exception} (o con un contenedor
  Redis detenido) y verifique los códigos de respuesta de los criterios
  1, 3 y 4; más un \texttt{Cache\allowbreak{}Error\allowbreak{}Handler} registrado en la
  configuración de caché para el criterio 4 y un manejador dedicado en
  \texttt{Global\allowbreak{}Exception\allowbreak{}Handler} (o en
  \texttt{Jwt\allowbreak{}Authentication\allowbreak{}Filter}) para los criterios 1 a 3. \textbf{No
  ejecutable hoy: no hay implementación.}
\item
  \textbf{Trazabilidad:}
  \texttt{com.\allowbreak{}biopet.\allowbreak{}security.\allowbreak{}Token\allowbreak{}Blacklist\allowbreak{}Service} (métodos
  \texttt{revoke} e \texttt{is\allowbreak{}Revoked}, ambos invocan
  \texttt{String\allowbreak{}Redis\allowbreak{}Template} sin captura de error),
  \texttt{com.\allowbreak{}biopet.\allowbreak{}security.\allowbreak{}Jwt\allowbreak{}Authentication\allowbreak{}Filter} (bloque
  \texttt{catch\ (Jwt\allowbreak{}Exception\ \textbar{}\ Illeg\allowbreak{}al\allowbreak{}Argument\allowbreak{}Exception\ ex)}),
  \texttt{com.\allowbreak{}biopet.\allowbreak{}integration.\allowbreak{}External\allowbreak{}Api\allowbreak{}Service.\allowbreak{}guardar\allowbreak{}En\allowbreak{}Cache}
  (única ruta con manejo de error de Redis ya implementado, probada por
  \texttt{External\allowbreak{}Api\allowbreak{}Service\allowbreak{}Test.\allowbreak{}falla\allowbreak{}Al\allowbreak{}Escribir\allowbreak{}En\allowbreak{}Redis\allowbreak{}No\allowbreak{}Rompe\allowbreak{}La\allowbreak{}Respuesta}),
  \texttt{com.\allowbreak{}biopet.\allowbreak{}service.\allowbreak{}Mascota\allowbreak{}Service} y \texttt{Consulta\allowbreak{}Service}
  (anotaciones \texttt{@Cacheable}/\texttt{@Cache\allowbreak{}Evict}), configuración
  \texttt{spring.\allowbreak{}data.\allowbreak{}redis.*} y \texttt{spring.\allowbreak{}cache.\allowbreak{}type:\ redis} en
  \texttt{application.\allowbreak{}yml}, servicio \texttt{redis} de
  \texttt{docker-\allowbreak{}compose.\allowbreak{}yml}, decisión
  \texttt{docs/\allowbreak{}adr/\allowbreak{}ADR-\allowbreak{}003-\allowbreak{}jwt-\allowbreak{}redis.\allowbreak{}md}; requisitos afectados
  REQ-F-005, REQ-F-025, REQ-NF-012.
\item
  \textbf{Estado:} pendiente
\item
  \textbf{Observaciones:} comportamiento \textbf{real} verificado en el
  código actual, sin suponer nada:
  \texttt{Token\allowbreak{}Blacklist\allowbreak{}Service.\allowbreak{}is\allowbreak{}Revoked} invoca
  \texttt{redis\allowbreak{}Template.\allowbreak{}has\allowbreak{}Key(...)} sin bloque \texttt{try}. Con Redis
  caído esa llamada lanza \texttt{Redis\allowbreak{}Connection\allowbreak{}Failure\allowbreak{}Exception}, que
  es una \texttt{Runtime\allowbreak{}Exception} de la jerarquía
  \texttt{Data\allowbreak{}Access\allowbreak{}Exception} de Spring y \textbf{no} es capturada por
  el
  \texttt{catch\ (Jwt\allowbreak{}Exception\ \textbar{}\ Illeg\allowbreak{}al\allowbreak{}Argument\allowbreak{}Exception\ ex)}
  de \texttt{Jwt\allowbreak{}Authentication\allowbreak{}Filter.\allowbreak{}do\allowbreak{}Filter\allowbreak{}Internal}. La excepción
  escapa del filtro antes de llegar al \texttt{Dispatcher\allowbreak{}Servlet}, por
  lo que tampoco la atiende \texttt{Global\allowbreak{}Exception\allowbreak{}Handler} (un
  \texttt{@Rest\allowbreak{}Controller\allowbreak{}Advice} no cubre la cadena de filtros): el
  resultado es un error \texttt{500} del contenedor, sin
  \texttt{Problem\allowbreak{}Detail}. Lo mismo ocurre en
  \texttt{Token\allowbreak{}Blacklist\allowbreak{}Service.\allowbreak{}revoke} durante el logout. Una búsqueda
  de \texttt{Redis\allowbreak{}Connection\allowbreak{}Failure}, \texttt{Redis\allowbreak{}System\allowbreak{}Exception},
  \texttt{Data\allowbreak{}Access\allowbreak{}Exception} o \texttt{Cache\allowbreak{}Error\allowbreak{}Handler} sobre
  \texttt{Backend/\allowbreak{}src/\allowbreak{}main/\allowbreak{}java/} no devuelve ninguna coincidencia:
  \textbf{no hay manejo alguno}. Por tanto no se afirma aquí que el
  sistema sea \emph{fail-open} ni \emph{fail-closed}: hoy es un fallo no
  controlado, y este requisito especifica la política que
  \textbf{deberá} implementarse. El estado es \texttt{pendiente} y no
  \texttt{implementado} porque no existe ningún código que intente
  aplicar esta política. La única excepción parcial ya implementada es
  la escritura en caché de \texttt{External\allowbreak{}Api\allowbreak{}Service.\allowbreak{}guardar\allowbreak{}En\allowbreak{}Cache},
  que sí captura el fallo y no rompe la respuesta; se cita en
  Trazabilidad como patrón de referencia para la implementación futura.
  \textbf{Corregir esto requiere modificar código de backend y queda
  explícitamente fuera del alcance de esta revisión documental.}
\end{itemize}

\subsubsection{REQ-NF-015 --- Aislamiento de datos por propietario
(requisito
transversal)}\label{req-nf-015-aislamiento-de-datos-por-propietario-requisito-transversal}

\begin{itemize}
\item
  \textbf{Tipo:} Seguridad
\item
  \textbf{Prioridad:} Must
\item
  \textbf{Enunciado:} Para todo recurso del sistema que tenga un
  propietario ---mascotas, citas, consultas médicas, vacunas y cualquier
  recurso de propietario que se incorpore en el futuro---, el sistema
  deberá garantizar que un usuario autenticado con rol
  \texttt{ROLE\_\allowbreak{}DUENO} acceda exclusivamente a los recursos asociados a
  su propia cuenta, tanto en las operaciones de listado como en las de
  consulta por identificador, devolviendo \texttt{403\ Forbidden} con
  cuerpo \texttt{Problem\allowbreak{}Detail} ante cualquier intento de acceso a un
  recurso de otro propietario y excluyendo del resultado, sin excepción,
  los elementos ajenos.
\item
  \textbf{Rationale:} requisito transversal nuevo incorporado en esta
  revisión. La regla de aislamiento existía dispersa en los requisitos
  individuales (REQ-F-009, REQ-F-010, REQ-F-015, REQ-F-024) y en la
  implementación de cada servicio, pero \textbf{no estaba especificada
  una sola vez de forma transversal}, lo que permitió que un módulo
  ---Consultas--- se entregara sin ella sin que ningún requisito lo
  detectara. Se formaliza aquí para que cualquier módulo nuevo herede la
  obligación explícitamente y para que el incumplimiento sea trazable a
  un requisito propio.
\item
  \textbf{Criterios de aceptación:}

  \begin{enumerate}
  \def\labelenumi{\arabic{enumi}.}
  \tightlist
  \item
    \textbf{Listados.} Para cada uno de \texttt{GET\ /\allowbreak{}api/\allowbreak{}mascotas},
    \texttt{GET\ /\allowbreak{}api/\allowbreak{}citas}, \texttt{GET\ /\allowbreak{}api/\allowbreak{}consultas} y
    \texttt{GET\ /\allowbreak{}api/\allowbreak{}vacunas}: dado un usuario \texttt{ROLE\_\allowbreak{}DUENO} A y
    existiendo al menos un recurso perteneciente a un dueño B distinto,
    cuando A invoca el listado con cualquier combinación de parámetros
    de paginación, entonces ningún elemento de ninguna página devuelta
    pertenece a B.
  \item
    \textbf{Consulta por identificador.} Para cada uno de
    \texttt{GET\ /\allowbreak{}api/\allowbreak{}mascotas/\allowbreak{}\{id\}}, \texttt{GET\ /\allowbreak{}api/\allowbreak{}citas/\allowbreak{}\{id\}},
    \texttt{GET\ /\allowbreak{}api/\allowbreak{}consultas/\allowbreak{}\{id\}} y
    \texttt{GET\ /\allowbreak{}api/\allowbreak{}vacunas/\allowbreak{}\{id\}}: dado el dueño A y un recurso de
    B, cuando A lo solicita por su identificador, entonces la respuesta
    es \texttt{403} con \texttt{Problem\allowbreak{}Detail} de tipo
    \texttt{urn:biopet:error:\allowbreak{}forbidden}, no \texttt{200} ni
    \texttt{404}.
  \item
    \textbf{Listado filtrado.} Dado el dueño A y una mascota de B,
    cuando A invoca \texttt{GET\ /\allowbreak{}api/\allowbreak{}vacunas/\allowbreak{}mascota/\allowbreak{}\{mascota\allowbreak{}Id\}} con
    el identificador de la mascota de B, entonces la respuesta es
    \texttt{403}.
  \item
    \textbf{Agregados.} Dado el dueño A, cuando invoca
    \texttt{GET\ /\allowbreak{}api/\allowbreak{}mascotas/\allowbreak{}resumen-\allowbreak{}especies} indicando el
    \texttt{duenio\allowbreak{}Id} de B, entonces el resultado se calcula sobre A: el
    parámetro ajeno se ignora y nunca se devuelven totales de B.
  \item
    \textbf{Caché.} Dos dueños distintos que solicitan la misma página
    con los mismos parámetros no comparten entrada de caché ni
    resultados.
  \item
    \textbf{Cobertura.} Cada recurso de propietario del sistema tiene al
    menos una prueba automatizada de acceso cruzado (dueño A intentando
    acceder a un recurso de B) para el listado y para la consulta por
    identificador.
  \end{enumerate}
\item
  \textbf{Método de verificación:} pruebas automatizadas de acceso
  cruzado:
  \texttt{Mascota\allowbreak{}Controller\allowbreak{}Test.\allowbreak{}dueno\allowbreak{}Solo\allowbreak{}Ve\allowbreak{}Sus\allowbreak{}Propias\allowbreak{}Mascotas\allowbreak{}En\allowbreak{}Listado},
  \texttt{dueno\allowbreak{}Consulta\allowbreak{}Mascota\allowbreak{}De\allowbreak{}Otro\allowbreak{}Duenio\allowbreak{}Devuelve403},
  \texttt{dos\allowbreak{}Duenos\allowbreak{}Con\allowbreak{}Misma\allowbreak{}Pagina\allowbreak{}No\allowbreak{}Comparen\allowbreak{}Resultados\allowbreak{}De\allowbreak{}Cache};
  \texttt{Cita\allowbreak{}Controller\allowbreak{}Test.\allowbreak{}dueno\allowbreak{}Solo\allowbreak{}Ve\allowbreak{}Sus\allowbreak{}Propias\allowbreak{}Citas\allowbreak{}En\allowbreak{}Listado},
  \texttt{dueno\allowbreak{}No\allowbreak{}Puede\allowbreak{}Consultar\allowbreak{}Cita\allowbreak{}De\allowbreak{}Mascota\allowbreak{}Ajena};
  \texttt{Vacuna\allowbreak{}Controller\allowbreak{}Test.\allowbreak{}dueno\allowbreak{}Solo\allowbreak{}Ve\allowbreak{}Vacunas\allowbreak{}De\allowbreak{}Sus\allowbreak{}Propias\allowbreak{}Mascotas},
  \texttt{dueno\allowbreak{}Consulta\allowbreak{}Vacuna\allowbreak{}De\allowbreak{}Mascota\allowbreak{}Ajena\allowbreak{}Devuelve403};
  \texttt{Consulta\allowbreak{}Controller\allowbreak{}Test.\allowbreak{}dueno\allowbreak{}De\allowbreak{}Otra\allowbreak{}Mascota\allowbreak{}No\allowbreak{}Puede\allowbreak{}Ver\allowbreak{}Consulta\allowbreak{}Ajena};
  más la evidencia HTTP real de acceso cruzado entre dos dueños sobre
  TLS en \texttt{docs/\allowbreak{}mediciones/\allowbreak{}sec/\allowbreak{}A01-\allowbreak{}access-\allowbreak{}control.\allowbreak{}md} y
  \texttt{docs/\allowbreak{}mediciones/\allowbreak{}sec/\allowbreak{}raw/\allowbreak{}A01-\allowbreak{}access-\allowbreak{}control.\allowbreak{}txt}. Criterio 4:
  inspección de \texttt{Mascota\allowbreak{}Service.\allowbreak{}resumen\allowbreak{}Por\allowbreak{}Especie}. Criterio 1
  para Consultas y criterio 3: \textbf{sin artefacto} --- ver
  Observaciones.
\item
  \textbf{Trazabilidad:} reglas de propiedad implementadas en
  \texttt{Mascota\allowbreak{}Service.\allowbreak{}verificar\allowbreak{}Propiedad},
  \texttt{Cita\allowbreak{}Service.\allowbreak{}verificar\allowbreak{}Acceso\allowbreak{}Lectura},
  \texttt{Vacuna\allowbreak{}Service.\allowbreak{}verificar\allowbreak{}Acceso} y
  \texttt{Consulta\allowbreak{}Service.\allowbreak{}verificar\allowbreak{}Acceso}; métodos de repositorio con
  filtro por propietario
  \texttt{Mascota\allowbreak{}Repository.\allowbreak{}find\allowbreak{}All\allowbreak{}By\allowbreak{}Duenio\allowbreak{}Id\allowbreak{}And\allowbreak{}Activo\allowbreak{}True},
  \texttt{Cita\allowbreak{}Repository.\allowbreak{}find\allowbreak{}All\allowbreak{}By\allowbreak{}Mascota\_\allowbreak{}Duenio\_\allowbreak{}Id\allowbreak{}And\allowbreak{}Activo\allowbreak{}True},
  \texttt{Vacuna\allowbreak{}Repository.\allowbreak{}find\allowbreak{}All\allowbreak{}By\allowbreak{}Mascota\_\allowbreak{}Duenio\_\allowbreak{}Id\allowbreak{}And\allowbreak{}Activo\allowbreak{}True}
  (\textbf{no existe el equivalente en \texttt{Consulta\allowbreak{}Repository}});
  matriz de permisos, sección 2.8; evidencia
  \texttt{docs/\allowbreak{}mediciones/\allowbreak{}sec/\allowbreak{}A01-\allowbreak{}access-\allowbreak{}control.\allowbreak{}md}; control OWASP A01;
  requisito funcional relacionado REQ-F-006.
\item
  \textbf{Estado:} implementado
\item
  \textbf{Observaciones:} \textbf{el requisito NO está cumplido.} El
  criterio 1 \textbf{falla para \texttt{GET\ /\allowbreak{}api/\allowbreak{}consultas}}, y el
  criterio 6 falla para ese mismo endpoint.
  \texttt{Consulta\allowbreak{}Service.\allowbreak{}listar()} devuelve
  \texttt{consulta\allowbreak{}Repository.\allowbreak{}find\allowbreak{}All\allowbreak{}By\allowbreak{}Activo\allowbreak{}True(pageable)} para todos
  los roles, incluido \texttt{ROLE\_\allowbreak{}DUENO}, y
  \texttt{Consulta\allowbreak{}Repository} no declara ningún método de filtrado por
  propietario. En consecuencia, un dueño autenticado puede ver en ese
  listado las consultas médicas de mascotas ajenas: el aislamiento
  transversal que exige el enunciado \textbf{no se cumple hoy}.

  \textbf{Por qué \texttt{implementado} y no \texttt{pendiente}.} Según
  la sección 1.3.1, \texttt{pendiente} significa que \emph{no existe
  implementación específica de lo que el requisito exige}, y eso sería
  falso aquí: tres de los cuatro recursos de propietario ---mascotas,
  citas y vacunas--- implementan la regla con métodos de repositorio
  filtrados por dueño y con pruebas de acceso cruzado que pasan
  (\texttt{dueno\allowbreak{}Solo\allowbreak{}Ve\allowbreak{}Sus\allowbreak{}Propias\allowbreak{}Mascotas\allowbreak{}En\allowbreak{}Listado},
  \texttt{dueno\allowbreak{}Solo\allowbreak{}Ve\allowbreak{}Sus\allowbreak{}Propias\allowbreak{}Citas\allowbreak{}En\allowbreak{}Listado},
  \texttt{dueno\allowbreak{}Solo\allowbreak{}Ve\allowbreak{}Vacunas\allowbreak{}De\allowbreak{}Sus\allowbreak{}Propias\allowbreak{}Mascotas}), además de la captura
  HTTP real de IDOR entre dos dueños. Declararlo \texttt{pendiente}
  ocultaría esa implementación real y describiría mal el defecto, que es
  \textbf{una omisión en un módulo concreto}, no la ausencia del
  mecanismo. Ambos estados significan igualmente ``requisito no
  cumplido'' (sección 1.3.1) y este requisito se contabiliza como
  \textbf{no cumplido} en las secciones 3.4.1 y 3.4.2;
  \texttt{implementado} es el que describe con exactitud la situación.

  Los criterios 2, 4 y 5 sí están verificados para los cuatro recursos;
  el criterio 3 se sostiene por inspección de
  \texttt{Vacuna\allowbreak{}Service.\allowbreak{}listar\allowbreak{}Por\allowbreak{}Mascota} (ver REQ-F-024), sin prueba
  dedicada. \textbf{La corrección del criterio 1 requiere modificar
  código de backend y queda explícitamente fuera del alcance de esta
  revisión documental}; se reporta como acción de seguimiento en la
  sección 7.4.
\end{itemize}

\subsubsection{REQ-NF-016 --- Política de
contraseñas}\label{req-nf-016-poluxedtica-de-contraseuxf1as}

\begin{itemize}
\tightlist
\item
  \textbf{Tipo:} Seguridad
\item
  \textbf{Prioridad:} Must
\item
  \textbf{Enunciado:} El sistema deberá exigir que toda contraseña
  establecida por un usuario ---en el registro público, en la creación
  administrativa de una cuenta y en su actualización--- tenga una
  longitud mínima de 8 y máxima de 80 caracteres, rechazando con
  \texttt{422\ Unprocessable\ E\allowbreak{}ntity} las que no cumplan, y no deberá
  devolver nunca la contraseña ni su hash en ninguna respuesta de la
  API.
\item
  \textbf{Rationale:} requisito nuevo incorporado en esta revisión.
  REQ-NF-008 cubría el \textbf{almacenamiento} cifrado de la contraseña,
  pero ninguna especificación recogía la \textbf{política de admisión}
  que el código ya aplica mediante
  \texttt{@Size(min\ =\ 8,\allowbreak{}\ max\ =\ 80)}. Se documenta la política real,
  sin atribuirle reglas de complejidad, expiración o comprobación contra
  listas de contraseñas filtradas que el sistema no implementa.
\item
  \textbf{Criterios de aceptación:}

  \begin{enumerate}
  \def\labelenumi{\arabic{enumi}.}
  \tightlist
  \item
    Dado un registro público con contraseña de menos de 8 caracteres,
    cuando se envía, entonces la respuesta es \texttt{422} con
    \texttt{Problem\allowbreak{}Detail} de validación y el objeto \texttt{errors}
    incluye la entrada \texttt{password}.
  \item
    Dado el mismo caso en la creación administrativa de usuarios
    (\texttt{POST\ /\allowbreak{}api/\allowbreak{}usuarios}), entonces la respuesta es
    \texttt{422} con la misma estructura.
  \item
    Dada una creación administrativa \textbf{sin} contraseña, entonces
    la respuesta es \texttt{400} con \texttt{Problem\allowbreak{}Detail}: la
    contraseña es obligatoria al crear.
  \item
    Dada una actualización administrativa con contraseña vacía o
    ausente, entonces la contraseña actual se conserva sin cambios y la
    operación responde \texttt{200\ OK}.
  \item
    Ninguna respuesta de la API incluye la contraseña ni el hash.
  \item
    La contraseña admitida se almacena cifrada según REQ-NF-008.
  \end{enumerate}
\item
  \textbf{Método de verificación:} pruebas automatizadas
  \texttt{Auth\allowbreak{}Controller\allowbreak{}Test.\allowbreak{}registro\allowbreak{}Con\allowbreak{}Campos\allowbreak{}Invalidos}, que envía
  \texttt{"password":"cort\allowbreak{}a"} y asserta \texttt{422} con
  \texttt{\$.\allowbreak{}errors.\allowbreak{}password} como arreglo (criterio 1);
  \texttt{Usuario\allowbreak{}Controller\allowbreak{}Test.\allowbreak{}crear\allowbreak{}Usuario\allowbreak{}Con\allowbreak{}Campos\allowbreak{}Invalidos\allowbreak{}Devuelve422}
  (criterio 2); \texttt{crear\allowbreak{}Usuario\allowbreak{}Sin\allowbreak{}Password\allowbreak{}Devuelve400} (criterio
  3); \texttt{actualizar\allowbreak{}Usuario\allowbreak{}Devuelve200}, que omite el campo
  \texttt{password} y obtiene \texttt{200\ OK}, más inspección de
  \texttt{Usuario\allowbreak{}Service.\allowbreak{}actualizar}, cuya guarda
  \texttt{if\ (request.\allowbreak{}password()\ !=\ null\ \&\&\ !\allowbreak{}request.\allowbreak{}password().\allowbreak{}is\allowbreak{}Blank())}
  deja intacto \texttt{password\allowbreak{}Hash} cuando no se envía contraseña
  (criterio 4). Criterio 5: inspección de \texttt{Usuario\allowbreak{}Response},
  \texttt{Auth\allowbreak{}Response} y \texttt{Auth\allowbreak{}Session\allowbreak{}Response}, ninguno de los
  cuales declara campo de contraseña. Criterio 6: ver REQ-NF-008.
  \textbf{El límite superior de 80 caracteres del criterio 2 no tiene
  ninguna prueba automatizada} --- ver Observaciones.
\item
  \textbf{Trazabilidad:} restricciones declarativas
  \texttt{@Not\allowbreak{}Blank\ @Size(min\ =\ 8,\allowbreak{}\ max\ =\ 80)} en
  \texttt{com.\allowbreak{}biopet.\allowbreak{}dto.\allowbreak{}Registro\allowbreak{}Request} y
  \texttt{@Size(min\ =\ 8,\allowbreak{}\ max\ =\ 80)} en
  \texttt{com.\allowbreak{}biopet.\allowbreak{}dto.\allowbreak{}Usuario\allowbreak{}Request} (sin \texttt{@Not\allowbreak{}Blank},
  deliberadamente, para permitir la actualización sin cambio de
  contraseña); \texttt{Usuario\allowbreak{}Service.\allowbreak{}crear} (validación explícita de
  contraseña obligatoria al crear);
  \texttt{Global\allowbreak{}Exception\allowbreak{}Handler.\allowbreak{}validacion} →
  \texttt{Problem\allowbreak{}Type.\allowbreak{}VALIDATION}; pruebas \texttt{Auth\allowbreak{}Controller\allowbreak{}Test} y
  \texttt{Usuario\allowbreak{}Controller\allowbreak{}Test}; requisito de almacenamiento
  REQ-NF-008; requisito dependiente REQ-F-026.
\item
  \textbf{Estado:} implementado
\item
  \textbf{Observaciones:} el requisito \textbf{no está cumplido} por
  falta de evidencia, no por defecto: la política está realmente
  implementada, pero una de las reglas que exige el enunciado no está
  probada. \textbf{Cobertura real, regla por regla:} longitud mínima de
  8 en el registro público →
  \texttt{Auth\allowbreak{}Controller\allowbreak{}Test.\allowbreak{}registro\allowbreak{}Con\allowbreak{}Campos\allowbreak{}Invalidos} (envía
  \texttt{"corta"} y asserta \texttt{422} con
  \texttt{\$.\allowbreak{}errors.\allowbreak{}password}); longitud mínima de 8 en la creación
  administrativa →
  \texttt{Usuario\allowbreak{}Controller\allowbreak{}Test.\allowbreak{}crear\allowbreak{}Usuario\allowbreak{}Con\allowbreak{}Campos\allowbreak{}Invalidos\allowbreak{}Devuelve422}
  (mismo valor y misma aserción); contraseña obligatoria al crear →
  \texttt{crear\allowbreak{}Usuario\allowbreak{}Sin\allowbreak{}Password\allowbreak{}Devuelve400}; conservación de la
  contraseña al actualizar → \texttt{actualizar\allowbreak{}Usuario\allowbreak{}Devuelve200} más
  inspección de \texttt{Usuario\allowbreak{}Service.\allowbreak{}actualizar}; ausencia de la
  contraseña en las respuestas → inspección de los tres records de
  salida; cifrado → REQ-NF-008. \textbf{El límite superior de 80
  caracteres no está cubierto por ninguna prueba}: sólo consta en la
  anotación \texttt{@Size(min\ =\ 8,\allowbreak{}\ max\ =\ 80)} de
  \texttt{Registro\allowbreak{}Request} y \texttt{Usuario\allowbreak{}Request}. No existe ningún
  caso que envíe una contraseña de más de 80 caracteres y compruebe el
  \texttt{422}. Mientras no exista esa prueba, el requisito no puede
  declararse \texttt{verificado}. Por lo demás, la política implementada
  es \textbf{solo de longitud}: el sistema \textbf{no} exige complejidad
  (mayúsculas, dígitos o símbolos), \textbf{no} aplica caducidad,
  \textbf{no} impide la reutilización de contraseñas anteriores y
  \textbf{no} comprueba la contraseña contra listas de credenciales
  filtradas. Se documenta lo que existe, sin declarar conformidad con
  ninguna guía externa de contraseñas que el proyecto no haya aplicado.
\end{itemize}

\subsubsection{REQ-NF-017 --- Respaldo y recuperación de la base de
datos}\label{req-nf-017-respaldo-y-recuperaciuxf3n-de-la-base-de-datos}

\begin{itemize}
\tightlist
\item
  \textbf{Tipo:} Operación
\item
  \textbf{Prioridad:} Should
\item
  \textbf{Enunciado:} El proyecto deberá mantener documentado un
  procedimiento de respaldo y restauración de la base de datos
  PostgreSQL que especifique frecuencia, formato, destino, retención y
  verificación de integridad, y deberá haber demostrado al menos una
  restauración real sobre un volcado generado con ese mismo
  procedimiento, comprobando que los datos y los procedimientos
  almacenados sobreviven a la restauración.
\item
  \textbf{Rationale:} requisito nuevo incorporado en esta revisión. El
  procedimiento y su prueba de restauración ya existían en el
  repositorio (\texttt{docs/\allowbreak{}despliegue/\allowbreak{}BACKUP.\allowbreak{}md}) pero no estaban
  recogidos como requisito, de modo que ni el SRS ni la matriz los
  declaraban. Se documenta lo que existe en \textbf{este} repositorio,
  sin tomar prácticas de ningún otro proyecto.
\item
  \textbf{Criterios de aceptación:}

  \begin{enumerate}
  \def\labelenumi{\arabic{enumi}.}
  \tightlist
  \item
    El procedimiento documentado especifica: qué se respalda,
    frecuencia, formato del volcado, destino separado del host de
    producción, retención y método de verificación por suma de
    comprobación.
  \item
    Existe evidencia fechada de una restauración real: los comandos
    exactos ejecutados y el resultado obtenido, sobre un volcado
    generado con el procedimiento documentado.
  \item
    Tras la restauración, la comprobación de integridad confirma el
    número de filas esperado en las tablas del sistema.
  \item
    Tras la restauración, los procedimientos almacenados siguen siendo
    invocables: al menos una rutina se ejecuta correctamente sobre la
    base restaurada.
  \item
    Se conserva en el repositorio un volcado de ejemplo del mismo
    formato, para revisión.
  \end{enumerate}
\item
  \textbf{Método de verificación:} inspección de
  \texttt{docs/\allowbreak{}despliegue/\allowbreak{}BACKUP.\allowbreak{}md}, secciones 0 y 1 (criterio 1) y
  sección 4 (criterios 2 a 4): restauración ejecutada el
  \textbf{2026-08-17} sobre PostgreSQL 16 en el contenedor
  \texttt{biopet-\allowbreak{}postgres}, con
  \texttt{pg\_\allowbreak{}dump\ -\/-\allowbreak{}clean\ -\/-\allowbreak{}if-\allowbreak{}exists\ -\/-\allowbreak{}no-\allowbreak{}owner} (64 762
  bytes, 6 tablas), restauración en una base de prueba
  \texttt{biopet\_\allowbreak{}restore\_\allowbreak{}tes\allowbreak{}t}, verificación de conteos
  (\texttt{usuarios}: 1 fila; \texttt{mascotas}, \texttt{citas},
  \texttt{consultas}, \texttt{vacunas}: 0 filas, consistente con la base
  de origen) e invocación posterior de
  \texttt{fn\_\allowbreak{}siguiente\_\allowbreak{}numer\allowbreak{}o\_\allowbreak{}ficha(\textquotesingle{}RST\textquotesingle{})}
  devolviendo \texttt{RST-\allowbreak{}000001}. Criterio 5: archivo
  \texttt{docs/\allowbreak{}despliegue/\allowbreak{}ejemplo-\allowbreak{}backup-\allowbreak{}20260817.\allowbreak{}sql} presente en el
  repositorio.
\item
  \textbf{Trazabilidad:} \texttt{docs/\allowbreak{}despliegue/\allowbreak{}BACKUP.\allowbreak{}md}
  (procedimiento y evidencia),
  \texttt{docs/\allowbreak{}despliegue/\allowbreak{}ejemplo-\allowbreak{}backup-\allowbreak{}20260817.\allowbreak{}sql} (volcado de
  ejemplo), \texttt{docs/\allowbreak{}despliegue/\allowbreak{}RUNBOOK.\allowbreak{}md} y \texttt{DEPLOYMENT.\allowbreak{}md}
  (operación del despliegue), \texttt{render.\allowbreak{}yaml} (servicio de base de
  datos gestionado), rutinas de \texttt{db/\allowbreak{}procs/} que sobreviven a la
  restauración (REQ-NF-013), migraciones Flyway en
  \texttt{Backend/\allowbreak{}src/\allowbreak{}main/\allowbreak{}resources/\allowbreak{}db/\allowbreak{}migration/}.
\item
  \textbf{Estado:} implementado
\item
  \textbf{Observaciones:} el requisito \textbf{no está cumplido}: dos de
  sus cinco criterios se apoyan en un relato, no en un artefacto
  reproducible. \textbf{Toda la evidencia citada pertenece a este
  repositorio (ENTREGA-FINAL-BIOPET); no se toma nada de ningún otro
  proyecto.} Comprobación independiente realizada en esta revisión sobre
  el artefacto \texttt{docs/\allowbreak{}despliegue/\allowbreak{}ejemplo-\allowbreak{}backup-\allowbreak{}20260817.\allowbreak{}sql} (64
  762 bytes, codificado en UTF-16): contiene exactamente las \textbf{6
  tablas} del esquema de este proyecto (\texttt{usuarios},
  \texttt{mascotas}, \texttt{citas}, \texttt{consultas},
  \texttt{vacunas}, \texttt{flyway\_\allowbreak{}schema\_\allowbreak{}hist\allowbreak{}ory}), \textbf{6
  secuencias} ---incluida \texttt{seq\_\allowbreak{}ficha\_\allowbreak{}biopet}---, \textbf{4
  triggers} y \textbf{7 rutinas}
  (\texttt{fn\_\allowbreak{}resumen\_\allowbreak{}mascota\allowbreak{}s\_\allowbreak{}por\_\allowbreak{}especie},
  \texttt{fn\_\allowbreak{}historial\_\allowbreak{}clini\allowbreak{}co\_\allowbreak{}mascota},
  \texttt{fn\_\allowbreak{}reporte\_\allowbreak{}dashboa\allowbreak{}rd},
  \texttt{fn\_\allowbreak{}siguiente\_\allowbreak{}numer\allowbreak{}o\_\allowbreak{}ficha},
  \texttt{sp\_\allowbreak{}actualizar\_\allowbreak{}esta\allowbreak{}do\_\allowbreak{}citas\_\allowbreak{}masivas},
  \texttt{sp\_\allowbreak{}registrar\_\allowbreak{}consu\allowbreak{}lta\_\allowbreak{}validada} y la función de trigger
  \texttt{set\_\allowbreak{}actualizado\_\allowbreak{}en}), y los datos son \texttt{usuarios}: 1
  fila y \texttt{mascotas}, \texttt{citas}, \texttt{consultas},
  \texttt{vacunas}: 0 filas. Esto confirma los criterios 3 y 5 y
  corrobora el contenido declarado en \texttt{BACKUP.\allowbreak{}md}. \textbf{Lo que
  no está respaldado:} los criterios 2 y 4 ---que la restauración se
  ejecutó realmente y que una rutina se invocó con éxito sobre la base
  restaurada--- constan únicamente como narración en \texttt{BACKUP.\allowbreak{}md}
  sección 4; \textbf{no existe ningún log crudo archivado} de esa
  ejecución, a diferencia de otras evidencias del proyecto
  (\texttt{docs/\allowbreak{}mediciones/\allowbreak{}sec/\allowbreak{}raw/}, las corridas k6 o las de
  Lighthouse). Mientras no se archive esa salida, el requisito no puede
  declararse \texttt{verificado}. Además, el respaldo \textbf{no está
  automatizado dentro del repositorio}: no existe ningún guion en
  \texttt{scripts/} ni tarea programada que lo ejecute.
  \texttt{BACKUP.\allowbreak{}md} describe un procedimiento manual del operador y
  cierra con una lista de comprobación (backup diario programado,
  registro de sumas, retención de 30 días, restauración de prueba
  mensual) que es responsabilidad operativa y \textbf{no} está marcada
  como cumplida de forma continua.
\end{itemize}

\subsubsection{REQ-NF-018 --- Accesibilidad de la
interfaz}\label{req-nf-018-accesibilidad-de-la-interfaz}

\begin{itemize}
\tightlist
\item
  \textbf{Tipo:} Usabilidad
\item
  \textbf{Prioridad:} Should
\item
  \textbf{Enunciado:} Las pantallas auditadas de la aplicación
  ---\texttt{/\allowbreak{}login} y \texttt{/\allowbreak{}mascotas}--- deberán obtener una
  puntuación de la categoría \emph{Accessibility} de Lighthouse igual o
  superior a 90 sobre 100, tanto en perfil móvil como en perfil de
  escritorio, y ese umbral deberá estar declarado como condición de
  fallo en la configuración de medición. Este requisito \textbf{no}
  exige ni declara conformidad con WCAG en ningún nivel.
\item
  \textbf{Rationale:} requisito nuevo incorporado en esta revisión. La
  accesibilidad se medía desde la Tercera Entrega como parte de las
  corridas Lighthouse (y su resultado se citaba dentro de REQ-NF-005),
  pero no existía un requisito propio con un umbral declarado, de modo
  que el valor medido no estaba contrastado contra ninguna exigencia.
\item
  \textbf{Criterios de aceptación:}

  \begin{enumerate}
  \def\labelenumi{\arabic{enumi}.}
  \tightlist
  \item
    En las corridas Lighthouse sobre la ruta \texttt{/\allowbreak{}login}, la
    puntuación de Accessibility es ≥ 90 en perfil móvil y en perfil de
    escritorio.
  \item
    En las corridas Lighthouse sobre la ruta \texttt{/\allowbreak{}mascotas} (vista
    autenticada), la puntuación de Accessibility es ≥ 90 en ambos
    perfiles.
  \item
    El umbral de 0,90 está declarado a nivel \texttt{error} en las dos
    configuraciones de medición (perfil móvil y perfil de escritorio),
    de modo que una regresión por debajo de 90 haga fallar la
    comprobación.
  \item
    Los archivos crudos de las corridas quedan archivados en el
    repositorio junto con un registro de procedencia que indique fecha,
    commit auditado, versiones de herramienta, URLs y perfiles.
  \end{enumerate}
\item
  \textbf{Método de verificación:} medición Lighthouse --- 12 archivos
  crudos archivados en
  \texttt{docs/\allowbreak{}mediciones/\allowbreak{}lighthouse/\allowbreak{}lhci-\allowbreak{}20260818-\allowbreak{}0538-*.\allowbreak{}json}
  (perfiles móvil y desktop × rutas \texttt{/\allowbreak{}login} y \texttt{/\allowbreak{}mascotas}
  × 3 corridas). Lectura directa del campo
  \texttt{categories.\allowbreak{}accessibility.\allowbreak{}score} de los 12 archivos: \textbf{91
  en los doce casos}, sin ninguna dispersión, por encima del umbral de
  90 (criterios 1 y 2). Criterio 3: inspección de
  \texttt{lighthouserc.\allowbreak{}js} (línea 47) y \texttt{lighthouserc.\allowbreak{}desktop.\allowbreak{}js}
  (línea 56), ambas con
  \texttt{\textquotesingle{}categories:acc\allowbreak{}essibility\textquotesingle{}:\ {[}\textquotesingle{}\allowbreak{}error\textquotesingle{},\allowbreak{}\ \{\ min\allowbreak{}Score:\ 0.\allowbreak{}9\ \}{]}}.
  Criterio 4: archivo de procedencia
  \texttt{docs/\allowbreak{}mediciones/\allowbreak{}lighthouse/\allowbreak{}lhci-\allowbreak{}20260818-\allowbreak{}0538.\allowbreak{}meta.\allowbreak{}txt}, que
  registra \texttt{fecha\_\allowbreak{}iso8601:\ 202\allowbreak{}6-\allowbreak{}08-\allowbreak{}18\allowbreak{}T05:38:06\allowbreak{}Z},
  \texttt{commit\_\allowbreak{}hash\_\allowbreak{}corto:\ e\allowbreak{}94bb72},
  \texttt{node\_\allowbreak{}version:\ v24.\allowbreak{}18.\allowbreak{}0},
  \texttt{lighthouse\_\allowbreak{}cli\_\allowbreak{}ver\allowbreak{}sion:\ 0.\allowbreak{}14.\allowbreak{}0}, las dos URLs auditadas y
  los dos perfiles con sus respectivas configuraciones.
\item
  \textbf{Trazabilidad:} configuración \texttt{lighthouserc.\allowbreak{}js} y
  \texttt{lighthouserc.\allowbreak{}desktop.\allowbreak{}js} (umbrales), guion
  \texttt{scripts/\allowbreak{}run-\allowbreak{}lighthouse.\allowbreak{}sh}, evidencias
  \texttt{docs/\allowbreak{}mediciones/\allowbreak{}lighthouse/} con \texttt{README.\allowbreak{}md}
  (procedencia y anonimización declarada) y
  \texttt{lhci-\allowbreak{}20260818-\allowbreak{}0538.\allowbreak{}meta.\allowbreak{}txt}; pantallas auditadas
  \texttt{frontend/\allowbreak{}src/\allowbreak{}app/\allowbreak{}features/\allowbreak{}login.\allowbreak{}component.\allowbreak{}ts} y
  \texttt{mascotas.\allowbreak{}component.\allowbreak{}ts}; requisito hermano de usabilidad
  REQ-NF-005.
\item
  \textbf{Estado:} verificado
\item
  \textbf{Observaciones:} el alcance de este requisito está
  deliberadamente limitado a lo que la evidencia existente demuestra, y
  no más: \textbf{(a) Sin declaración de norma.} \textbf{No se declara
  conformidad con WCAG 2.1 en ningún nivel (A, AA o AAA)}, ni con
  ninguna otra norma de accesibilidad. No se ha realizado auditoría
  manual, ni pruebas con lector de pantalla, ni revisión completa de
  navegación por teclado; Lighthouse cubre sólo un subconjunto
  automatizable de los criterios WCAG. La métrica obtenida (91/100) es
  la puntuación de esa herramienta, no un nivel de conformidad.
  \textbf{(b) Alcance de pantallas.} Las rutas auditadas son exactamente
  \texttt{/\allowbreak{}login} y \texttt{/\allowbreak{}mascotas}. La pantalla de vacunas
  (\texttt{frontend/\allowbreak{}src/\allowbreak{}app/\allowbreak{}features/\allowbreak{}vacunas.\allowbreak{}component.\allowbreak{}ts}) \textbf{no}
  está incluida en las corridas archivadas, por lo que el requisito no
  afirma nada sobre ella. \textbf{(c) Entorno de medición.} Las 12
  corridas se ejecutaron contra \texttt{http:/\allowbreak{}/\allowbreak{}localhost:4200} sobre el
  commit \texttt{e94bb72}, no contra el despliegue de producción.
  \textbf{(d) Sumas de comprobación.} Los archivos
  \texttt{docs/\allowbreak{}mediciones/\allowbreak{}lighthouse/\allowbreak{}SHA256\allowbreak{}SUMS.\allowbreak{}txt} y
  \texttt{SHA256\allowbreak{}SUMS-\allowbreak{}ORIGINAL.\allowbreak{}txt} corresponden a la corrida
  \textbf{original del 2026-08-01} (carpeta \texttt{raw/}) y \textbf{no}
  cubren los 12 archivos \texttt{lhci-\allowbreak{}20260818-*} citados aquí; por eso
  el criterio 4 se apoya en el registro de procedencia
  \texttt{lhci-\allowbreak{}20260818-\allowbreak{}0538.\allowbreak{}meta.\allowbreak{}txt} y no en esas sumas. Archivar las
  sumas de la re-corrida del 2026-08-18 queda como mejora de
  trazabilidad.
\end{itemize}

\subsubsection{3.3. Criterios INVEST aplicados a las historias de
usuario}\label{criterios-invest-aplicados-a-las-historias-de-usuario}

Las 24 historias de usuario de
\texttt{docs/\allowbreak{}requisitos/\allowbreak{}historias/\allowbreak{}Historias\allowbreak{}Usuario.\allowbreak{}md} se evalúan aquí
explícitamente contra los seis criterios INVEST de Cohn (Independent,
Negotiable, Valuable, Estimable, Small, Testable), con justificación
breve y verificable contra campos que ya existen en cada historia (no se
introduce información nueva para esta evaluación).

\begin{quote}
\textbf{Nota de vocabulario.} La taxonomía de tres estados de la sección
1.3.1 se aplica a los \textbf{requisitos}
(\texttt{REQ-\allowbreak{}F}/\texttt{REQ-\allowbreak{}NF}) de las secciones 3.1 y 3.2. Las
historias de usuario conservan su propio campo de estado en
\texttt{Historias\allowbreak{}Usuario.\allowbreak{}md}; las dos escalas no se mezclan ni se
contradicen.
\end{quote}

\textbf{Negotiable, Valuable y Testable se cumplen en las 24 historias,
sin excepción}, verificado estructuralmente así:

\begin{itemize}
\tightlist
\item
  \textbf{Negotiable:} las 24 historias siguen el formato Connextra
  (``Como\ldots{} quiero\ldots{} para\ldots{}''), que expresa intención
  y valor, no una solución de diseño fija; el ``cómo'' (campos exactos,
  validaciones) queda abierto a negociación en el caso de uso y en la
  implementación.
\item
  \textbf{Valuable:} cada historia declara un rol beneficiario explícito
  (\texttt{Como\ \textless{}rol\textgreater{}}) y está trazada a un
  requisito funcional con \emph{rationale} propio en la sección 3.1,
  nunca a una tarea puramente técnica sin beneficiario.
\item
  \textbf{Testable:} las 24 historias incluyen un bloque
  \texttt{gherkin} con al menos un escenario \texttt{Given/\allowbreak{}When/\allowbreak{}Then}
  verificable; las que aún no están implementadas lo declaran
  explícitamente como ``comportamiento esperado, no implementado'', sin
  fingir que ya existe evidencia de ejecución.
\end{itemize}

\textbf{Independent, Estimable y Small varían por historia}
(justificación por fila, tomada del campo \emph{Dependencias} y del
estado real de cada historia):

{\def\LTcaptype{none} % do not increment counter
\begin{longtable}[]{@{}
  >{\raggedright\arraybackslash}p{(\linewidth - 6\tabcolsep) * \real{0.2500}}
  >{\raggedright\arraybackslash}p{(\linewidth - 6\tabcolsep) * \real{0.2500}}
  >{\raggedright\arraybackslash}p{(\linewidth - 6\tabcolsep) * \real{0.2500}}
  >{\raggedright\arraybackslash}p{(\linewidth - 6\tabcolsep) * \real{0.2500}}@{}}
\toprule\noalign{}
\begin{minipage}[b]{\linewidth}\raggedright
HU
\end{minipage} & \begin{minipage}[b]{\linewidth}\raggedright
Independent
\end{minipage} & \begin{minipage}[b]{\linewidth}\raggedright
Estimable
\end{minipage} & \begin{minipage}[b]{\linewidth}\raggedright
Small
\end{minipage} \\
\midrule\noalign{}
\endhead
\bottomrule\noalign{}
\endlastfoot
HU-001 & Sí & Sí --- implementada, esfuerzo ya observado & Nota --- 2
REQ-F relacionados \\
HU-002 & Sí & Sí --- implementada, esfuerzo ya observado & Sí --- 1
REQ-F \\
HU-003 & Parcial --- depende de HU-002 (requiere haber iniciado sesión
previamente) & Sí --- implementada, esfuerzo ya observado & Sí --- 1
REQ-F \\
HU-004 & Parcial --- depende de HU-002 & Sí --- implementada, esfuerzo
ya observado & Sí --- 1 REQ-F \\
HU-005 & Parcial --- depende de HU-002 & Sí --- implementada, esfuerzo
ya observado & Sí --- 1 REQ-F \\
HU-006 & Parcial --- depende de HU-002 & Sí --- implementada, esfuerzo
ya observado & Sí --- 1 REQ-F \\
HU-007 & Parcial --- depende de HU-005 (control de acceso por rol) & Sí
--- implementada, esfuerzo ya observado & Sí --- 1 REQ-F \\
HU-008 & Parcial --- depende de HU-005 & Sí --- implementada, esfuerzo
ya observado & Sí --- 1 REQ-F \\
HU-009 & Parcial --- depende de HU-005, HU-008 & Sí --- implementada,
esfuerzo ya observado & Sí --- 1 REQ-F \\
HU-010 & Parcial --- depende de HU-005, HU-009 & Sí --- implementada,
esfuerzo ya observado & Sí --- 1 REQ-F \\
HU-011 & Parcial --- depende de HU-005, HU-009 & Sí --- implementada,
esfuerzo ya observado & Sí --- 1 REQ-F \\
HU-012 & Parcial --- depende de HU-007 (requiere que la mascota exista)
& Parcial --- hay código y pruebas reales para la obligación A de
REQ-F-013; el esfuerzo de la obligación B sigue siendo estimación & Nota
--- 1 REQ-F con 2 obligaciones (REQ-F-013) \\
HU-013 & Parcial --- depende de HU-012 & Parcial --- sin artefacto
entregado que medir & Sí --- 1 REQ-F \\
HU-014 & Parcial --- depende de HU-007 & Parcial --- la obligación A de
REQ-F-015 está cubierta y probada; el esfuerzo del calendario
(obligación B) sigue siendo estimación & Nota --- 1 REQ-F con 2
obligaciones (REQ-F-015) \\
HU-015 & Parcial --- depende de HU-007 & Parcial --- sin artefacto
entregado que medir & Sí --- 1 REQ-F \\
HU-016 & Parcial --- depende de HU-012 (requiere historial clínico) &
Parcial --- sin artefacto entregado que medir & Sí --- 1 REQ-F \\
HU-017 & Parcial --- depende de HU-007 & Parcial --- sin artefacto
entregado que medir & Sí --- 1 REQ-F \\
HU-018 & Sí & Parcial --- sin artefacto entregado que medir & Sí --- 1
REQ-F \\
HU-019 & Parcial --- depende de HU-002 (requiere usuario autenticado) &
Parcial --- sólo la parte de autenticación tiene artefacto medible & Sí
--- 1 REQ-F \\
HU-020 & Parcial --- depende de HU-005 y HU-007 & Sí --- implementada,
esfuerzo ya observado & Sí --- 1 REQ-F \\
HU-021 & Parcial --- depende de HU-001 (registro de usuario) & Parcial
--- sin artefacto entregado que medir & Sí --- 1 REQ-F \\
HU-022 & Parcial --- depende de HU-005 (control de acceso por rol) & Sí
--- implementada, esfuerzo ya observado & Sí --- 1 REQ-F \\
HU-023 & Parcial --- depende de HU-007 (requiere que la mascota exista)
& Sí --- implementada, esfuerzo ya observado & Sí --- 1 REQ-F \\
HU-024 & Sí & Sí --- implementada, esfuerzo ya observado & Sí --- 1
REQ-F \\
\end{longtable}
}

\textbf{Lectura honesta de ``Independent''.} La mayoría de historias
depende de al menos otra (típicamente HU-002 ``sesión iniciada'' o
HU-005/HU-007 ``rol autorizado''/``mascota existente''), lo cual es
normal y esperado en un sistema con autenticación y CRUD relacional:
ninguna historia de gestión de un recurso es realmente independiente de
``existe sesión'' y ``el recurso padre existe''. Se documenta la
dependencia real en vez de declarar ``Independent: Sí'' de forma
genérica sin sustento, que habría sido la alternativa más fácil pero
menos honesta.

\textbf{Lectura honesta de ``Estimable''.} Las historias ya
implementadas son estimables con certeza retrospectiva: el código y las
pruebas existen, su esfuerzo ya fue observado. Las parciales (HU-012,
HU-014, HU-019) son estimables sólo para el subconjunto ya entregado ---
el esfuerzo del resto (historial consolidado expuesto por API,
calendario interactivo, auditoría genérica de CRUD) sigue siendo una
estimación no verificada. Las pendientes no tienen ningún artefacto que
permita estimar con base empírica; su tamaño es una suposición heredada
de la Entrega 1A, no una medición.

\textbf{Nota sobre ``Small''.} Tres historias agrupan más de un
requisito. HU-001 agrupa REQ-F-001 y REQ-F-002 porque ambos comparten el
mismo flujo de entrada (\texttt{POST\ /\allowbreak{}api/\allowbreak{}auth/\allowbreak{}registro}) y serían
artificialmente pequeños por separado. HU-012 y HU-014 corresponden cada
una a un único requisito (REQ-F-013 y REQ-F-015) que agrupa dos
obligaciones: ni la historia ni el requisito se dividen en
identificadores nuevos ---el valor para el usuario sigue siendo uno
solo---; lo que se separa explícitamente dentro del bloque son sus
obligaciones A y B, que es lo que permite describir con exactitud qué
parte está hecha y cuál no.

\textbf{Nota sobre REQ-F-026.} El requisito de recuperación de
contraseña incorporado en esta revisión \textbf{no tiene historia de
usuario ni caso de uso todavía}; se deberán crear la historia y el caso
de uso correspondientes al planificar su implementación. No se inventa
aquí una historia que no exista en \texttt{Historias\allowbreak{}Usuario.\allowbreak{}md}, para no
introducir una referencia colgante.

\subsubsection{3.4. Tabla de control de completitud de
requisitos}\label{tabla-de-control-de-completitud-de-requisitos}

Auditoría mecánica de los campos obligatorios de la plantilla (sección
1.3.2) sobre los \textbf{44 requisitos} de las secciones 3.1 y 3.2 ---
26 funcionales y 18 no funcionales, uno por identificador. Cada celda
indica si el campo está presente y con contenido sustantivo
(\textbf{OK}) o ausente (\textbf{FALTA}). El campo \emph{Observaciones}
no se audita aquí porque es opcional por diseño.

Abreviaturas de columna: \textbf{Tip} = Tipo · \textbf{Pri} = Prioridad
· \textbf{Enu} = Enunciado · \textbf{Rat} = Rationale · \textbf{Cri} =
Criterios de aceptación · \textbf{Ver} = Método de verificación ·
\textbf{Tra} = Trazabilidad · \textbf{Est} = Estado (valor declarado).

{\def\LTcaptype{none} % do not increment counter
\begin{longtable}[]{@{}lllllllll@{}}
\toprule\noalign{}
ID & Tip & Pri & Enu & Rat & Cri & Ver & Tra & Est \\
\midrule\noalign{}
\endhead
\bottomrule\noalign{}
\endlastfoot
REQ-F-001 & OK & OK & OK & OK & OK & OK & OK & verificado \\
REQ-F-002 & OK & OK & OK & OK & OK & OK & OK & verificado \\
REQ-F-003 & OK & OK & OK & OK & OK & OK & OK & verificado \\
REQ-F-004 & OK & OK & OK & OK & OK & OK & OK & verificado \\
REQ-F-005 & OK & OK & OK & OK & OK & OK & OK & verificado \\
REQ-F-006 & OK & OK & OK & OK & OK & OK & OK & implementado \\
REQ-F-007 & OK & OK & OK & OK & OK & OK & OK & verificado \\
REQ-F-008 & OK & OK & OK & OK & OK & OK & OK & verificado \\
REQ-F-009 & OK & OK & OK & OK & OK & OK & OK & verificado \\
REQ-F-010 & OK & OK & OK & OK & OK & OK & OK & verificado \\
REQ-F-011 & OK & OK & OK & OK & OK & OK & OK & verificado \\
REQ-F-012 & OK & OK & OK & OK & OK & OK & OK & verificado \\
REQ-F-013 & OK & OK & OK & OK & OK & OK & OK & implementado \\
REQ-F-014 & OK & OK & OK & OK & OK & OK & OK & pendiente \\
REQ-F-015 & OK & OK & OK & OK & OK & OK & OK & implementado \\
REQ-F-016 & OK & OK & OK & OK & OK & OK & OK & pendiente \\
REQ-F-017 & OK & OK & OK & OK & OK & OK & OK & pendiente \\
REQ-F-018 & OK & OK & OK & OK & OK & OK & OK & pendiente \\
REQ-F-019 & OK & OK & OK & OK & OK & OK & OK & pendiente \\
REQ-F-020 & OK & OK & OK & OK & OK & OK & OK & implementado \\
REQ-F-021 & OK & OK & OK & OK & OK & OK & OK & verificado \\
REQ-F-022 & OK & OK & OK & OK & OK & OK & OK & pendiente \\
REQ-F-023 & OK & OK & OK & OK & OK & OK & OK & verificado \\
REQ-F-024 & OK & OK & OK & OK & OK & OK & OK & verificado \\
REQ-F-025 & OK & OK & OK & OK & OK & OK & OK & verificado \\
REQ-F-026 & OK & OK & OK & OK & OK & OK & OK & pendiente \\
REQ-NF-001 & OK & OK & OK & OK & OK & OK & OK & verificado \\
REQ-NF-002 & OK & OK & OK & OK & OK & OK & OK & verificado \\
REQ-NF-003 & OK & OK & OK & OK & OK & OK & OK & verificado \\
REQ-NF-004 & OK & OK & OK & OK & OK & OK & OK & verificado \\
REQ-NF-005 & OK & OK & OK & OK & OK & OK & OK & verificado \\
REQ-NF-006 & OK & OK & OK & OK & OK & OK & OK & pendiente \\
REQ-NF-007 & OK & OK & OK & OK & OK & OK & OK & implementado \\
REQ-NF-008 & OK & OK & OK & OK & OK & OK & OK & verificado \\
REQ-NF-009 & OK & OK & OK & OK & OK & OK & OK & verificado \\
REQ-NF-010 & OK & OK & OK & OK & OK & OK & OK & verificado \\
REQ-NF-011 & OK & OK & OK & OK & OK & OK & OK & verificado \\
REQ-NF-012 & OK & OK & OK & OK & OK & OK & OK & implementado \\
REQ-NF-013 & OK & OK & OK & OK & OK & OK & OK & verificado \\
REQ-NF-014 & OK & OK & OK & OK & OK & OK & OK & pendiente \\
REQ-NF-015 & OK & OK & OK & OK & OK & OK & OK & implementado \\
REQ-NF-016 & OK & OK & OK & OK & OK & OK & OK & implementado \\
REQ-NF-017 & OK & OK & OK & OK & OK & OK & OK & implementado \\
REQ-NF-018 & OK & OK & OK & OK & OK & OK & OK & verificado \\
\end{longtable}
}

\textbf{Resultado de la auditoría: 0 celdas FALTA sobre 44 requisitos ×
8 campos obligatorios (352 celdas).}

\paragraph{3.4.1. Resumen por estado}\label{resumen-por-estado}

{\def\LTcaptype{none} % do not increment counter
\begin{longtable}[]{@{}llll@{}}
\toprule\noalign{}
Estado & REQ-F & REQ-NF & Total \\
\midrule\noalign{}
\endhead
\bottomrule\noalign{}
\endlastfoot
\texttt{verificado} (cumplidos) & 15 & 11 & 26 \\
\texttt{implementado} (no cumplidos) & 4 & 5 & 9 \\
\texttt{pendiente} (no cumplidos) & 7 & 2 & 9 \\
\textbf{Total de requisitos} & \textbf{26} & \textbf{18} &
\textbf{44} \\
\end{longtable}
}

\textbf{El número oficial de requisitos del SRS es 44}:
\texttt{REQ-\allowbreak{}F-\allowbreak{}001} a \texttt{REQ-\allowbreak{}F-\allowbreak{}026} (26 funcionales) y
\texttt{REQ-\allowbreak{}NF-\allowbreak{}001} a \texttt{REQ-\allowbreak{}NF-\allowbreak{}018} (18 no funcionales). Un
requisito, un identificador, una fila en esta tabla y una fila en
\texttt{docs/\allowbreak{}trazabilidad/\allowbreak{}matriz.\allowbreak{}csv}: los conjuntos de identificadores
de ambos documentos son \textbf{idénticos}, sin excedentes ni faltantes
en ninguna dirección. Tres de esos 44 (\texttt{REQ-\allowbreak{}F-\allowbreak{}013},
\texttt{REQ-\allowbreak{}F-\allowbreak{}015} y \texttt{REQ-\allowbreak{}NF-\allowbreak{}012}) agrupan dos obligaciones cada
uno, separadas dentro de su bloque (sección 1.3.2) sin generar
identificadores adicionales.

De los 44, \textbf{26 están cumplidos} (\texttt{verificado}) y
\textbf{18 no lo están} (9 \texttt{implementado} y 9
\texttt{pendiente}), conforme a la lectura obligada de la sección 1.3.1.

\paragraph{3.4.2. Cobertura de los requisitos de prioridad
Must}\label{cobertura-de-los-requisitos-de-prioridad-must}

Hay \textbf{25 requisitos \texttt{Must}} (11 funcionales y 14 no
funcionales). De ellos, \textbf{20 están \texttt{verificado}} y
\textbf{5 no}. No se declara una cobertura de 25/25:

{\def\LTcaptype{none} % do not increment counter
\begin{longtable}[]{@{}
  >{\raggedright\arraybackslash}p{(\linewidth - 4\tabcolsep) * \real{0.3333}}
  >{\raggedright\arraybackslash}p{(\linewidth - 4\tabcolsep) * \real{0.3333}}
  >{\raggedright\arraybackslash}p{(\linewidth - 4\tabcolsep) * \real{0.3333}}@{}}
\toprule\noalign{}
\begin{minipage}[b]{\linewidth}\raggedright
Requisito Must
\end{minipage} & \begin{minipage}[b]{\linewidth}\raggedright
Estado
\end{minipage} & \begin{minipage}[b]{\linewidth}\raggedright
Por qué no está \texttt{verificado}
\end{minipage} \\
\midrule\noalign{}
\endhead
\bottomrule\noalign{}
\endlastfoot
REQ-F-006 & implementado & \texttt{Consulta\allowbreak{}Service.\allowbreak{}listar()} no filtra
por propietario para \texttt{ROLE\_\allowbreak{}DUENO}; el criterio 3 de aislamiento
falla. \\
REQ-NF-007 & implementado & El endpoint de salud está verificado, pero
no existe registro fechado de sondeos que permita calcular un porcentaje
de disponibilidad (criterio 4). \\
REQ-NF-014 & pendiente & No existe ninguna política implementada ante
indisponibilidad de Redis: hoy el fallo no está controlado. \\
REQ-NF-015 & implementado & El aislamiento falla en el listado de
consultas médicas (misma causa que REQ-F-006). \\
REQ-NF-016 & implementado & El límite superior de 80 caracteres de la
política de contraseñas no tiene ninguna prueba automatizada. \\
\end{longtable}
}

Dos de los cinco comparten una única causa raíz ---la falta de filtrado
por propietario en \texttt{Consulta\allowbreak{}Service.\allowbreak{}listar()}--- y un tercero es
la ausencia de política ante fallo de Redis; ambas están registradas en
la sección 7.4 como acciones que \textbf{requieren modificar código de
backend} y quedan fuera del alcance de esta revisión documental. Los
otros dos (REQ-NF-007 y REQ-NF-016) se cierran archivando evidencia, sin
tocar código.

\subsection{4. Trazabilidad (resumen)}\label{trazabilidad-resumen}

La matriz completa vive en \texttt{docs/\allowbreak{}trazabilidad/\allowbreak{}matriz.\allowbreak{}csv} (bloque
A.3.3 de la Guía), con una fila por requisito y las columnas
\texttt{id\_\allowbreak{}requisito}, \texttt{tipo}, \texttt{prioridad\_\allowbreak{}moscow},
\texttt{historia\_\allowbreak{}usuario}, \texttt{caso\_\allowbreak{}de\_\allowbreak{}uso},
\texttt{modulo\_\allowbreak{}codigo}, \texttt{endpoint\_\allowbreak{}api},
\texttt{prueba\_\allowbreak{}automatiza\allowbreak{}da}, \texttt{tipo\_\allowbreak{}acceso},
\texttt{evidencia\_\allowbreak{}empiric\allowbreak{}a} y \texttt{estado}. Este SRS es la fuente de
verdad para las columnas \texttt{id\_\allowbreak{}requisito}, \texttt{tipo},
\texttt{prioridad\_\allowbreak{}moscow} y \texttt{estado}; la matriz no debe declarar
un requisito que no exista aquí, ni omitir ninguno de los declarados en
la sección 3.

La consistencia entre ambos documentos la comprueba automáticamente
\texttt{scripts/\allowbreak{}validate-\allowbreak{}traceability.\allowbreak{}sh} (job \texttt{traceability} del
CI), que verifica que todo \texttt{REQ-\allowbreak{}F-}/\texttt{REQ-\allowbreak{}NF-} del SRS
tenga fila en la matriz y viceversa, que ninguna fila quede sin
historia, caso de uso ni prueba, que toda referencia
\texttt{HU-}/\texttt{CU-} exista realmente en
\texttt{Historias\allowbreak{}Usuario.\allowbreak{}md}/\texttt{Casos\allowbreak{}De\allowbreak{}Uso.\allowbreak{}md}, y que no haya
identificadores duplicados.

\textbf{Nota sobre los requisitos con dos obligaciones.}
\texttt{REQ-\allowbreak{}F-\allowbreak{}013}, \texttt{REQ-\allowbreak{}F-\allowbreak{}015} y \texttt{REQ-\allowbreak{}NF-\allowbreak{}012} agrupan
cada uno dos obligaciones, separadas dentro de su bloque como
\emph{Obligación A} y \emph{Obligación B} (sección 1.3.2). \textbf{No
generan identificadores adicionales.} Cada uno tiene un único
identificador, un único bloque en la sección 3, una única fila en la
tabla de control 3.4 y una única fila en
\texttt{docs/\allowbreak{}trazabilidad/\allowbreak{}matriz.\allowbreak{}csv}, con \textbf{un solo} valor de
Estado, que describe el requisito completo: basta con que una de sus dos
obligaciones no esté cumplida para que el requisito no pueda declararse
\texttt{verificado}. El campo \texttt{evidencia\_\allowbreak{}empiric\allowbreak{}a} de la fila de
la matriz detalla la situación de cada obligación por separado. Así la
trazabilidad histórica hacia las historias, los casos de uso y los
identificadores \texttt{RF-\allowbreak{}NN} de la Entrega 1A se mantiene intacta.

\textbf{Correspondencia exacta de identificadores.} El conjunto de
identificadores de la sección 3 de este SRS y el conjunto de la columna
\texttt{id\_\allowbreak{}requisito} de \texttt{docs/\allowbreak{}trazabilidad/\allowbreak{}matriz.\allowbreak{}csv} son
\textbf{el mismo conjunto de 44 elementos}: la diferencia SRS − matriz
es vacía y la diferencia matriz − SRS es vacía. Esta igualdad la
comprueba \texttt{scripts/\allowbreak{}validate-\allowbreak{}traceability.\allowbreak{}sh} en el CI y, campo a
campo, \texttt{docs/\allowbreak{}requisitos/\allowbreak{}tools/\allowbreak{}auditoria-\allowbreak{}srs.\allowbreak{}py}.

\textbf{Correspondencia REQ-F ↔ HU ↔ CU:}

{\def\LTcaptype{none} % do not increment counter
\begin{longtable}[]{@{}llllll@{}}
\toprule\noalign{}
REQ-F & HU & CU & REQ-F & HU & CU \\
\midrule\noalign{}
\endhead
\bottomrule\noalign{}
\endlastfoot
001 & HU-001 & CU-01 & 014 & HU-013 & CU-13 \\
002 & HU-001 & CU-01 & 015 & HU-014 & CU-14 \\
003 & HU-002 & CU-02 & 016 & HU-015 & CU-15 \\
004 & HU-003 & CU-03 & 017 & HU-016 & CU-16 \\
005 & HU-004 & CU-04 & 018 & HU-017 & CU-17 \\
006 & HU-005 & CU-05 & 019 & HU-018 & CU-18 \\
007 & HU-006 & CU-06 & 020 & HU-019 & CU-19 \\
008 & HU-007 & CU-07 & 021 & HU-020 & CU-20 \\
009 & HU-008 & CU-08 & 022 & HU-021 & CU-21 \\
010 & HU-009 & CU-09 & 023 & HU-022 & CU-22 \\
011 & HU-010 & CU-10 & 024 & HU-023 & CU-23 \\
012 & HU-011 & CU-11 & 025 & HU-024 & CU-24 \\
013 & HU-012 & CU-12 & 026 & \emph{(sin HU/CU aún)} & \emph{(sin HU/CU
aún)} \\
\end{longtable}
}

\texttt{REQ-\allowbreak{}F-\allowbreak{}026} es el único requisito funcional sin historia ni caso
de uso: se incorporó en esta revisión y su historia y caso de uso
deberán crearse al planificar su implementación. No se declara ninguna
referencia a historias o casos de uso inexistentes.

\textbf{Trazabilidad de los requisitos no funcionales.} Ningún
\texttt{REQ-\allowbreak{}NF} declara historia de usuario, porque ninguno la tiene: su
trazabilidad apunta a decisiones de arquitectura, configuración, clases,
pruebas, guiones, mediciones, evidencias y normas, tal como se detalla
en el campo \emph{Trazabilidad} de cada bloque de la sección 3.2. Este
es un cambio deliberado respecto de la práctica de inventar historias
para rellenar la columna.

\textbf{Trazabilidad de los procedimientos almacenados.} Las seis
rutinas de \texttt{db/\allowbreak{}procs/} se declaran en la trazabilidad del
requisito al que sirven, no como requisitos propios:

{\def\LTcaptype{none} % do not increment counter
\begin{longtable}[]{@{}
  >{\raggedright\arraybackslash}p{(\linewidth - 4\tabcolsep) * \real{0.3333}}
  >{\raggedright\arraybackslash}p{(\linewidth - 4\tabcolsep) * \real{0.3333}}
  >{\raggedright\arraybackslash}p{(\linewidth - 4\tabcolsep) * \real{0.3333}}@{}}
\toprule\noalign{}
\begin{minipage}[b]{\linewidth}\raggedright
Rutina
\end{minipage} & \begin{minipage}[b]{\linewidth}\raggedright
Requisito al que sirve
\end{minipage} & \begin{minipage}[b]{\linewidth}\raggedright
Situación
\end{minipage} \\
\midrule\noalign{}
\endhead
\bottomrule\noalign{}
\endlastfoot
\texttt{fn\_\allowbreak{}resumen\_\allowbreak{}mascota\allowbreak{}s\_\allowbreak{}por\_\allowbreak{}especie} & REQ-F-021 & Consumida por
\texttt{Mascota\allowbreak{}Service.\allowbreak{}resumen\allowbreak{}Por\allowbreak{}Especie} y expuesta en
\texttt{GET\ /\allowbreak{}api/\allowbreak{}mascotas/\allowbreak{}resumen-\allowbreak{}especies}. \\
\texttt{fn\_\allowbreak{}historial\_\allowbreak{}clini\allowbreak{}co\_\allowbreak{}mascota} & REQ-F-013 (obligación B) &
Invocable desde el repositorio y probada; sin endpoint que la
exponga. \\
\texttt{fn\_\allowbreak{}reporte\_\allowbreak{}dashboa\allowbreak{}rd} & REQ-F-019 & Invocable desde el
repositorio y probada; sin endpoint ni exportación. \\
\texttt{sp\_\allowbreak{}registrar\_\allowbreak{}consu\allowbreak{}lta\_\allowbreak{}validada} & REQ-F-013 (obligación A) &
Validación cruzada de mascota y veterinario; probada, no invocada desde
\texttt{Consulta\allowbreak{}Service}. \\
\texttt{sp\_\allowbreak{}actualizar\_\allowbreak{}esta\allowbreak{}do\_\allowbreak{}citas\_\allowbreak{}masivas} & REQ-F-015 (obligación
A) & Cierre masivo de citas por veterinario y fecha límite; probada, no
expuesta como endpoint. \\
\texttt{fn\_\allowbreak{}siguiente\_\allowbreak{}numer\allowbreak{}o\_\allowbreak{}ficha} & REQ-NF-013 (y, a futuro,
REQ-F-018) & Generador de códigos correlativos; probado, sin consumidor
en la aplicación todavía. \\
\end{longtable}
}

Ninguna de las seis genera un requisito funcional nuevo: cuatro
implementan la capa de datos de requisitos que ya existen, una es
consumida en producción y la sexta es un objeto disponible sin
consumidor, declarado como tal en las Observaciones de REQ-NF-013.

\subsubsection{4.1. Trazabilidad histórica: identificadores originales →
identificadores
actuales}\label{trazabilidad-histuxf3rica-identificadores-originales-identificadores-actuales}

La retroalimentación oficial del SGA de la Entrega 1A señaló que ``los
RF-WEB se remapean a RF-16/RF-17 sin matriz de trazabilidad explícita en
esta entrega'' (ver \texttt{docs/\allowbreak{}observaciones/\allowbreak{}OBSERVACIONES.\allowbreak{}md},
OBS-03). La siguiente tabla consolida el vínculo entre los
identificadores originales (\texttt{RF-\allowbreak{}NN}/\texttt{RF-\allowbreak{}WEB-\allowbreak{}NN}, Entrega
1A) y los actuales (\texttt{REQ-\allowbreak{}F-\allowbreak{}NNN}), sin omitir ningún requisito
funcional y sin inventar ningún origen.

{\def\LTcaptype{none} % do not increment counter
\begin{longtable}[]{@{}
  >{\raggedright\arraybackslash}p{(\linewidth - 10\tabcolsep) * \real{0.1667}}
  >{\raggedright\arraybackslash}p{(\linewidth - 10\tabcolsep) * \real{0.1667}}
  >{\raggedright\arraybackslash}p{(\linewidth - 10\tabcolsep) * \real{0.1667}}
  >{\raggedright\arraybackslash}p{(\linewidth - 10\tabcolsep) * \real{0.1667}}
  >{\raggedright\arraybackslash}p{(\linewidth - 10\tabcolsep) * \real{0.1667}}
  >{\raggedright\arraybackslash}p{(\linewidth - 10\tabcolsep) * \real{0.1667}}@{}}
\toprule\noalign{}
\begin{minipage}[b]{\linewidth}\raggedright
Identificador anterior (Entrega 1A)
\end{minipage} & \begin{minipage}[b]{\linewidth}\raggedright
Identificador actual
\end{minipage} & \begin{minipage}[b]{\linewidth}\raggedright
Descripción
\end{minipage} & \begin{minipage}[b]{\linewidth}\raggedright
CU
\end{minipage} & \begin{minipage}[b]{\linewidth}\raggedright
HU
\end{minipage} & \begin{minipage}[b]{\linewidth}\raggedright
Estado
\end{minipage} \\
\midrule\noalign{}
\endhead
\bottomrule\noalign{}
\endlastfoot
RF-01 & REQ-F-001 & Registro de usuario dueño de mascota & CU-01 &
HU-001 & verificado \\
RF-01, RF-02 & REQ-F-008 & Creación de mascota asociada a un dueño
existente & CU-07 & HU-007 & verificado \\
RF-02 & REQ-F-009 & Listado paginado de mascotas activas por
propietario/rol & CU-08 & HU-008 & verificado \\
RF-02 & REQ-F-011 & Actualización de mascota con verificación de
propiedad & CU-10 & HU-010 & verificado \\
RF-03, RF-04 & REQ-F-013 & Registro de atención médica (obligación A) y
consulta del historial clínico consolidado (obligación B) & CU-12 &
HU-012 & implementado \\
RF-05 & REQ-F-014 & Registro de medicamentos prescritos & CU-13 & HU-013
& pendiente \\
RF-06 & REQ-F-015 & Gestión de citas por API (obligación A) y calendario
interactivo (obligación B) & CU-14 & HU-014 & implementado \\
RF-07 & REQ-F-022 & Notificaciones por correo electrónico (cierre de
OBS-02) & CU-21 & HU-021 & pendiente \\
RF-08, RF-09 & REQ-F-016 & API de recepción de telemetría de
dispositivos IoT & CU-15 & HU-015 & pendiente \\
RF-10 & REQ-F-017 & Recomendaciones clínicas informativas (redacción
cerrada en OBS-04) & CU-16 & HU-016 & pendiente \\
RF-11, RF-12 & REQ-F-018 & Comprobantes de pago digitales & CU-17 &
HU-017 & pendiente \\
RF-13, RF-WEB-02 & REQ-F-006 & Control de acceso por rol y por
propietario & CU-05 & HU-005 & implementado \\
RF-14 & REQ-F-019 & Reportes estadísticos exportables & CU-18 & HU-018 &
pendiente \\
RF-15 & REQ-F-020 & Auditoría de operaciones del sistema & CU-19 &
HU-019 & implementado \\
RF-16, RF-WEB-01 & REQ-F-003 & Autenticación mediante usuario y
contraseña & CU-02 & HU-002 & verificado \\
RF-17, RF-WEB-04 & REQ-F-005 & Cierre de sesión con revocación de token
& CU-04 & HU-004 & verificado \\
RNF-03, RNF-WEB-03 & REQ-F-004 & Renovación de sesión (refresh) & CU-03
& HU-003 & verificado \\
\end{longtable}
}

\textbf{Requisitos funcionales sin origen en la Entrega 1A.} Agregados
durante la Tercera Entrega para documentar funcionalidad ya presente en
el código: \texttt{REQ-\allowbreak{}F-\allowbreak{}002} (rechazo de correo duplicado),
\texttt{REQ-\allowbreak{}F-\allowbreak{}007} (consulta del perfil propio), \texttt{REQ-\allowbreak{}F-\allowbreak{}010}
(consulta de mascota por id), \texttt{REQ-\allowbreak{}F-\allowbreak{}012} (baja lógica de
mascota) y \texttt{REQ-\allowbreak{}F-\allowbreak{}021} (resumen de mascotas por especie, exigido
por el bloque A.2.2 de la Guía). Agregados durante la Unidad IV:
\texttt{REQ-\allowbreak{}F-\allowbreak{}023} (gestión administrativa de usuarios),
\texttt{REQ-\allowbreak{}F-\allowbreak{}024} (gestión de vacunas) y \texttt{REQ-\allowbreak{}F-\allowbreak{}025} (consulta
externa de especies con caché). Agregado en la revisión 2026-09-11:
\texttt{REQ-\allowbreak{}F-\allowbreak{}026} (recuperación de contraseña). Ninguno sustituye ni
duplica un identificador \texttt{RF-\allowbreak{}NN} original.

\textbf{Requisitos no funcionales sin origen en la Entrega 1A.}
\texttt{REQ-\allowbreak{}NF-\allowbreak{}004} (claims RFC 7519), \texttt{REQ-\allowbreak{}NF-\allowbreak{}009} (OWASP A09),
\texttt{REQ-\allowbreak{}NF-\allowbreak{}010} (OWASP A07) y \texttt{REQ-\allowbreak{}NF-\allowbreak{}013} (estrategia
híbrida de acceso a datos), exigidos por los bloques A.1, A.2 y C.2 de
la Guía. Agregados en la revisión 2026-09-11: \texttt{REQ-\allowbreak{}NF-\allowbreak{}014}
(indisponibilidad de Redis), \texttt{REQ-\allowbreak{}NF-\allowbreak{}015} (aislamiento de datos
por propietario), \texttt{REQ-\allowbreak{}NF-\allowbreak{}016} (política de contraseñas),
\texttt{REQ-\allowbreak{}NF-\allowbreak{}017} (respaldo y recuperación) y \texttt{REQ-\allowbreak{}NF-\allowbreak{}018}
(accesibilidad).

\subsection{5. Modelo de datos
(referencia)}\label{modelo-de-datos-referencia}

El diccionario de datos completo de las entidades implementadas
(\texttt{usuarios}, \texttt{mascotas}, \texttt{citas},
\texttt{consultas}, \texttt{vacunas}) está en
\texttt{docs/\allowbreak{}diccionario\_\allowbreak{}datos.\allowbreak{}md}, sincronizado con las migraciones
Flyway de \texttt{Backend/\allowbreak{}src/\allowbreak{}main/\allowbreak{}resources/\allowbreak{}db/\allowbreak{}migration/} y con
\texttt{database/\allowbreak{}migrations/\allowbreak{}V1\_\allowbreak{}\_\allowbreak{}schema\_\allowbreak{}inicial.\allowbreak{}sql}. El modelo
entidad-relación conceptual completo (incluyendo las entidades
pendientes: \texttt{Dispositivo\_\allowbreak{}Io\allowbreak{}T}, \texttt{Chat\_\allowbreak{}Triage},
\texttt{Producto\_\allowbreak{}Servicio}, \texttt{Factura},
\texttt{Detalle\_\allowbreak{}Factura}, \texttt{Proveedor}, \texttt{Mercaderia},
\texttt{Detalle\_\allowbreak{}Ingreso}) proviene de la Entrega 1A
(\texttt{PFC\_\allowbreak{}Entrega1\allowbreak{}A\_\allowbreak{}BMT.\allowbreak{}pdf}, sección 6) y se conserva como visión
de producto pendiente de materialización, sin duplicar aquí el DDL
completo. Los estados y transiciones de las entidades implementadas
están en la sección 2.9.

\textbf{Diagrama entidad-relación (DER) --- dos artefactos distintos, no
intercambiables (cierre de OBS-05):}

\begin{itemize}
\tightlist
\item
  \texttt{docs/\allowbreak{}diagrams/\allowbreak{}der-\allowbreak{}biopet/\allowbreak{}der-\allowbreak{}biopet.\allowbreak{}png} --- renderizado
  generado a partir de la fuente Graphviz \texttt{der-\allowbreak{}biopet.\allowbreak{}dot}. Es un
  diagrama \textbf{dibujado}, no una exportación de una herramienta de
  modelado de base de datos.
\item
  \href{../observaciones/evidencias/DER-BIOPET-pgAdmin-ERD-Tool.png}{\texttt{docs/\allowbreak{}observaciones/\allowbreak{}evidencias/\allowbreak{}DER-\allowbreak{}BIOPET-\allowbreak{}pg\allowbreak{}Admin-\allowbreak{}ERD-\allowbreak{}Tool.\allowbreak{}png}}
  --- exportación \textbf{real} generada desde \textbf{pgAdmin 4,
  herramienta ERD Tool}, solicitada explícitamente por la
  retroalimentación oficial del SGA de la Entrega 1B: \emph{``Exportar
  el DER desde pgAdmin 4 (ERD Tool) como PNG de alta resolución para el
  informe final''} (ver \texttt{docs/\allowbreak{}observaciones/\allowbreak{}OBSERVACIONES.\allowbreak{}md},
  OBS-05). PNG válido (firma de archivo verificada), 768×883 px,
  generado directamente sobre el esquema real de PostgreSQL, no sobre la
  fuente \texttt{.\allowbreak{}dot}.
\end{itemize}

\subsection{6. Interfaces de usuario
(referencia)}\label{interfaces-de-usuario-referencia}

Los wireframes conceptuales (login unificado, dashboard por rol,
historial clínico, gestión de citas) están documentados en la Entrega 1A
(\texttt{PFC\_\allowbreak{}Entrega1\allowbreak{}A\_\allowbreak{}BMT.\allowbreak{}pdf}, sección 7). La interfaz realmente
implementada en v1.0.0 está en \texttt{frontend/\allowbreak{}src/\allowbreak{}app/\allowbreak{}features/}:

{\def\LTcaptype{none} % do not increment counter
\begin{longtable}[]{@{}
  >{\raggedright\arraybackslash}p{(\linewidth - 4\tabcolsep) * \real{0.3333}}
  >{\raggedright\arraybackslash}p{(\linewidth - 4\tabcolsep) * \real{0.3333}}
  >{\raggedright\arraybackslash}p{(\linewidth - 4\tabcolsep) * \real{0.3333}}@{}}
\toprule\noalign{}
\begin{minipage}[b]{\linewidth}\raggedright
Componente
\end{minipage} & \begin{minipage}[b]{\linewidth}\raggedright
Cubre
\end{minipage} & \begin{minipage}[b]{\linewidth}\raggedright
Requisitos
\end{minipage} \\
\midrule\noalign{}
\endhead
\bottomrule\noalign{}
\endlastfoot
\texttt{login.\allowbreak{}component.\allowbreak{}ts} & Autenticación & REQ-F-003, REQ-F-004,
REQ-F-005 \\
\texttt{mascotas.\allowbreak{}component.\allowbreak{}ts} & CRUD de mascotas con paginación,
formularios, confirmación de borrado, resumen por especie y atributos de
accesibilidad & REQ-F-008 a REQ-F-012, REQ-F-021 \\
\texttt{vacunas.\allowbreak{}component.\allowbreak{}ts} & CRUD de vacunas & REQ-F-024 \\
\end{longtable}
}

\textbf{Módulos con API real y sin pantalla propia.} Citas (REQ-F-015,
obligación A), consultas médicas (REQ-F-013, obligación A), gestión
administrativa de usuarios (REQ-F-023) y consulta externa de especies
(REQ-F-025) tienen backend real y verificado, pero \textbf{no} tienen
componente en \texttt{frontend/\allowbreak{}src/\allowbreak{}app/\allowbreak{}features/}. El calendario
interactivo de citas es la obligación B de REQ-F-015, no implementada.
Ninguna de estas pantallas cuenta con wireframe formal actualizado; se
deja como observación abierta en la sección 7.

Las interfaces del sistema con su entorno (software, hardware y
comunicaciones) están en la sección 2.7.

\subsection{7. Observaciones e información
pendiente}\label{observaciones-e-informaciuxf3n-pendiente}

Esta sección se revisó contra el estado real del repositorio en el
commit de revisión \texttt{8130ee00b0083808\allowbreak{}fc60567018589e38\allowbreak{}302b375d}. Se
mantiene el mismo principio que en versiones anteriores: no se marca
nada como resuelto sin evidencia verificable, y no se inventa contenido
para cerrar un pendiente.

\subsubsection{7.1. Cerrado en revisiones
anteriores}\label{cerrado-en-revisiones-anteriores}

\begin{itemize}
\tightlist
\item
  \textbf{ADR de cambio de pila tecnológica} --- CERRADO.
  \texttt{docs/\allowbreak{}adr/\allowbreak{}ADR-\allowbreak{}002-\allowbreak{}pila-\allowbreak{}tecnologica.\allowbreak{}md} existe, está en estado
  ``Aceptado'' y documenta explícitamente la migración de ASP.NET Core 8
  (ADR-001, Entrega 1A) a Java 21 / Spring Boot 3.2, con contexto,
  evidencia verificable en \texttt{Backend/\allowbreak{}pom.\allowbreak{}xml} y alternativas
  consideradas.
\item
  \textbf{Evidencia empírica de rendimiento (REQ-NF-001)} --- CERRADO.
  10 corridas k6 reales (5 en caliente, 5 en frío) del 2026-09-03 en
  \texttt{docs/\allowbreak{}mediciones/\allowbreak{}perf/}, con p95 entre 6,04 ms y 18,56 ms, muy
  por debajo del umbral de 200/500 ms.
\item
  \textbf{Evidencia OWASP A07/A09 (REQ-NF-009, REQ-NF-010)} --- CERRADO.
  \texttt{docs/\allowbreak{}mediciones/\allowbreak{}sec/\allowbreak{}A07-\allowbreak{}authentication.\allowbreak{}md} y
  \texttt{A09-\allowbreak{}logging.\allowbreak{}md} contienen capturas HTTP (\texttt{curl}) y de
  log reales, no sólo pruebas JUnit.
\item
  \textbf{Lighthouse --- perfil de escritorio y re-ejecución tras el fix
  de SEO} --- CERRADO. Re-corrida real del 2026-08-18
  (\texttt{docs/\allowbreak{}mediciones/\allowbreak{}lighthouse/\allowbreak{}lhci-\allowbreak{}20260818-\allowbreak{}0538-*.\allowbreak{}json}, 12
  archivos). SEO pasó de 82/100 a 100/100 en las 12 corridas;
  Performance ≥ 90 en todos los casos; Accessibility 91 en todos (ver
  REQ-NF-005 y REQ-NF-018).
\item
  \textbf{PRISMA (trabajos relacionados)} --- CERRADO con 10 estudios y
  39 referencias BibTeX integradas en
  \texttt{docs/\allowbreak{}informe/\allowbreak{}referencias.\allowbreak{}bib}.
\item
  \textbf{INVEST} --- CERRADO. Sección 3.3, con los seis criterios
  aplicados explícitamente a las 24 historias de usuario.
\item
  \textbf{Bitácora de observaciones}
  (\texttt{docs/\allowbreak{}observaciones/\allowbreak{}OBSERVACIONES.\allowbreak{}md}): cubre OBS-01 a OBS-15
  con fuente primaria (capturas SGA) y verificación cruzada contra
  \texttt{git\ log}.
\item
  \textbf{REQ-NF-013 (estrategia híbrida de acceso a datos)} ---
  CERRADO. \texttt{scripts/\allowbreak{}audit-\allowbreak{}sql-\allowbreak{}dynamic.\allowbreak{}sh} se reprodujo sobre el
  commit exacto del tag histórico \texttt{v1.\allowbreak{}0.\allowbreak{}0}
  (\texttt{0d5cd525ce648cca\allowbreak{}7219da204e16fa62\allowbreak{}2e671a87}) con resultado ``0
  hallazgos en 7 archivos. PASA''; evidencia archivada en
  \texttt{docs/\allowbreak{}mediciones/\allowbreak{}sec/\allowbreak{}audit-\allowbreak{}sql-\allowbreak{}dynamic-\allowbreak{}v1.\allowbreak{}0.\allowbreak{}0.\allowbreak{}txt}.
\end{itemize}

\subsubsection{7.2. Cerrado en la revisión 2026-09-11 (esta
revisión)}\label{cerrado-en-la-revisiuxf3n-2026-09-11-esta-revisiuxf3n}

\begin{itemize}
\tightlist
\item
  \textbf{Criterios de aceptación en todos los requisitos} --- CERRADO.
  Los 47 bloques de las secciones 3.1 y 3.2 declaran criterios de
  aceptación observables en forma \emph{Dado/cuando/entonces}, que
  permiten responder CUMPLE / NO CUMPLE. Antes sólo dos requisitos los
  tenían.
\item
  \textbf{Bloques individuales para REQ-F-013 a REQ-F-020} --- CERRADO.
  Los ocho requisitos que aparecían únicamente como filas de una tabla
  resumen tienen ahora bloque completo con la plantilla única.
\item
  \textbf{Taxonomía de estados} --- CERRADO. La sección 1.3.1 define
  exactamente tres estados y todo el documento los usa sin excepción.
  Los valores \texttt{parcial}, \texttt{verificado\ parcia\allowbreak{}lmente},
  \texttt{implementado\ pero\ n\allowbreak{}o\ verificado},
  \texttt{pendiente\ de\ evide\allowbreak{}ncia\ archivada} y
  \texttt{verificado\ en\ conf\allowbreak{}iguración\ TLS\ de\ de\allowbreak{}sarrollo} han
  desaparecido del campo Estado; sus matices se trasladaron al campo
  Observaciones del requisito correspondiente.
\item
  \textbf{Requisitos compuestos} --- CERRADO. \texttt{REQ-\allowbreak{}F-\allowbreak{}013},
  \texttt{REQ-\allowbreak{}F-\allowbreak{}015} y \texttt{REQ-\allowbreak{}NF-\allowbreak{}012} se descomponen en
  subrequisitos con estado propio, conservando el identificador base y
  la trazabilidad histórica.
\item
  \textbf{Sección 2.6 (servicios externos)} --- CERRADO. Se corrige la
  contradicción: la integración con API Ninjas \emph{sí} existe y se
  declara explícitamente, junto con las tres integraciones que no
  existen.
\item
  \textbf{REQ-NF-007 (disponibilidad)} --- CERRADO como redacción. El
  enunciado no verificable se reformula en criterios medibles y el
  compromiso académico se traslada a la sección 2.5 como restricción
  operacional.
\item
  \textbf{Interfaces externas, matriz de permisos y estados del dominio}
  --- CERRADO. Secciones 2.7, 2.8 y 2.9, construidas exclusivamente
  desde el código de este repositorio.
\item
  \textbf{Control documental} --- CERRADO. Fecha de emisión explícita,
  historial de revisiones respaldado por commits reales (sección 1.6) y
  eliminación del índice manual duplicado.
\item
  \textbf{Correspondencia exacta SRS ↔ matriz} --- CERRADO. Una versión
  intermedia de esta revisión había descompuesto \texttt{REQ-\allowbreak{}F-\allowbreak{}013},
  \texttt{REQ-\allowbreak{}F-\allowbreak{}015} y \texttt{REQ-\allowbreak{}NF-\allowbreak{}012} en subrequisitos con sufijo
  de letra, lo que dejaba al SRS con identificadores que la matriz no
  contenía. Se revirtió: los tres conservan un identificador único y
  separan sus obligaciones dentro del bloque, de modo que ambos
  documentos declaran exactamente los mismos 44 identificadores.
\item
  \textbf{Referencia a prueba inexistente en REQ-F-001} --- CERRADO. Se
  retira la cita al test ``\texttt{Auth\allowbreak{}Controller\allowbreak{}Test} (registro
  exitoso)'', que no existe, y se sustituye por la evidencia HTTP real
  que sí existe.
\end{itemize}

\subsubsection{7.3. Limitaciones declaradas que siguen
abiertas}\label{limitaciones-declaradas-que-siguen-abiertas}

Ninguna de estas se cierra con datos inventados.

\begin{enumerate}
\def\labelenumi{\arabic{enumi}.}
\tightlist
\item
  \textbf{Aislamiento de datos incompleto en Consultas (REQ-F-006,
  REQ-F-013, REQ-NF-015).} \texttt{Consulta\allowbreak{}Service.\allowbreak{}listar()} no filtra
  por propietario para \texttt{ROLE\_\allowbreak{}DUENO}: ambas ramas del \texttt{if}
  devuelven \texttt{consulta\allowbreak{}Repository.\allowbreak{}find\allowbreak{}All\allowbreak{}By\allowbreak{}Activo\allowbreak{}True(pageable)}, y
  \texttt{Consulta\allowbreak{}Repository} no declara el método de filtrado que sí
  tienen \texttt{Cita\allowbreak{}Repository} y \texttt{Vacuna\allowbreak{}Repository}. Un dueño
  autenticado puede ver consultas médicas de mascotas ajenas a través de
  \texttt{GET\ /\allowbreak{}api/\allowbreak{}consultas}. \textbf{Requiere cambio de código de
  backend}, fuera del alcance de esta revisión documental.
\item
  \textbf{Sin política definida ante indisponibilidad de Redis
  (REQ-NF-014).} Con Redis caído,
  \texttt{Token\allowbreak{}Blacklist\allowbreak{}Service.\allowbreak{}is\allowbreak{}Revoked} lanza
  \texttt{Redis\allowbreak{}Connection\allowbreak{}Failure\allowbreak{}Exception}, que no captura ni
  \texttt{Jwt\allowbreak{}Authentication\allowbreak{}Filter} ni \texttt{Global\allowbreak{}Exception\allowbreak{}Handler},
  produciendo un error \texttt{500} sin \texttt{Problem\allowbreak{}Detail}. No hay
  \texttt{Cache\allowbreak{}Error\allowbreak{}Handler} registrado. El requisito especifica la
  política que deberá implementarse. \textbf{Requiere cambio de código
  de backend.}
\item
  \textbf{Cobertura de pruebas del módulo de Consultas (REQ-F-013).}
  \texttt{GET\ /\allowbreak{}api/\allowbreak{}consultas} (listado) y
  \texttt{PUT\ /\allowbreak{}api/\allowbreak{}consultas/\allowbreak{}\{id\}} no tienen prueba automatizada
  dedicada: \texttt{Consulta\allowbreak{}Controller\allowbreak{}Test} tiene 6 pruebas frente a las
  25 de \texttt{Cita\allowbreak{}Controller\allowbreak{}Test} para el módulo equivalente.
\item
  \textbf{Hit ratio de caché Redis (REQ-NF-012, obligación B).} Sigue
  sin una medición de \texttt{keyspace\_\allowbreak{}hits}/\texttt{keyspace\_\allowbreak{}misses};
  sólo hay evidencia de TTL y existencia de clave bajo carga.
\item
  \textbf{Compatibilidad de navegadores (REQ-NF-006).} Sin evidencia
  archivada de ejecución de los flujos críticos en Chrome, Firefox y
  Edge, y sin configuración específica de compatibilidad en el
  repositorio (no hay \texttt{browserslist} propio). Ninguno de sus
  criterios tiene artefacto, por lo que el requisito queda
  \texttt{pendiente}.
\item
  \textbf{Disponibilidad sostenida (REQ-NF-007).} No existe monitoreo
  continuo ni registro fechado de sondeos a \texttt{/\allowbreak{}actuator/\allowbreak{}health}
  durante una ventana de evaluación; no se declara ninguna cifra de
  disponibilidad histórica.
\item
  \textbf{Auditoría de operaciones de negocio (REQ-F-020).} Sólo se
  auditan los eventos de autenticación; no hay registro de las
  operaciones CRUD sobre mascotas, citas, consultas, vacunas ni
  usuarios.
\item
  \textbf{Historial clínico consolidado sin endpoint (REQ-F-013,
  obligación B).} La rutina \texttt{fn\_\allowbreak{}historial\_\allowbreak{}clini\allowbreak{}co\_\allowbreak{}mascota} y
  su repositorio existen y están probados, pero ningún controlador los
  expone; además la rutina devuelve un resumen consolidado, no la lista
  cronológica completa de eventos.
\item
  \textbf{Reportes exportables sin endpoint (REQ-F-019).}
  \texttt{fn\_\allowbreak{}reporte\_\allowbreak{}dashboa\allowbreak{}rd} existe y está probada; la exportación
  a PDF y Excel y el endpoint no existen.
\item
  \textbf{\texttt{fn\_\allowbreak{}siguiente\_\allowbreak{}numer\allowbreak{}o\_\allowbreak{}ficha} sin consumidor
  (REQ-NF-013).} El generador de códigos correlativos está implementado
  y probado, pero ningún servicio de la aplicación lo invoca; su
  consumidor natural es REQ-F-018.
\item
  \textbf{Cobertura de \texttt{GET\ /\allowbreak{}api/\allowbreak{}vacunas/\allowbreak{}mascota/\allowbreak{}\{mascota\allowbreak{}Id\}}
  (REQ-F-024).} Sin prueba automatizada dedicada; la corrida Newman
  archivada (\texttt{docs/\allowbreak{}mediciones/\allowbreak{}postman/\allowbreak{}newman-\allowbreak{}report.\allowbreak{}json},
  2026-07-31) es anterior a la colección de vacunas y no puede citarse
  como evidencia de su ejecución.
\item
  \textbf{Recuperación de contraseña ausente (REQ-F-026).} No existe
  implementación alguna; hoy la única vía de restablecimiento es que un
  \texttt{ROLE\_\allowbreak{}ADMIN} sobrescriba la contraseña. Se incorpora como
  requisito \texttt{pendiente}, con su historia de usuario y su caso de
  uso por crear.
\item
  \textbf{Accesibilidad (REQ-NF-018).} Se declara únicamente el umbral
  de la métrica automatizada de Lighthouse (≥ 90, medido 91 en las 12
  corridas). \textbf{No se declara conformidad con WCAG 2.1 en ningún
  nivel}: no hay auditoría manual, pruebas con lector de pantalla ni
  revisión completa de navegación por teclado. La pantalla de vacunas no
  está entre las rutas auditadas, la medición se hizo contra
  \texttt{http:/\allowbreak{}/\allowbreak{}localhost:4200} y las sumas SHA-256 archivadas
  corresponden a la corrida del 2026-08-01, no a las 12 corridas del
  2026-08-18 que se citan como evidencia.
\item
  \textbf{Restauración sin log archivado y respaldo no automatizado
  (REQ-NF-017).} El procedimiento está documentado y el volcado de
  ejemplo \texttt{docs/\allowbreak{}despliegue/\allowbreak{}ejemplo-\allowbreak{}backup-\allowbreak{}20260817.\allowbreak{}sql} es real y
  verificable (6 tablas, 6 secuencias, 4 triggers y 7 rutinas de este
  proyecto), pero la ejecución de la restauración y la invocación de una
  rutina sobre la base restaurada constan sólo como narración en
  \texttt{BACKUP.\allowbreak{}md}, sin log crudo archivado. Además no existe guion ni
  tarea programada en el repositorio que ejecute el respaldo diario: es
  un procedimiento manual del operador.
\item
  \textbf{Límite de intentos de login por instancia (REQ-NF-010).} El
  contador vive en memoria del proceso; en un despliegue multiinstancia
  el límite se aplicaría por instancia. El despliegue actual es de
  instancia única.
\item
  \textbf{Separación de capas sin control automatizado (REQ-NF-011).} No
  existe una regla de arquitectura automatizada (por ejemplo ArchUnit)
  que impida introducir una dependencia de controlador a repositorio.
\item
  \textbf{Wireframes desactualizados.} Los de la pantalla de Mascotas y
  de los módulos pendientes (historial clínico, citas, facturación)
  siguen siendo los de la Entrega 1A. El producto real y su
  documentación de arquitectura (C4, DER, este SRS) sí reflejan el
  sistema implementado.
\item
  \textbf{Pantallas sin interfaz propia.} Citas, historial clínico
  consolidado, gestión administrativa de usuarios y consulta externa de
  especies tienen backend real, pero ninguna tiene componente en
  \texttt{frontend/\allowbreak{}src/\allowbreak{}app/\allowbreak{}features/}.
\item
  \textbf{Sin SIEM centralizado (REQ-NF-009).} Los eventos de auditoría
  de autenticación quedan en el log de la aplicación.
\item
  \textbf{Límite superior de contraseña sin prueba (REQ-NF-016).} El
  máximo de 80 caracteres consta en la anotación
  \texttt{@Size(min\ =\ 8,\allowbreak{}\ max\ =\ 80)} de \texttt{Registro\allowbreak{}Request} y
  \texttt{Usuario\allowbreak{}Request}, pero ninguna prueba envía una contraseña más
  larga para comprobar el \texttt{422}. Es la única regla de la política
  que carece de cobertura automatizada.
\end{enumerate}

\subsubsection{7.4. Acciones que requieren cambio de código (fuera del
alcance de esta
revisión)}\label{acciones-que-requieren-cambio-de-cuxf3digo-fuera-del-alcance-de-esta-revisiuxf3n}

Los puntos 1 y 2 de la lista anterior son los únicos que impiden que un
requisito \texttt{Must} alcance el estado \texttt{verificado} y que
\textbf{no} pueden resolverse documentalmente:

{\def\LTcaptype{none} % do not increment counter
\begin{longtable}[]{@{}
  >{\raggedright\arraybackslash}p{(\linewidth - 4\tabcolsep) * \real{0.3333}}
  >{\raggedright\arraybackslash}p{(\linewidth - 4\tabcolsep) * \real{0.3333}}
  >{\raggedright\arraybackslash}p{(\linewidth - 4\tabcolsep) * \real{0.3333}}@{}}
\toprule\noalign{}
\begin{minipage}[b]{\linewidth}\raggedright
\#
\end{minipage} & \begin{minipage}[b]{\linewidth}\raggedright
Cambio necesario
\end{minipage} & \begin{minipage}[b]{\linewidth}\raggedright
Requisitos que desbloquea
\end{minipage} \\
\midrule\noalign{}
\endhead
\bottomrule\noalign{}
\endlastfoot
1 & Añadir \texttt{find\allowbreak{}All\allowbreak{}By\allowbreak{}Mascota\_\allowbreak{}Duenio\_\allowbreak{}Id\allowbreak{}And\allowbreak{}Activo\allowbreak{}True} a
\texttt{Consulta\allowbreak{}Repository} y usarlo en la rama \texttt{ROLE\_\allowbreak{}DUENO} de
\texttt{Consulta\allowbreak{}Service.\allowbreak{}listar()}, con prueba de acceso cruzado en
\texttt{Consulta\allowbreak{}Controller\allowbreak{}Test}. & REQ-F-006, REQ-F-013, REQ-NF-015 \\
2 & Capturar el fallo de conexión con Redis en
\texttt{Token\allowbreak{}Blacklist\allowbreak{}Service}/\texttt{Jwt\allowbreak{}Authentication\allowbreak{}Filter} y
traducirlo a \texttt{503} con \texttt{Problem\allowbreak{}Detail} (\emph{fail-closed}
para la revocación), y registrar un \texttt{Cache\allowbreak{}Error\allowbreak{}Handler} que
degrade la caché a consulta directa (\emph{fail-open} limitado a la
caché), con la prueba correspondiente. & REQ-NF-014 \\
\end{longtable}
}

Ambas quedan registradas como acciones de seguimiento para el
propietario de los módulos correspondientes. Este documento \textbf{no
modifica código}.

\section{APROBACIÓN DEL DOCENTE-DIRECTOR}\label{aprobacion-del-docente-director}

\begin{itemize}
\tightlist
\item
  \textbf{Nombre:}
\item
  \textbf{Cargo:}
\item
  \textbf{Firma electrónica:}
\item
  \textbf{Fecha:}
\end{itemize}


\end{document}
